% Leçon 30 — Expression écrite : facteurs de la corruption ; caractères complexes — Chinois 1ère A — Cours signé OnBuch+
% Programme officiel MINESEC. Compiler avec : tectonic ZHA30-ecrit-facteurs-de-la-corruption-caracteres-complex.tex
\newcommand{\DOCMATIERE}{Chinois}
\newcommand{\DOCNIVEAU}{1ère A}
\newcommand{\DOCMODULE}{Module 3 : Citoyenneté et ouverture au monde}
\newcommand{\DOCLECON}{ZHA30}
\newcommand{\DOCTITRE}{Leçon 30 — Expression écrite : facteurs de la corruption ; caractères complexes}
% OnBuch+ — préambule « cours signés » ENRICHI (compilé avec tectonic / XeLaTeX)
% Système visuel « Pop » d'OnBuch+ : fond crème (jamais blanc), bordures encre
% épaisses, ombres « dures » décalées, titres Archivo Black en capitales.
% Version MATHÉMATIQUES : graphes (pgfplots/tikz), boîtes pédagogiques
% (définition, propriété, méthode, attention, le savais-tu, exercice résolu,
% exercice, TP, bilan), page de garde et signature OnBuch+ ancrée en bas.
%
% Macros attendues dans main.tex AVANT \input{preamble} :
%   \DOCMATIERE  \DOCNIVEAU  \DOCMODULE  \DOCLECON (numéro)  \DOCTITRE
\documentclass[11pt]{article}

\usepackage{fontspec}
\usepackage[french]{babel} % français = langue principale ; le chinois via \zh{..} / textebox (xeCJK)
\usepackage{xeCJK}
\usepackage{pifont} % \ding{51} \ding{55} utilisés par les modèles
\usepackage[a4paper,margin=2cm,top=3.4cm,bottom=2.8cm,headheight=1.4cm,footskip=1.3cm]{geometry}
\usepackage{amsmath,amssymb,amsfonts}
\usepackage{mathtools} % \xmapsto, \coloneqq... souvent utilisés par les modèles
\usepackage{cancel} % simplifications barrées (bilans de forces)
% \abs{z} (module, valeur absolue) et \norm{u} (norme), utilisés par les modèles
\providecommand{\abs}{}\let\abs\relax\DeclarePairedDelimiter{\abs}{\lvert}{\rvert}
\providecommand{\norm}{}\let\norm\relax\DeclarePairedDelimiter{\norm}{\lVert}{\rVert}
\usepackage[table]{xcolor} % [table] : fournit aussi \rowcolors (alternance de lignes)
\usepackage{pagecolor}
\usepackage{tikz}
\usetikzlibrary{arrows.meta,positioning,calc,shapes.geometric,shapes.misc,shapes.arrows,decorations.pathmorphing,decorations.pathreplacing,decorations.markings,patterns,shadows,mindmap,trees,fit,backgrounds,matrix,intersections,angles,quotes}
\usepackage{pgfplots}
\usepackage[version=4]{mhchem} % équations nucléaires \ce{^{235}_{92}U}
\usepackage[european,siunitx]{circuitikz} % schémas électriques
\pgfplotsset{compat=1.18}
\usepgfplotslibrary{fillbetween,groupplots} % aires entre courbes (calcul intégral)
\usepackage{enumitem}
\usepackage{fancyhdr}
% [most] charge toutes les bibliothèques tcolorbox (listings, minted, raster,
% documentation...) et gonfle inutilement la mémoire TeX sur un document long
% (~300 boîtes) : on ne charge que ce qui est réellement utilisé.
\usepackage[skins,breakable]{tcolorbox}
\usepackage{graphicx}
\usepackage{colortbl}
\usepackage{booktabs}
\usepackage{tabularx}
\usepackage{diagbox} % case d'en-tête diagonale des tableaux à double entrée
% Colonnes extensibles centrée / gauche / droite (C, L, R), souvent utilisées par les modèles
\newcolumntype{C}{>{\centering\arraybackslash}X}
\newcolumntype{L}{>{\raggedright\arraybackslash}X}
\newcolumntype{R}{>{\raggedleft\arraybackslash}X}
\usepackage{array}
\usepackage{multicol}
\usepackage{multirow}
\usepackage{siunitx}
% parse-numbers=false : accepte aussi \SI{2,5 \times 10^{-3}}{...}
\sisetup{locale=FR,output-decimal-marker={,},group-minimum-digits=4,parse-numbers=false}
\DeclareSIUnit\ohm{\ensuremath{\symup{\Omega}}} % U+2126 absent de la police
\DeclareSIUnit\division{div} % réglages d'oscilloscope (ms/div, V/div)
\DeclareSIUnit\tr{tr} % tours (vitesse de rotation, ex. \SI{1500}{\tr\per\minute})
\DeclareSIUnit\molar{M} % concentration molaire (ex. \SI{3}{\molar}, \SI{1}{\micro\molar})
\DeclareSIUnit\an{an} % année (le modèle confond parfois avec \year, un compteur TeX) — ex. \SI{5}{\kilo\meter\per\an}
\DeclareSIUnit\FCFA{FCFA} % franc CFA (ex. \SI{500}{\FCFA\per\kilo\gram})
\DeclareSIUnit\degreeBrix{\textdegree Brix} % teneur en sucre d'un sirop/jus (réfractométrie)
\DeclareSIUnit\month{mois} % le modèle utilise parfois \month, un compteur TeX, comme unité de durée
\usepackage{qrcode}
% \degree : commande fréquemment utilisée par les modèles pour un angle ou une
% température en dehors de \SI{}{} (non fournie par les paquets chargés).
\providecommand{\degree}{^{\circ}}
% \qed (fin de démonstration), utilisé par les modèles sans amsthm
\providecommand{\qed}{\hfill$\square$}
% \barre : double barre des valeurs interdites dans un tableau de variations (tkz-tab)
\providecommand{\barre}{\ensuremath{\|}}
% environnement proof (amsthm non chargé) utilisé par les modèles
\ifdefined\proof\else\newenvironment{proof}[1][Démonstration]{\par\noindent\textit{#1.}\ }{\qed\par}\fi
\usepackage{lastpage}
\usepackage{hyperref}
\hypersetup{hidelinks}

% ============================================================
% Palette « Pop » OnBuch+ — mêmes hex que la classe P (plus_ui.dart)
% ============================================================
\definecolor{poporange}{HTML}{F59321}
\definecolor{popdark}{HTML}{E07A0C}
\definecolor{popcream}{HTML}{FFF6E8}
\definecolor{popink}{HTML}{1C1714}
\definecolor{poppurple}{HTML}{7A5AE0}
\definecolor{popgreen}{HTML}{2E9E6A}
\definecolor{poppink}{HTML}{E0457B}
\definecolor{popblue}{HTML}{2F7DE1}
\definecolor{popgold}{HTML}{C98A12}
\definecolor{popmuted}{HTML}{8A7F76}
\definecolor{popline}{HTML}{EADFD2}
\definecolor{popgray}{HTML}{8A7F76}
\colorlet{popgrey}{popgray}
% Alias tolérants (le modèle nomme parfois les couleurs différemment)
\colorlet{popwhite}{white}
\colorlet{popblack}{popink}
\colorlet{popred}{poppink}
\colorlet{popyellow}{popgold}
% Teintes claires (fonds de tableaux/encadrés) — le modèle nomme parfois
% les couleurs avec un suffixe L (ex. popblueL) sans les avoir définies.
\colorlet{poporangeL}{poporange!20}
\colorlet{popdarkL}{popdark!20}
\colorlet{poppurpleL}{poppurple!20}
\colorlet{popgreenL}{popgreen!20}
\colorlet{poppinkL}{poppink!20}
\colorlet{popblueL}{popblue!20}
\colorlet{popgoldL}{popgold!20}
\colorlet{popredL}{popred!20}
\colorlet{popgrayL}{popgray!20}
% Teintes claires pour fonds de tableaux / zones
\colorlet{popblueL}{popblue!10!white}
\colorlet{poporangeL}{poporange!14!white}
\colorlet{popgreenL}{popgreen!12!white}
\colorlet{poppurpleL}{poppurple!12!white}
\colorlet{poppinkL}{poppink!12!white}
\colorlet{popgoldL}{popgold!14!white}

\pagecolor{popcream}

% ============================================================
% Typographie
% ============================================================
\defaultfontfeatures{Ligatures=TeX}
\setmainfont{PlusJakartaSans-Medium.ttf}[Path=../../../fonts/,BoldFont=PlusJakartaSans-Bold.ttf]
% Symboles, API et emojis absents de la police principale : repli sur DejaVu Sans (caractères actifs)
\newfontface\symfont{DejaVuSans.ttf}[Path=../../../fonts/]
\makeatletter
\newcommand\symactive[1]{\catcode#1=13 \begingroup\lccode`\~=#1 \lowercase{\endgroup\def~}{{\symfont\char#1}}}
\newcommand\symblank[1]{\catcode#1=13 \begingroup\lccode`\~=#1 \lowercase{\endgroup\def~}{}}
\newcommand\symhyph[1]{\catcode#1=13 \begingroup\lccode`\~=#1 \lowercase{\endgroup\def~}{\mbox{-}}}
\newcount\symc
\newcommand\symrange[2]{\symc=#1 \@whilenum\symc<#2 \do{\expandafter\symactive\expandafter{\the\symc}\advance\symc by 1 }}
\newcommand\symblankrange[2]{\symc=#1 \@whilenum\symc<#2 \do{\expandafter\symblank\expandafter{\the\symc}\advance\symc by 1 }}
\makeatother
\symrange{"0250}{"02FF}\symactive{"2713}\symactive{"2714}\symactive{"2717}\symactive{"2718}\symactive{"2610}\symactive{"2611}\symactive{"2612}
\symrange{"2600}{"27BF}
\catcode"2705=13 \begingroup\lccode`\~="2705 \lowercase{\endgroup\def~}{{\symfont\char"2713}}\symblank{"26BD}\symblank{"26A1}
\symrange{"2460}{"257F}\symactive{"223C}\symactive{"2248}\symactive{"2260}
\symactive{"01F8}\symactive{"01F9}\symactive{"1E3E}\symactive{"1E3F}
\symhyph{"2011}\symhyph{"2E1A}\symblankrange{"1F300}{"1FAFF}\symblank{"FE0F}
% Caractères chinois : WenQuanYi Zen Hei (sous-ensemble GB2312 + pinyin), gras/italique simulés
\setCJKmainfont{WenQuanYiZenHei-subset.ttf}[Path=../../../fonts/,AutoFakeBold=3,AutoFakeSlant=0.2]
\xeCJKsetup{CJKspace=false}
% Pinyin : voyelles à caron (ǎ ǐ ǒ ǔ ǖ ǘ ǚ ǜ) absentes de la police principale -> caractères actifs
\newfontface\pyfont{WenQuanYiZenHei-subset.ttf}[Path=../../../fonts/]
\newcommand\pyactive[1]{\catcode#1=13 \begingroup\lccode`\~=#1 \lowercase{\endgroup\def~}{{\pyfont\char#1}}}
\pyactive{"01CD}\pyactive{"01CE}\pyactive{"01CF}\pyactive{"01D0}\pyactive{"01D1}\pyactive{"01D2}\pyactive{"01D3}\pyactive{"01D4}\pyactive{"01D5}\pyactive{"01D6}\pyactive{"01D7}\pyactive{"01D8}\pyactive{"01D9}\pyactive{"01DA}\pyactive{"01DB}\pyactive{"01DC}
% Maths en sans-serif (Fira Math) avec chiffres et lettres droites de Plus
% Jakarta, pour que les formules se fondent dans le texte.
\usepackage[math-style=ISO,bold-style=ISO]{unicode-math}
\setmathfont{FiraMath-Regular.otf}[Path=../../../fonts/]
\setmathfont{PlusJakartaSans-Medium.ttf}[Path=../../../fonts/,range={up/num,up/{Latin,latin},\mathup}]
\usepackage{icomma} % virgule décimale sans espace en mode math
% Caractères mathématiques Unicode absents des polices de texte (Plus Jakarta,
% Archivo Black) : rendus par leur équivalent mathématique, en texte comme en
% formule — sinon ils s'affichent en carré vide (ex. « Arithmétique dans ℤ »).
\usepackage{newunicodechar}
\newunicodechar{ℝ}{\ensuremath{\mathbb{R}}}
\newunicodechar{ℤ}{\ensuremath{\mathbb{Z}}}
\newunicodechar{ℂ}{\ensuremath{\mathbb{C}}}
\newunicodechar{ℕ}{\ensuremath{\mathbb{N}}}
\newunicodechar{ℚ}{\ensuremath{\mathbb{Q}}}
\newunicodechar{𝔻}{\ensuremath{\mathbb{D}}}
\newunicodechar{α}{\ensuremath{\alpha}}
\newunicodechar{β}{\ensuremath{\beta}}
\newunicodechar{γ}{\ensuremath{\gamma}}
\newunicodechar{δ}{\ensuremath{\delta}}
\newunicodechar{ε}{\ensuremath{\varepsilon}}
\newunicodechar{θ}{\ensuremath{\theta}}
\newunicodechar{λ}{\ensuremath{\lambda}}
\newunicodechar{μ}{\ensuremath{\mu}}
\newunicodechar{σ}{\ensuremath{\sigma}}
\newunicodechar{τ}{\ensuremath{\tau}}
\newunicodechar{φ}{\ensuremath{\varphi}}
\newunicodechar{ψ}{\ensuremath{\psi}}
\newunicodechar{ω}{\ensuremath{\omega}}
\newunicodechar{Σ}{\ensuremath{\Sigma}}
% majuscules produites par \MakeUppercase dans les titres (α-aminés -> Α)
\newunicodechar{Α}{\ensuremath{\alpha}}
\newunicodechar{Β}{\ensuremath{\beta}}
\newunicodechar{Γ}{\ensuremath{\Gamma}}
\newunicodechar{Δ}{\ensuremath{\Delta}}
\newunicodechar{Ε}{\ensuremath{\varepsilon}}
\newunicodechar{Θ}{\ensuremath{\Theta}}
\newunicodechar{Λ}{\ensuremath{\Lambda}}
\newunicodechar{Μ}{\ensuremath{\mu}}
\newunicodechar{Τ}{\ensuremath{\tau}}
\newunicodechar{Φ}{\ensuremath{\Phi}}
\newunicodechar{Ψ}{\ensuremath{\Psi}}
\newunicodechar{Ω}{\ensuremath{\Omega}}
\newunicodechar{↦}{\ensuremath{\mapsto}}
\newunicodechar{∈}{\ensuremath{\in}}
\newunicodechar{∉}{\ensuremath{\notin}}
\newunicodechar{∧}{\ensuremath{\wedge}}
\newunicodechar{⇔}{\ensuremath{\Leftrightarrow}}
\newunicodechar{⟺}{\ensuremath{\Longleftrightarrow}}
\newunicodechar{⇒}{\ensuremath{\Rightarrow}}
\newunicodechar{⟹}{\ensuremath{\Longrightarrow}}
\newunicodechar{∘}{\ensuremath{\circ}}
\newunicodechar{∪}{\ensuremath{\cup}}
\newunicodechar{∩}{\ensuremath{\cap}}
\newunicodechar{≡}{\ensuremath{\equiv}}
\newunicodechar{⊂}{\ensuremath{\subset}}
\newunicodechar{⊥}{\ensuremath{\perp}}
\newunicodechar{∃}{\ensuremath{\exists}}
\newunicodechar{∀}{\ensuremath{\forall}}
\newunicodechar{∛}{\ensuremath{\sqrt[3]{\,}}}
\newunicodechar{∜}{\ensuremath{\sqrt[4]{\,}}}
\newunicodechar{⃗}{\ensuremath{{}^{\to}}}
\newunicodechar{ⁿ}{\ifmmode^{n}\else\textsuperscript{\ensuremath{n}}\fi}
\newunicodechar{ˣ}{\ifmmode^{x}\else\textsuperscript{\ensuremath{x}}\fi}
\newunicodechar{ᵇ}{\ifmmode^{b}\else\textsuperscript{\ensuremath{b}}\fi}
\newunicodechar{ʳ}{\ifmmode^{r}\else\textsuperscript{\ensuremath{r}}\fi}
\newunicodechar{ᵘ}{\ifmmode^{u}\else\textsuperscript{\ensuremath{u}}\fi}
\newunicodechar{ʸ}{\ifmmode^{y}\else\textsuperscript{\ensuremath{y}}\fi}
\newunicodechar{ᵃ}{\ifmmode^{a}\else\textsuperscript{\ensuremath{a}}\fi}
\newunicodechar{ᵉ}{\ifmmode^{e}\else\textsuperscript{\ensuremath{e}}\fi}
\newunicodechar{ᵏ}{\ifmmode^{k}\else\textsuperscript{\ensuremath{k}}\fi}
\newunicodechar{ᵖ}{\ifmmode^{p}\else\textsuperscript{\ensuremath{p}}\fi}
\newunicodechar{ᴺ}{\ifmmode^{N}\else\textsuperscript{\ensuremath{N}}\fi}
\newunicodechar{ᶜ}{\ifmmode^{c}\else\textsuperscript{\ensuremath{c}}\fi}
\newunicodechar{ʰ}{\ifmmode^{h}\else\textsuperscript{\ensuremath{h}}\fi}
\newunicodechar{ᵠ}{\ifmmode^{q}\else\textsuperscript{\ensuremath{q}}\fi}
\newunicodechar{ₙ}{\ifmmode_{n}\else\textsubscript{\ensuremath{n}}\fi}
\newunicodechar{ₐ}{\ifmmode_{a}\else\textsubscript{\ensuremath{a}}\fi}
\newunicodechar{ᵢ}{\ifmmode_{i}\else\textsubscript{\ensuremath{i}}\fi}
\newunicodechar{ᵣ}{\ifmmode_{r}\else\textsubscript{\ensuremath{r}}\fi}
\newunicodechar{ₖ}{\ifmmode_{k}\else\textsubscript{\ensuremath{k}}\fi}
\newunicodechar{ᵦ}{\ifmmode_{\beta}\else\textsubscript{\ensuremath{\beta}}\fi}
\newunicodechar{ₓ}{\ifmmode_{x}\else\textsubscript{\ensuremath{x}}\fi}
\newfontfamily\popdisplay{ArchivoBlack-Regular.ttf}[Path=../../../fonts/]
\newfontfamily\poplogo{SpaceGrotesk-Bold.ttf}[Path=../../../fonts/]
\newfontfamily\popsemi{PlusJakartaSans-SemiBold.ttf}[Path=../../../fonts/]
\newfontfamily\popxbold{PlusJakartaSans-ExtraBold.ttf}[Path=../../../fonts/]
\newfontfamily\popxboldls{PlusJakartaSans-ExtraBold.ttf}[Path=../../../fonts/,LetterSpace=3.0]

\setlist[enumerate]{leftmargin=1.7em,itemsep=5pt,topsep=4pt}
\setlist[itemize]{leftmargin=1.7em,itemsep=4pt,topsep=4pt,label={\color{poporange}\textbullet}}
\setlength{\parindent}{0pt}
\setlength{\parskip}{5pt}
\renewcommand{\arraystretch}{1.25}

% pgfplots : style Pop par défaut
\pgfplotsset{
  every axis/.append style={axis line style={popink,line width=1pt},tick style={popink},
    grid style={popline,line width=0.5pt},label style={font=\small},tick label style={font=\footnotesize}},
  popcurve/.style={poporange,line width=1.8pt,no markers,smooth},
}
% Même style disponible dans un \draw TikZ simple (hors pgfplots)
\tikzset{popcurve/.style={poporange,line width=1.8pt}}
\tikzset{popgrid/.style={popline,very thin}}
\colorlet{popcurve}{poporange} % le nom du style est parfois utilisé comme couleur

% ============================================================
% Pill / squiggle
% ============================================================
\newcommand{\poppill}[3]{%
  \tikz[baseline=-0.55ex]\node[draw=popink,line width=1.2pt,rounded corners=2.6pt,%
    fill=#2,inner xsep=6.5pt,inner ysep=3pt]{%
    \popxboldls\fontsize{7}{7}\selectfont\color{#3}\MakeUppercase{#1}};%
}
\newcommand{\popsquiggle}[1]{%
  \tikz[baseline=-0.4ex]{
    \draw[#1,line width=2.6pt,line cap=round]
      plot [smooth] coordinates {(0,0.09) (0.34,0.01) (0.68,0.21) (1.02,0.01) (1.36,0.10)};
  }%
}
% Mot-clé surligné (marqueur orange sous le texte)
\newcommand{\cle}[1]{{\popxbold\color{popdark}#1}}

% ============================================================
% En-tête / pied
% ============================================================
\pagestyle{fancy}
\fancyhf{}
\renewcommand{\headrulewidth}{0pt}
\renewcommand{\footrulewidth}{0pt}
\lhead{\raisebox{-0.15ex}{\poplogo\fontsize{13}{13}\selectfont\color{popink}OnBuch\color{poporange}+}}
\rhead{\poppill{\DOCMATIERE\ \textbullet\ \DOCNIVEAU\ \textbullet\ Leçon \DOCLECON}{white}{popink}}
\lfoot{\popxbold\fontsize{7.6}{7.6}\selectfont\color{popmuted}\MakeUppercase{Cours signé OnBuch+}\ \textbullet\ {\popsemi\DOCTITRE}}
\rfoot{\tikz[baseline=-0.6ex]\node[rounded corners=8.5pt,fill=popink,minimum height=17pt,minimum width=17pt,inner xsep=5pt,inner sep=0]{\popxbold\fontsize{8}{8}\selectfont\color{popcream}\thepage\,/\,\pageref*{LastPage}};}

% ============================================================
% Titres
% ============================================================
\newcommand{\chaptitre}[2]{%
  \begin{tcolorbox}[enhanced,colback=poporange,colframe=popink,boxrule=2.6pt,arc=11pt,
      drop shadow=popink,left=15pt,right=15pt,top=12pt,bottom=11pt,before skip=3pt,after skip=18pt]
    \poppill{#2}{popink}{popcream}\par\vspace{9pt}
    {\raggedright\hyphenpenalty=10000\exhyphenpenalty=10000\popdisplay\fontsize{22}{24}\selectfont\color{popcream}\MakeUppercase{#1}\par}
  \end{tcolorbox}
}
% Section numérotée : pastille orange avec le numéro + titre Archivo
\newcounter{popsec}
\newcommand{\coursec}[1]{%
  \stepcounter{popsec}\setcounter{popsub}{0}%
  \par\vspace{16pt}\noindent
  \begin{minipage}{\linewidth}
  \tikz[baseline=-0.7ex]\node[draw=popink,line width=1.6pt,fill=poporange,rounded corners=4pt,minimum size=22pt,inner sep=0]{\popdisplay\fontsize{12}{12}\selectfont\color{popcream}\thepopsec};%
  \hspace{8pt}{\popdisplay\fontsize{15}{17}\selectfont\color{popink}\MakeUppercase{#1}}\par
  \vspace{4pt}\noindent{\color{poporange}\rule{36pt}{3.6pt}}
  \end{minipage}\par\nopagebreak\vspace{8pt}
}
\newcounter{popsub}
\newcommand{\courssub}[1]{%
  \stepcounter{popsub}%
  \par\vspace{9pt}\noindent{\popxbold\fontsize{11.5}{13}\selectfont\color{popink}{\color{poporange}\thepopsec.\thepopsub}\hspace{6pt}#1}\par\nopagebreak\vspace{4pt}
}

% ============================================================
% Boîtes pédagogiques — même recette : fond blanc, bordure épaisse colorée,
% ombre dure encre, badge plein. Titre optionnel [#1].
% ============================================================
\tcbset{popbase/.style={breakable,enhanced,colback=white,boxrule=2.3pt,arc=9pt,
  shadow={3pt}{-3pt}{0pt}{popink},left=14pt,right=14pt,top=11pt,bottom=11pt,before skip=11pt,after skip=15pt,
  fontupper=\popsemi\fontsize{10.6}{14.5}\selectfont\color{popink},
  fontlower=\popsemi\fontsize{10.6}{14.5}\selectfont\color{popink}}}
% Coche dessinée (le glyphe ✓ n'existe pas dans les polices du document)
\newcommand{\popcheck}[1]{\tikz[baseline=-0.5ex]\draw[#1,line width=1.6pt,line cap=round,line join=round](-0.13,0.0)--(-0.04,-0.09)--(0.13,0.1);}
\providecommand{\coche}{\popcheck{popgreen}}
\providecommand{\resultat}[1]{\par\smallskip\noindent\fcolorbox{popgreen}{popgreenL}{\parbox{\dimexpr\linewidth-2\fboxsep-2\fboxrule\relax}{\textbf{Résultat.} #1}}\par\smallskip}
\newcommand{\popboxhead}[3]{\poppill{#1}{#2}{white}\ifx\relax#3\relax\else\hspace{6pt}{\popxbold\fontsize{10.6}{12}\selectfont\color{popink}#3}\fi\par\vspace{7pt}}

\newtcolorbox{aretenir}[1][]{popbase,colframe=popgreen,before upper={\popboxhead{À retenir}{popgreen}{#1}}}
% Texte en chinois (document de lecture, dialogue, extrait) : cadre
% d'étude. Un titre optionnel (source, auteur) s'affiche dans l'en-tête.
\newtcolorbox{textebox}[1][]{popbase,colframe=popblue,before upper={\popboxhead{文本 · Texte}{popblue}{#1}},after upper={}}
% Marques ✓ / ✗ (forme correcte / fautive) souvent utilisées par les modèles
\providecommand{\cross}{\textcolor{poppink}{\ensuremath{\times}}}
\providecommand{\xmark}{\cross}\providecommand{\crossmark}{\cross}\providecommand{\cmark}{\ensuremath{\checkmark}}
% Ligne de réponse / justification à compléter (inventée par les modèles)
\providecommand{\justification}[1]{\par\noindent\textit{Justification~:} \dotfill\par}
% \textipa (package tipa non chargé) : la police ne couvre pas l'API -> simple passage
\providecommand{\textipa}[1]{#1}
% Soulignement (ulem non chargé) : \uline, \ul
\providecommand{\uline}[1]{\underline{#1}}\providecommand{\ul}[1]{\underline{#1}}
% \glossaire{\item ...} inventée par les modèles : liste de glossaire
\providecommand{\glossaire}[1]{\begin{itemize}#1\end{itemize}}
% \alert (beamer) inventée par les modèles : texte mis en évidence
\providecommand{\alert}[1]{\textbf{\textcolor{poppink}{#1}}}
% variantes ASCII d'ordinaux écrites en macro
\providecommand{\iere}{\textsuperscript{re}}\providecommand{\ieme}{\textsuperscript{e}}\providecommand{\ere}{\textsuperscript{re}}\providecommand{\eme}{\textsuperscript{e}}\providecommand{\er}{\textsuperscript{er}}\providecommand{\ie}{\textsuperscript{e}}
% Ordinaux français écrits en macro par les modèles : 1\ière, 3\ième
\newcommand{\ière}{\textsuperscript{re}}\newcommand{\ième}{\textsuperscript{e}}
% \exemple (inventée par les modèles) : étiquette « Exemple : »
\providecommand{\exemple}{\textbf{Exemple~:}\ }
% \square utilisable aussi hors mode math (cases à cocher)
\AtBeginDocument{\let\squareMath\square\renewcommand{\square}{\ensuremath{\squareMath}}}
% Mot ou phrase en chinois à l'intérieur d'un paragraphe français
\newcommand{\zh}[1]{\ifmmode\text{#1}\else{#1}\fi}
\let\eng\zh \let\esp\zh \let\all\zh % alias tolérés
% flèches d'intonation inventées par les modèles
\providecommand{\textsearrow}{}\renewcommand{\textsearrow}{\ensuremath{\searrow}}\providecommand{\textnearrow}{}\renewcommand{\textnearrow}{\ensuremath{\nearrow}}
% \fr{..} : mot français (inventé par les modèles)
\providecommand{\fr}[1]{\textit{#1}}
% zhbox : encadré de texte chinois inventé par les modèles
\newenvironment{zhbox}{\begin{quote}}{\end{quote}}
% \courssubsub : titre de niveau 3 inventé par les modèles
\providecommand{\courssubsub}[1]{\par\medskip\noindent\textbf{#1}\par\smallskip}
% \zhbf : texte chinois en gras inventé par les modèles
\providecommand{\zhbf}[1]{\textbf{#1}}
% \vs : « versus » inventé par les modèles
\providecommand{\vs}{\textit{vs}\ }
% \blank : case à compléter inventée par les modèles
\providecommand{\blank}[1][]{\underline{\hspace{1.6em}}}
\newtcolorbox{exemplebox}[1][]{popbase,colframe=poporange,before upper={\popboxhead{Exemple}{poporange}{#1}}}
\newtcolorbox{definition}[1][]{popbase,colframe=popblue,before upper={\popboxhead{Définition}{popblue}{#1}}}
\newtcolorbox{propriete}[1][]{popbase,colframe=poppurple,before upper={\popboxhead{Propriété}{poppurple}{#1}}}
\newtcolorbox{methode}[1][]{popbase,colframe=popgold,before upper={\popboxhead{Méthode}{popgold}{#1}}}
\newtcolorbox{attention}[1][]{popbase,colframe=poppink,before upper={\popboxhead{Attention}{poppink}{#1}}}
\newtcolorbox{experience}[1][]{popbase,colframe=popblue,colback=popblueL,before upper={\popboxhead{Expérience}{popblue}{#1}}}
% « Le savais-tu ? » : carte violette pleine (contexte, culture, Cameroun)
\newtcolorbox{savaistu}[1][]{popbase,colback=poppurple,colframe=popink,
  fontupper=\popsemi\fontsize{10.4}{14.2}\selectfont\color{white},
  before upper={\poppill{Le savais-tu\,?}{white}{poppurple}\ifx\relax#1\relax\else\hspace{6pt}{\popxbold\color{white}#1}\fi\par\vspace{7pt}}}
% Exercice résolu : énoncé en haut, \tcblower puis la solution sur fond crème
\newcounter{popexo}\newcounter{popexr}
\newtcolorbox{exoresolu}[1][]{popbase,colframe=poporange,colbacklower=popcream,
  segmentation style={popink,line width=1.2pt,dashed},
  before upper={\stepcounter{popexr}\popboxhead{Exercice résolu \thepopexr}{poporange}{#1}},
  before lower={\poppill{Solution}{popink}{popcream}\par\vspace{6pt}}}
% Exercice (non résolu en place) — la correction est dans la partie Corrigés
\newtcolorbox{exercice}[1][]{popbase,colframe=popink,
  before upper={\stepcounter{popexo}\popboxhead{Exercice \thepopexo}{popink}{#1}}}
% Corrigé
\newtcolorbox{corrige}[1][]{popbase,colframe=popgreen,colback=popgreenL,
  before upper={\popboxhead{Corrigé}{popgreen}{#1}}}
% Objectifs / prérequis / bilan
\newtcolorbox{objectifs}{popbase,colframe=popink,colback=poporangeL,
  before upper={\popboxhead{Objectifs de la leçon}{popink}{}}}
\newtcolorbox{prerequis}{popbase,colframe=popink,colback=popblueL,
  before upper={\popboxhead{Prérequis}{popblue}{}}}
\newtcolorbox{bilan}{popbase,colframe=popink,colback=white,boxrule=2.8pt,
  before upper={\popboxhead{Fiche bilan}{poporange}{L'essentiel à connaître pour l'examen}}}
% Figure encadrée avec légende
\newtcolorbox{popfigure}[1][]{enhanced,breakable=false,colback=white,colframe=popline,boxrule=1.4pt,arc=8pt,
  left=10pt,right=10pt,top=10pt,bottom=8pt,before skip=10pt,after skip=12pt,halign=center,
  lower separated=false,halign lower=center,
  fontlower=\popsemi\fontsize{9}{11.5}\selectfont\color{popmuted},
  #1}
\newcommand{\legende}[1]{\par\vspace{6pt}{\popsemi\fontsize{9}{11.5}\selectfont\color{popmuted}#1}}

% ============================================================
% Signature OnBuch+
% ============================================================
\newcommand{\poplogoinline}[1]{{\poplogo\fontsize{#1}{#1}\selectfont\color{popink}OnBuch\color{poporange}+}}
\newcommand{\signatureonbuch}{%
  \begin{tcolorbox}[enhanced,colback=popink,colframe=popink,boxrule=2.2pt,arc=10pt,
      drop shadow=poporange,left=14pt,right=14pt,top=10pt,bottom=10pt,before skip=0pt,after skip=0pt]
    \tikz[baseline=-0.7ex]\node[circle,fill=popgreen,draw=popcream,line width=1.3pt,minimum size=22pt,inner sep=0]{\popcheck{white}};%
    \hspace{9pt}\begin{minipage}[c]{0.62\linewidth}
      {\popxbold\fontsize{12}{13}\selectfont\color{popcream}Signé\ \textbullet\ {\poplogo OnBuch\color{poporange}+}}\par\vspace{2pt}
      {\raggedright\popsemi\fontsize{8}{10.5}\selectfont\color{popline}\MakeUppercase{Équipe pédagogique OnBuch+}\\\MakeUppercase{Conforme au programme officiel MINESEC}\par}
    \end{minipage}\hfill
    {\poplogo\fontsize{22}{22}\selectfont\color{popcream}OnBuch\color{poporange}+}
  \end{tcolorbox}%
}

% ============================================================
% Page de garde
% #1 titre · #2 matière · #3 niveau · #4 module · #5 n° leçon · #6 durée ·
% #7 description / objectifs (texte ou itemize)
% ============================================================
\newcommand{\pagegarde}[7]{%
  \thispagestyle{empty}
  \newgeometry{a4paper,margin=1.7cm,top=1.6cm,bottom=1.4cm}%
  \begin{tikzpicture}[remember picture,overlay]
    \fill[poporange!18] ([xshift=-3.2cm,yshift=-3.6cm]current page.north east) circle (4.6cm);
    \fill[poppurple!14] ([xshift=2.4cm,yshift=7.6cm]current page.south west) circle (3.8cm);
  \end{tikzpicture}%
  {\poplogo\fontsize{30}{30}\selectfont\color{popink}OnBuch\color{poporange}+}\hfill
  \raisebox{0.6ex}{\poppill{Cours signé \textbullet\ Premium}{popink}{popcream}}\par
  \vspace{1pt}\popsquiggle{poppurple}\par
  \vspace{8pt}
  \begin{tcolorbox}[enhanced,colback=poporange,colframe=popink,boxrule=2.8pt,arc=13pt,
      drop shadow=popink,left=18pt,right=18pt,top=12pt,bottom=13pt,after skip=10pt]
    \poppill{#2 \textbullet\ #3}{popink}{popcream}\hspace{5pt}\poppill{#4}{white}{popink}\par\vspace{8pt}
    {\popdisplay\fontsize{14}{14}\selectfont\color{popink}LEÇON #5}\par\vspace{4pt}
    {\raggedright\hyphenpenalty=10000\exhyphenpenalty=10000\popdisplay\fontsize{25}{27}\selectfont\color{popcream}\MakeUppercase{#1}\par}
    \vspace{8pt}
    {\popxbold\fontsize{9.5}{9.5}\selectfont\color{popink}\MakeUppercase{Durée conseillée : #6}}
  \end{tcolorbox}
  \begin{tcolorbox}[enhanced,colback=white,colframe=popink,boxrule=2.2pt,arc=10pt,
      drop shadow=popink,left=16pt,right=16pt,top=10pt,bottom=10pt,after skip=8pt]
    \poppill{Dans cette leçon}{poppurple}{white}\par\vspace{5pt}
    \popsemi\fontsize{9.3}{12.6}\selectfont\color{popink}#7
  \end{tcolorbox}
  \noindent\begin{minipage}[t]{0.487\linewidth}
    \begin{tcolorbox}[enhanced,colback=popgreen,colframe=popink,boxrule=2.2pt,arc=10pt,
        drop shadow=popink,left=11pt,right=11pt,top=10pt,bottom=10pt,height=3.1cm,valign=center]
      \tcbox[colback=white,colframe=popink,boxrule=1.3pt,arc=5pt,boxsep=0pt,nobeforeafter,
        left=3pt,right=3pt,top=3pt,bottom=3pt,valign=center]{\qrcode[height=1.9cm]{https://play.google.com/store/apps/details?id=com.luvvix.onbuch&hl=fr}}%
      \hspace{8pt}\begin{minipage}[c]{0.5\linewidth}
        {\popdisplay\fontsize{11}{12.5}\selectfont\color{white}\MakeUppercase{Télécharge OnBuch}}\par\vspace{4pt}
        {\raggedright\popsemi\fontsize{8.4}{11}\selectfont\color{white}Gratuit \textbullet\ Google Play\\Tuteur IA, annales, concours, résultats\par}
      \end{minipage}
    \end{tcolorbox}
  \end{minipage}\hfill
  \begin{minipage}[t]{0.487\linewidth}
    \begin{tcolorbox}[enhanced,colback=poppurple,colframe=popink,boxrule=2.2pt,arc=10pt,
        drop shadow=popink,left=11pt,right=11pt,top=10pt,bottom=10pt,height=3.1cm,valign=center]
      \tikz[baseline=-0.7ex]\node[circle,fill=white,draw=popink,line width=1.3pt,minimum size=1.6cm,inner sep=0]{\popdisplay\fontsize{24}{24}\selectfont\color{poppurple}+};%
      \hspace{8pt}\begin{minipage}[c]{0.6\linewidth}
        {\popdisplay\fontsize{11}{12.5}\selectfont\color{white}\MakeUppercase{Passe à OnBuch+}}\par\vspace{4pt}
        {\raggedright\popsemi\fontsize{8.4}{11}\selectfont\color{white}Tous les cours signés, Tuteur IA illimité, fiches PDF — onglet OnBuch+ de l'app\par}
      \end{minipage}
    \end{tcolorbox}
  \end{minipage}
  \vspace{8pt}
  \popcontenu
  \vfill
  \signatureonbuch
  \restoregeometry
  \newpage
}

% Rangée « Ce cours contient » (page de garde)
\newcommand{\poptile}[3]{% #1 couleur #2 gros texte #3 légende
  \begin{tcolorbox}[enhanced,colback=#1,colframe=popink,boxrule=2pt,arc=9pt,shadow={2.5pt}{-2.5pt}{0pt}{popink},
      left=6pt,right=6pt,top=9pt,bottom=9pt,halign=center,valign=center,height=1.75cm,width=\linewidth,nobeforeafter]
    {\popdisplay\fontsize{12.5}{13}\selectfont\color{white}\MakeUppercase{#2}}\par\vspace{3pt}
    {\centering\popsemi\fontsize{7.8}{9.5}\selectfont\color{white}#3\par}
  \end{tcolorbox}}
\newcommand{\popcontenu}{%
  \par\noindent{\popxboldls\fontsize{7.5}{7.5}\selectfont\color{popmuted}CE COURS SIGNÉ CONTIENT}\par\vspace{7pt}
  \noindent\begin{minipage}[t]{0.235\linewidth}\poptile{popblue}{Cours}{complet, illustré, pas à pas}\end{minipage}\hfill
  \begin{minipage}[t]{0.235\linewidth}\poptile{poporange}{Méthodes}{et exercices résolus}\end{minipage}\hfill
  \begin{minipage}[t]{0.235\linewidth}\poptile{poppink}{Exercices}{type examen + corrigés}\end{minipage}\hfill
  \begin{minipage}[t]{0.235\linewidth}\poptile{popgreen}{Fiche bilan}{l'essentiel à retenir}\end{minipage}\par}

% Sommaire
\newtcolorbox{sommaire}{enhanced,colback=white,colframe=popink,boxrule=2.2pt,arc=10pt,
  drop shadow=popink,left=18pt,right=18pt,top=13pt,bottom=13pt,before skip=6pt,after skip=20pt,
  before upper={\poppill{Sommaire}{poppurple}{white}\par\vspace{9pt}\popsemi\fontsize{10.6}{14.5}\selectfont\color{popink}}}

% Signature de fin de cours — poussée tout en bas de la dernière page
\newcommand{\findecours}{%
  \par\vfill\nopagebreak
  \begin{center}\popsquiggle{poporange}\end{center}\vspace{4pt}
  \signatureonbuch
}
\AtBeginDocument{\def\llbracket{\mathopen{[\mkern-3.5mu[}}\def\rrbracket{\mathclose{]\mkern-3.5mu]}}} % intervalles d'entiers (glyphe ⟦ absent de la police)
% Glyphes absents de Fira Math : redéfinis à partir de glyphes disponibles
\AtBeginDocument{%
  \def\setminus{\mathbin{\backslash}}\let\smallsetminus\setminus
  \def\vdots{\mathord{\vbox{\baselineskip3.2pt\lineskiplimit0pt\kern4pt\hbox{.}\hbox{.}\hbox{.}}}}%
  \def\ddots{\mathinner{\mkern1mu\raise6.5pt\hbox{.}\mkern2mu\raise3.6pt\hbox{.}\mkern2mu\raise0.7pt\hbox{.}\mkern1mu}}%
  \def\triangle{\mathord{\scalebox{0.9}{$\symup{\Delta}$}}}%
  \def\nabla{\mathord{\raisebox{\depth}{\scalebox{1}[-1]{$\symup{\Delta}$}}}}%
  \def\checkmark{\popcheck{popdark}}%
  \def\star{\mathbin{\ast}}\def\bigstar{\mathord{\ast}}%
  \def\diamond{\mathbin{\scalebox{0.6}{$\Diamond$}}}%
  \def\bigcup{\mathop{\mathchoice{\raisebox{-0.3em}{\scalebox{1.7}{$\cup$}}}{\scalebox{1.25}{$\cup$}}{\cup}{\cup}}\displaylimits}%
  \def\bigcap{\mathop{\mathchoice{\raisebox{-0.3em}{\scalebox{1.7}{$\cap$}}}{\scalebox{1.25}{$\cap$}}{\cap}{\cap}}\displaylimits}%
  \def\lesseqqgtr{\mathrel{\lessgtr}}\def\bigoplus{\mathop{\oplus}\displaylimits}%
}
\providecommand{\ssi}{si et seulement si} % abréviation fréquente
\AtBeginDocument{\providecommand{\wideparen}{\overparen}} % arc de cercle

%% FIGFIX-BEGIN v3 — correctifs globaux des figures (docs/onbuch-generation/tools/figfix)
\usepackage{adjustbox}\usepackage{etoolbox}
% toute figure TikZ/pgfplots plus large que la ligne est réduite à la largeur disponible
\BeforeBeginEnvironment{tikzpicture}{\begin{adjustbox}{max width=\linewidth}}
\AfterEndEnvironment{tikzpicture}{\end{adjustbox}}
% légendes pgfplots lisibles par-dessus les courbes
\pgfplotsset{every axis/.append style={legend style={fill=white,fill opacity=0.92,text opacity=1,draw=black!15,inner sep=3pt}}}
% étiquettes d'arêtes (syntaxe edge["..."]) avec fond blanc
\tikzset{every edge quotes/.append style={fill=white,inner sep=1.2pt}}
% bulles de cartes mentales : texte à la ligne dans la bulle, bulle un peu plus grande
\tikzset{every concept/.append style={text width=2.0cm,align=center,minimum size=2.3cm,inner sep=2pt}}
%% FIGFIX-END
\begin{document}

\pagegarde{Leçon 30 — Expression écrite : facteurs de la corruption ; caractères complexes}{Chinois}{1ère A}{Module 3 : Citoyenneté et ouverture au monde}{ZHA30}{2 h}{Cette leçon d'expression écrite apprend à rédiger un texte simple en chinois sur les facteurs de la corruption, à partir d'un modèle commenté et de connecteurs utiles. Elle entraîne aussi l'écriture correcte des caractères complexes du thème (traits, radicaux, caractères proches) et la traduction dans les deux sens (thème et version).
\vspace{4pt}
\begin{multicols}{2}\small
\begin{itemize}
\item À la fin de cette leçon, je sais identifier la structure d'un texte argumentatif simple en chinois (introduction, causes, conclusion).
\item À la fin de cette leçon, je sais utiliser les connecteurs 因为、所以、首先、其次、最后 pour enchaîner mes idées.
\item À la fin de cette leçon, je sais employer le vocabulaire du thème : 腐败、贪污、原因、法律、教育、诚实.
\item À la fin de cette leçon, je sais écrire correctement les caractères complexes du thème en respectant l'ordre des traits et les radicaux.
\item À la fin de cette leçon, je sais distinguer les caractères proches (败/贩, 贪/贫, 律/津) pour éviter les confusions.
\item À la fin de cette leçon, je sais traduire des phrases simples du français vers le chinois (thème) sur la corruption.
\end{itemize}\end{multicols}}

\begin{sommaire}
\begin{multicols}{2}
\begin{enumerate}
\item Modèle de texte commenté
\item Structure et connecteurs utiles
\item Écriture des caractères complexes du thème
\item Caractères proches : ne pas confondre
\item Thème : français → chinois
\item Version : chinois → français
\item Production guidée et grille d'évaluation
\item Activité d'intégration
\item Méthodes et exercices résolus
\item Exercices
\item Corrigés des exercices
\item Fiche bilan
\end{enumerate}
\end{multicols}
\end{sommaire}

\begin{objectifs}
\begin{itemize}
\item À la fin de cette leçon, je sais identifier la structure d'un texte argumentatif simple en chinois (introduction, causes, conclusion).
\item À la fin de cette leçon, je sais utiliser les connecteurs 因为、所以、首先、其次、最后 pour enchaîner mes idées.
\item À la fin de cette leçon, je sais employer le vocabulaire du thème : 腐败、贪污、原因、法律、教育、诚实.
\item À la fin de cette leçon, je sais écrire correctement les caractères complexes du thème en respectant l'ordre des traits et les radicaux.
\item À la fin de cette leçon, je sais distinguer les caractères proches (败/贩, 贪/贫, 律/津) pour éviter les confusions.
\item À la fin de cette leçon, je sais traduire des phrases simples du français vers le chinois (thème) sur la corruption.
\item À la fin de cette leçon, je sais traduire un court texte du chinois vers le français (version).
\item À la fin de cette leçon, je sais produire un texte de 5 à 8 phrases sur les facteurs de la corruption et les moyens de la combattre.
\end{itemize}
\end{objectifs}

\begin{prerequis}
\begin{itemize}
\item Structure de base de la phrase chinoise Sujet + Verbe + Complément (Seconde)
\item Emploi de 是、有、很 et de la négation 不 / 没有
\item Connecteurs simples déjà vus : 和、但是、因为……所以……
\item Radicaux fréquents et ordre général des traits (de haut en bas, de gauche à droite)
\item Vocabulaire de la citoyenneté vu dans le Module 3 (社会、国家、人民)
\end{itemize}
\end{prerequis}

\newpage

\coursec{Modèle de texte commenté}

\courssub{Situation d'accroche : un courriel de Shanghai}

À la sortie du lycée, un élève de Première A de Yaoundé entend à la radio un débat animé sur la lutte contre la corruption au Cameroun. Le soir même, il ouvre sa boîte mail et découvre un message de son correspondant chinois de Shanghai : « Dans mon pays, on parle beaucoup de la lutte contre la corruption. Et dans le tien ? Quelles sont, selon toi, les causes de la corruption ? »

Pour répondre à cette question en chinois, il ne suffit pas de connaître du vocabulaire. Il faut savoir \cle{organiser un texte argumentatif simple} : présenter le problème, donner des causes dans un ordre logique, proposer une solution. Il faut aussi maîtriser les \cle{connecteurs} (首先、其次、最后、因为、所以、为了) et écrire correctement des caractères complexes comme 腐, 败 ou 廉. C'est exactement le défi de cette leçon : rédiger un court texte clair et correct sur un sujet de citoyenneté.

\textbf{Problématique :} comment structurer en chinois un texte simple de type « problème → causes → solution » ?

\courssub{Lecture du texte modèle : 腐败的原因}

Lisons d'abord le texte modèle. Il est court (huit phrases), écrit dans un chinois simple de niveau HSK 2-3, et il servira de \cle{modèle} pour ta propre production.

\begin{textebox}[腐败的原因 — Les causes de la corruption]
腐败是一个严重的问题。\\
Fǔbài shì yí ge yánzhòng de wèntí.\\
La corruption est un problème grave.\\[4pt]
首先，因为一些人想很快地有钱，所以他们贪污。\\
Shǒuxiān, yīnwèi yìxiē rén xiǎng hěn kuài de yǒu qián, suǒyǐ tāmen tānwū.\\
D'abord, parce que certaines personnes veulent avoir de l'argent très vite, elles sont corrompues.\\[4pt]
其次，法律不够严格。\\
Qícì, fǎlǜ bù gòu yángé.\\
Ensuite, les lois ne sont pas assez strictes.\\[4pt]
还有，一些人不知道诚实很重要。\\
Hái yǒu, yìxiē rén bù zhīdào chéngshí hěn zhòngyào.\\
De plus, certaines personnes ne savent pas que l'honnêteté est très importante.\\[4pt]
腐败让国家变穷，也让人民不相信政府。\\
Fǔbài ràng guójiā biàn qióng, yě ràng rénmín bù xiāngxìn zhèngfǔ.\\
La corruption appauvrit le pays et fait que le peuple ne fait plus confiance au gouvernement.\\[4pt]
最后，为了解决这个问题，我们应该加强教育，做一个诚实的人。\\
Zuìhòu, wèile jiějué zhège wèntí, wǒmen yīnggāi jiāqiáng jiàoyù, zuò yí ge chéngshí de rén.\\
Enfin, pour résoudre ce problème, nous devons renforcer l'éducation et être des personnes honnêtes.\\[4pt]
如果每个人都诚实，我们的国家会越来越好。\\
Rúguǒ měi ge rén dōu chéngshí, wǒmen de guójiā huì yuè lái yuè hǎo.\\
Si chacun est honnête, notre pays ira de mieux en mieux.\\[4pt]
我相信明天会更好。\\
Wǒ xiāngxìn míngtiān huì gèng hǎo.\\
Je crois que demain sera meilleur.
\end{textebox}

\textbf{Glossaire du texte}

\begin{center}
\begin{tabularx}{\linewidth}{|X|X|X|X|}
\hline
\rowcolor{popblueL} Caractères & Pinyin & Nature & Traduction \\
\hline
腐败 & fǔbài & nom, verbe, adj. & corruption ; corrompre ; corrompu \\
\hline
严重 & yánzhòng & adj. & grave, sérieux \\
\hline
问题 & wèntí & nom & problème, question \\
\hline
贪污 & tānwū & verbe & être corrompu, détourner de l'argent public \\
\hline
法律 & fǎlǜ & nom & loi, législation \\
\hline
严格 & yángé & adj. & strict, rigoureux \\
\hline
诚实 & chéngshí & adj. & honnête, intègre \\
\hline
政府 & zhèngfǔ & nom & gouvernement \\
\hline
解决 & jiějué & verbe & résoudre \\
\hline
加强 & jiāqiáng & verbe & renforcer \\
\hline
相信 & xiāngxìn & verbe & croire, faire confiance \\
\hline
越来越 & yuè lái yuè & structure & de plus en plus \\
\hline
\end{tabularx}
\end{center}

\courssub{Repérage de la structure en trois parties}

Un bon texte argumentatif simple en chinois suit presque toujours le même plan en trois parties, que les Chinois appellent souvent 总—分—总 (zǒng — fēn — zǒng, « général — détaillé — général ») :

\begin{enumerate}
\item \cle{引言} (yǐnyán, introduction) : on présente le problème en une phrase. Ici : 腐败是一个严重的问题。
\item \cle{原因} (yuányīn, les causes) : on donne deux ou trois causes, introduites par 首先、其次、还有. Ici : l'envie d'argent rapide, des lois pas assez strictes, l'ignorance de l'importance de l'honnêteté. On peut ajouter les conséquences (结果) : le pays s'appauvrit, le peuple perd confiance.
\item \cle{结论} (jiélùn, conclusion) : on propose une solution avec 为了…应该…, puis on termine par un espoir (如果…会… / 我相信…).
\end{enumerate}

\begin{popfigure}
\begin{center}
\begin{tikzpicture}[node distance=0.9cm,
boite/.style={rectangle, rounded corners, draw=popblue, fill=popblueL, text width=5.5cm, align=center, minimum height=1cm, font=\small},
fleche/.style={-{Stealth[length=3mm]}, thick, popdark}]
\node[boite, align=center] (intro){\textbf{引言} yǐnyán\\Présenter le problème :\\腐败是一个严重的问题。};
\node[boite, below=of intro, fill=poporangeL, draw=poporange, align=center] (cause1){\textbf{原因一} yuányīn yī\\首先，因为…所以…\\(l'argent facile)};
\node[boite, below=of cause1, fill=poporangeL, draw=poporange, align=center] (cause2){\textbf{原因二} yuányīn èr\\其次，法律不够严格。\\(lois pas assez strictes)};
\node[boite, below=of cause2, fill=popgreenL, draw=popgreen, align=center] (conclu){\textbf{Conclusion} jiélùn\\最后，为了…我们应该…\\(solution + espoir)};
\draw[fleche] (intro) -- (cause1);
\draw[fleche] (cause1) -- (cause2);
\draw[fleche] (cause2) -- (conclu);
\end{tikzpicture}
\end{center}
\legende{Organigramme du texte modèle : introduction, causes, conclusion.}
\end{popfigure}

\courssub{Observation du ton : phrases courtes, sujet clair}

Observe le texte modèle : chaque phrase est courte (6 à 15 caractères), le sujet est toujours clairement exprimé (一些人、法律、我们、每个人), et il n'y a aucune subordonnée compliquée. C'est exactement le style attendu au lycée :

\begin{exemplebox}[Deux phrases modèles à imiter]
\zh{腐败是一个严重的问题。} \emph{Fǔbài shì yí ge yánzhòng de wèntí.} « La corruption est un problème grave. » — Phrase d'introduction type : sujet + 是 + nom. Simple, directe, efficace.

\zh{为了解决这个问题，我们应该加强教育。} \emph{Wèile jiějué zhège wèntí, wǒmen yīnggāi jiāqiáng jiàoyù.} « Pour résoudre ce problème, nous devons renforcer l'éducation. » — Phrase de conclusion type : 为了 + but，sujet + 应该 + verbe.
\end{exemplebox}

\begin{methode}[Comment exploiter un texte modèle]
\begin{enumerate}
\item \textbf{Lire deux fois} : une première fois pour le sens général, une deuxième fois en traduisant mentalement chaque phrase.
\item \textbf{Encadrer les connecteurs} : surligne 首先、其次、还有、最后、因为、所以、为了. Ce sont les piliers du texte.
\item \textbf{Numéroter les idées} : à côté de chaque phrase, note sa fonction (problème ? cause ? conséquence ? solution ? espoir ?).
\item \textbf{Recopier le plan avant d'écrire} : écris sur ton brouillon 引言 → 原因一 → 原因二 → 结论, puis remplis chaque case avec tes propres idées.
\end{enumerate}
\end{methode}

\begin{exoresolu}[Reformuler le texte en trois phrases]
\textbf{Énoncé.} Résume le texte modèle en trois phrases avec tes propres mots : une phrase pour le problème, une pour les causes, une pour la solution.

\tcblower
\textbf{Solution détaillée.}

1. Problème : \zh{腐败是一个严重的问题，很多国家都有这个问题。} \emph{Fǔbài shì yí ge yánzhòng de wèntí, hěn duō guójiā dōu yǒu zhège wèntí.} « La corruption est un problème grave, beaucoup de pays ont ce problème. »

2. Causes : \zh{因为一些人想很快地有钱，法律也不够严格，所以有腐败。} \emph{Yīnwèi yìxiē rén xiǎng hěn kuài de yǒu qián, fǎlǜ yě bù gòu yángé, suǒyǐ yǒu fǔbài.} « Parce que certaines personnes veulent s'enrichir vite et que les lois ne sont pas assez strictes, il y a de la corruption. »

3. Solution : \zh{为了解决这个问题，我们应该加强教育，做诚实的人。} \emph{Wèile jiějué zhège wèntí, wǒmen yīnggāi jiāqiáng jiàoyù, zuò chéngshí de rén.} « Pour résoudre ce problème, nous devons renforcer l'éducation et être honnêtes. »

\textbf{Méthode} : garde les connecteurs du modèle (因为…所以…, 为了…应该…), mais change les détails pour montrer que tu as compris et que tu sais réutiliser les structures.
\end{exoresolu}

\begin{savaistu}[反腐倡廉 : un grand slogan chinois]
En Chine, la lutte contre la corruption, \zh{反腐败} \emph{fǎn fǔbài}, est un très grand thème de société. On la résume souvent par le slogan \zh{反腐倡廉} \emph{fǎn fǔ chàng lián}, « combattre la corruption et promouvoir l'intégrité ». Le caractère \zh{廉} (lián, « intègre, honnête ») est un caractère complexe et ancien : il apparaît déjà dans les textes classiques pour décrire le fonctionnaire modèle. Au Cameroun comme en Chine, l'honnêteté des citoyens et des responsables est présentée comme la base du développement du pays : un beau point commun à exploiter dans ton texte !
\end{savaistu}

\begin{aretenir}[L'essentiel de cette section]
\begin{itemize}
\item Un texte argumentatif simple en chinois suit le plan : \cle{引言} (problème) → \cle{原因} (deux ou trois causes) → \cle{结论} (solution + espoir).
\item Les connecteurs obligatoires : 首先、其次、还有、最后 pour ordonner ; 因为…所以… pour la cause ; 为了…应该… pour le but et la solution ; 如果…会… pour l'espoir.
\item Style attendu : phrases courtes, sujet clair, pas de subordonnées compliquées.
\item Avant d'écrire : lis le modèle deux fois, encadre les connecteurs, numérote les idées, recopie le plan.
\end{itemize}
\end{aretenir}

\coursec{Structure et connecteurs utiles}

\courssub{Les connecteurs du texte : les articulations du raisonnement}

Relis le texte et souligne mentalement les mots qui relient les phrases. Ce sont les \cle{connecteurs}. Sans eux, le texte est une suite de phrases coupées ; avec eux, le raisonnement avance.

\begin{center}
\begin{tabularx}{\linewidth}{|X|X|X|}
\hline
\rowcolor{popblueL} Connecteur & Fonction & Exemple dans le texte \\
\hline
首先 shǒuxiān & première idée & 首先，因为一些人想很快地有钱… \\
\hline
其次 qícì & idée suivante & 其次，法律不够严格。 \\
\hline
还有 hái yǒu & ajout & 还有，一些人不知道诚实很重要。 \\
\hline
最后 zuìhòu & dernière idée & 最后，为了解决这个问题… \\
\hline
因为…所以… yīnwèi… suǒyǐ… & cause → conséquence & 因为一些人想很快地有钱，所以他们贪污。 \\
\hline
为了 wèile & but & 为了解决这个问题，我们应该加强教育。 \\
\hline
如果…会… rúguǒ… huì… & hypothèse & 如果每个人都诚实，国家会越来越好。 \\
\hline
\end{tabularx}
\end{center}

\begin{propriete}[因为…所以… et 为了]
\textbf{1. 因为 A，所以 B。} = « parce que A, donc B ». Contrairement au français, où l'on évite « parce que… donc… » dans la même phrase (pléonasme), le chinois \emph{aime} cette double marque : 因为 ouvre la cause, 所以 annonce le résultat. On peut garder les deux ou n'en garder qu'un, mais jamais traduire mot à mot « parce que… donc… » en français.

\zh{因为法律不够严格，所以有人贪污。} \emph{Yīnwèi fǎlǜ bù gòu yángé, suǒyǐ yǒu rén tānwū.} « Parce que les lois ne sont pas assez strictes, des gens sont corrompus. »

\textbf{2. 为了 + but，sujet + 应该/要 + verbe。} = « pour + infinitif ». 为了 se place en tête de phrase, suivi du but (un verbe ou un groupe verbal).

\zh{为了解决这个问题，我们应该加强教育。} \emph{Wèile jiějué zhège wèntí, wǒmen yīnggāi jiāqiáng jiàoyù.} « Pour résoudre ce problème, nous devons renforcer l'éducation. »
\end{propriete}

\begin{attention}[Pièges fréquents des francophones]
\begin{itemize}
\item Ne traduis pas « parce que… donc… » en français : c'est lourd. En chinois, 因为…所以… est naturel ; en français, choisis l'un des deux.
\item 首先、其次、最后 se placent en \emph{début de phrase}, suivis d'une virgule chinoise （，）. Ne les mets pas au milieu de la phrase.
\item 为了 exprime le \emph{but}, pas la cause. Pour « à cause de », le chinois utilise 因为. Ne confonds pas : 为了钱 (pour l'argent, but) ≠ 因为钱 (à cause de l'argent, cause).
\item Après 应该 (« devoir »), mets directement le verbe, sans 是 : 我们应该加强教育 (et non 我们应该 \emph{是} 加强教育).
\item \textbf{Sandhi des tons} : 不 (bù) devient \emph{bú} devant un 4\`eme ton (ex: \emph{bú gòu} 不够) ; 一 (yī) devient \emph{yí} devant un 4\`eme ton (ex: \emph{yí ge} 一个) et \emph{yì} devant les autres tons (ex: \emph{yìxiē} 一些).
\end{itemize}
\end{attention}

\begin{popfigure}
\begin{center}
\begin{tikzpicture}[mindmap, grow cyclic, every node/.style={concept, align=center, font=\small},
root concept/.append style={concept color=popblue, text=white, font=\small\bfseries},
level 1/.append style={level distance=3.2cm, sibling angle=120},
level 2/.append style={level distance=2.4cm, sibling angle=50, font=\footnotesize}]
\node[root concept, align=center]{腐败\\fǔbài}
  child[concept color=poporange]{ node[align=center]{原因\\yuányīn\\causes}
    child { node[align=center]{想很快地有钱\\argent facile} }
    child { node[align=center]{法律不够严格\\lois laxistes} }
    child { node[align=center]{不懂诚实\\ignorance} }
  }
  child[concept color=poppink]{ node[align=center]{结果\\jiéguǒ\\conséquences}
    child { node[align=center]{国家变穷\\pays appauvri} }
    child { node[align=center]{不相信政府\\méfiance} }
  }
  child[concept color=popgreen]{ node[align=center]{Solutions\\bànfǎ}
    child { node[align=center]{加强教育\\éducation} }
    child { node[align=center]{做诚实的人\\être honnête} }
    child { node[align=center]{加强监督\\contrôle} }
  };
\end{tikzpicture}
\end{center}
\legende{Carte mentale du thème : causes, conséquences et solutions de la corruption.}
\end{popfigure}

\coursec{Écriture des caractères complexes du thème}

\courssub{Analyse graphique : radicaux, phonétiques et ordre des traits}

Les mots-clés du thème comportent des caractères à structure phono-idéographique (forme + son) ou idéographique. Les maîtriser permet d'écrire vite et juste, et de deviner le sens de mots nouveaux.

\begin{center}
\begin{tabularx}{\linewidth}{|X|X|X|X|X|}
\hline
\rowcolor{popblueL} Caractère & Pinyin & Radical (clé) & Structure & Ordre des traits (résumé) \\
\hline
腐 & fǔ & \zh{月} (viande/charnière, 4 traits) & Gauche-Droite : 月 + 府 (fǔ, phonétique) & 1. 月 (haut→bas, gauche→droite) 2. 府 (大→土→寸) \\
\hline
败 & bài & \zh{卩} (jié, sceau, 2 traits) & Haut-Bas (simplifié) : 卩 + 贝 + 又 & 1. 卩 (haut→bas) 2. 贝 (haut→bas) 3. 又 (horiz.→crochet) \\
\hline
贪 & tān & \zh{贝} (bei, coquillage/argent, 4 traits) & Haut-Bas : 今 (jīn, phonétique) + 贝 & 1. 今 (haut→bas) 2. 贝 (bas) \\
\hline
污 & wū & \zh{氵} (shuǐ, eau, 3 traits) & Gauche-Droite : 氵+ 于 (yú, phonétique) & 1. 氵(2 points + trait montant) 2. 于 (haut→bas) \\
\hline
廉 & lián & \zh{广} (guǎng, toit/maison, 3 traits) & Haut-Bas : 广 + 兼 (jiān, phonétique) & 1. 广 (point, horiz., boucle) 2. 兼 (haut→bas) \\
\hline
严 & yán & \zh{厂} (hǎn, falaise, 2 traits) & Haut-Bas : 厂 + 良 (liáng) → 土+口+小 & 1. 厂 (horiz.→vertical) 2. 内部 (haut→bas) \\
\hline
格 & gé & \zh{木} (mù, bois, 4 traits) & Gauche-Droite : 木 + 各 (gè, phonétique) & 1. 木 (horiz., vert., g., d.) 2. 各 (口→夂) \\
\hline
\end{tabularx}
\end{center}

\begin{definition}[Rappel : les 8 règles de base de l'ordre des traits]
\begin{enumerate}
\item Horizontal avant vertical (十)
\item Gauche avant droite (人)
\item Haut avant bas (三)
\item Extérieur avant intérieur (月)
\item Fermer en dernier (回)
\item Trait central avant les ailes (小)
\item Point en haut à droite à la fin (瓦)
\item Trait traversant en dernier (中)
\end{enumerate}
\end{definition}

\courssub{Entraînement guidé : décomposition et recomposition}

\begin{experience}[Atelier « Caractères en puzzle »]
En binôme. L'élève A dicte l'ordre des composants (ex: « 月, puis 府 »), l'élève B écrit le caractère sur ardoise sans voir le modèle. Inverser les rôles. Vérifier ensemble : radical correct ? ordre respecté ? proportion équilibrée (le radical 月 est étroit, 府 large) ?
\end{experience}

\begin{exercice}[Écriture dictée]
Écris chaque caractère 5 fois en respectant l'ordre des traits, puis écris le mot complet 3 fois.
\begin{enumerate}
\item 腐败 (fǔbài)
\item 贪污 (tānwū)
\item 诚实 (chéngshí)
\item 廉洁 (liánjié) — \emph{intégrité, probité}
\item 严格 (yángé)
\item 监督 (jiāndū) — \emph{superviser, contrôle} (vocabulaire d'extension)
\end{enumerate}
\end{exercice}

\coursec{Caractères proches : ne pas confondre}

\courssub{Pièges visuels et phonétiques}

Le chinois regorge de caractères qui se ressemblent (même composants, ordre différent) ou qui sonnent pareil (homophones). Les distinguer est crucial pour l'expression écrite notée.

\begin{center}
\begin{tabularx}{\linewidth}{|X|X|X|X|}
\hline
\rowcolor{popblueL} Caractère cible & Caractère confus & Différence clé & Astuce mnémotechnique \\
\hline
腐 (fǔ, pourrir) & 腹 (fù, ventre/abdomen) & Droite : 府 (fu) vs 复 (fù) & \zh{腐} : la viande (\zh{月}) dans le magasin (\zh{府}) pourrit. \zh{腹} : la viande (\zh{月}) est complexe (\zh{复}) à digérer. \\
\hline
败 (bài, échouer) & 拜 (bài, saluer) & Haut : 卩/贝/又 vs 手 (main) & \zh{败} : le sceau (\zh{卩}) et l'argent (\zh{贝}) sont perdus (\zh{又} main qui lâche). \zh{拜} : les deux mains (\zh{手}) au sol. \\
\hline
贪 (tān, avide) & 谈 (tán, discuter) & Gauche : 贝 (argent) vs 讠 (parole) & \zh{贪} : vouloir l'argent (\zh{贝}) maintenant (\zh{今}). \zh{谈} : parole (\zh{讠}) + feu (\zh{炎}) = discussion animée. \\
\hline
污 (wū, sale) & 乌 (wū, corbeau/noir) & Gauche : 氵(eau) vs rien & \zh{污} : l'eau (\zh{氵}) devient noire (\zh{于} ≈ 乌). \zh{乌} : oiseau sans œil (noir). \\
\hline
廉 (lián, intègre) & 兼 (jiān, cumuler) & Haut : 广 (toit) vs 兼 (tout le caractère) & \zh{廉} : sous le toit (\zh{广}), on cumule (\zh{兼}) les vertus. \zh{兼} : deux épis (\zh{禾}) en main (\zh{兼}). \\
\hline
严 (yán, strict) & 研 (yán, rechercher) & Bas : 土+口+小 vs 石 (pierre) & \zh{严} : la falaise (\zh{厂}) sur la terre (\zh{土}) bouche (\zh{口}) le petit (\zh{小}). \zh{研} : broyer la pierre (\zh{石}) avec un mortier. \\
\hline
格 (gé, norme) & 核 (hé, noyau) & Droite : 各 (gè) vs 亥 (hài) & \zh{格} : le bois (\zh{木}) arrive (\zh{各}) à maturité. \zh{核} : le noyau (\zh{亥}) dans le bois (\zh{木}). \\
\hline
\end{tabularx}
\end{center}

\begin{attention}[Erreurs fatales en examen]
\begin{itemize}
\item Écrire \zh{腐败} avec \zh{腹} ou \zh{拜} → 0 point au mot.
\item Confondre \zh{贪} et \zh{谈} change le sens : « il est avide » vs « il discute ».
\item Oublier le trait d'eau \zh{氵} à \zh{污} → caractère inexistant.
\item Inverser l'ordre des traits (ex: écrire le phonétique avant le radical) rend l'écriture lente et déformée.
\end{itemize}
\end{attention}

\begin{exoresolu}[Repérage d'erreurs]
\textbf{Énoncé.} Chaque ligne contient une erreur de caractère. Identifie-la, corrige-la et explique la confusion.

\tcblower
\textbf{Solution détaillée.}

1. \zh{因为腹败，国家变穷。} → \textbf{Correction} : \zh{因为腐败…} \emph{(腐 fǔ « pourrir » ≠ 腹 fù « ventre »)}.
2. \zh{他很贪污，想和别人谈。} → \textbf{Correction} : \zh{…想和别人贪。} (absurde) ou phrase originale : \zh{他很贪污，想贪钱。} \emph{Le contexte « argent » impose 贪 (argent 贝), pas 谈 (parole 讠).}
3. \zh{我们要加强监独。} → \textbf{Correction} : \zh{我们要加强监督。} \emph{督 (dū, superviser) ≠ 独 (dú, seul). 督 = 目 (œil) + 叔 ; 独 = 虫 (insecte) + 犬.}
4. \zh{廉洁的官员不研刻。} → \textbf{Correction} : \zh{…不严刻。} \emph{严 (yán, strict) ≠ 研 (yán, rechercher). « Strict et incisif » = 严刻.}
\end{exoresolu}

\coursec{Thème : français → chinois}

\courssub{Méthode de traduction (thème)}

\begin{methode}[5 étapes pour un thème réussi]
\begin{enumerate}
\item \textbf{Analyser} : identifier le sujet, le verbe, les compléments, le temps/aspect.
\item \textbf{Structurer} : ordonner les mots selon la syntaxe chinoise (Sujet + Temps + Lieu + Manière + Verbe + Complément).
\item \textbf{Lexicaliser} : choisir le mot chinois précis (attention aux classificateurs, aux verbes d'action vs état).
\item \textbf{Grammaticaliser} : ajouter les mots-outils (的, 得, 地, 了, 过, 把, 被, 比, 连…).
\item \textbf{Relire} : vérifier l'ordre des mots, les tons (pinyin), l'écriture des caractères, la ponctuation chinoise.
\end{enumerate}
\end{methode}

\courssub{Exercices de thème progressifs}

\begin{exercice}[Phrases isolées — Niveau HSK 2-3]
Traduis en chinois (caractères + pinyin).
\begin{enumerate}
\item La corruption est un problème grave dans beaucoup de pays.
\item Parce que certains fonctionnaires veulent de l'argent facile, ils sont corrompus.
\item Les lois ne sont pas assez strictes, donc le peuple ne fait pas confiance au gouvernement.
\item Pour résoudre ce problème, nous devons renforcer l'éducation et promouvoir l'honnêteté.
\item Si chacun est honnête, le pays ira de mieux en mieux.
\end{enumerate}
\end{exercice}

\begin{corrige}[Exercice Thème 1]
\begin{enumerate}
\item \zh{腐败在很多国家是一个严重的问题。} \emph{Fǔbài zài hěn duō guójiā shì yí ge yánzhòng de wèntí.} (Structure : Sujet + 在 + Lieu + 是 + Nom). \textbf{Note} : \zh{在很多国家} placé avant le verbe \zh{是}.
\item \zh{因为一些官员想要容易的钱，所以他们贪污。} \emph{Yīnwèi yìxiē guānyuán xiǎng yào róngyì de qián, suǒyǐ tāmen tānwū.} \textbf{Vocab} : 官员 (guānyuán, fonctionnaire), 容易 (róngyì, facile). \textbf{Grammaire} : 因为…所以… ; 想要 + nom.
\item \zh{法律不够严格，所以人民不相信政府。} \emph{Fǎlǜ bù gòu yángé, suǒyǐ rénmín bù xiāngxìn zhèngfǔ.} \textbf{Attention} : « donc » = 所以 (souvent après une virgule, sans 因为 explicite si le contexte est clair). \zh{不够} (bù gòu) = pas assez.
\item \zh{为了解决这个问题，我们应该加强教育，提倡诚实。} \emph{Wèile jiějué zhège wèntí, wǒmen yīnggāi jiāqiáng jiàoyù, tíchàng chéngshí.} \textbf{Vocab} : 提倡 (tíchàng, promouvoir/préconiser). \textbf{Structure} : 为了…，应该… ; coordination par virgule.
\item \zh{如果每个人都诚实，国家会越来越好。} \emph{Rúguǒ měi ge rén dōu chéngshí, guójiā huì yuè lái yuè hǎo.} \textbf{Grammaire} : 都 (dōu) après 每个人 pour « chacun » ; 越来越 + adj.
\end{enumerate}
\end{corrige}

\begin{exercice}[Texte court — Niveau HSK 3-4]
Traduis le paragraphe suivant en un texte chinois cohérent (environ 80-100 caractères). Utilise les connecteurs vus en section 2.

« Au Cameroun, la corruption est un problème sérieux. D'abord, certains veulent s'enrichir vite. Ensuite, le contrôle n'est pas assez strict. Enfin, pour changer la situation, nous devons renforcer la loi et l'éducation. Je crois que l'avenir sera meilleur. »
\end{exercice}

\begin{corrige}[Exercice Thème 2]
\zh{在喀麦隆，腐败是一个严重的问题。首先，一些人想很快发财。其次，监督不够严格。最后，为了改变这种情况，我们应该加强法治和教育。我相信未来会更好。}

\emph{Zài Kāmàilóng, fǔbài shì yí ge yánzhòng de wèntí. Shǒuxiān, yìxiē rén xiǎng hěn kuài fācái. Qícì, jiāndū bù gòu yángé. Zuìhòu, wèile gǎibiàn zhè zhǒng qíngkuàng, wǒmen yīnggāi jiāqiáng fǎzhì hé jiàoyù. Wǒ xiāngxìn wèilái huì gèng hǎo.}

\textbf{Points clés} :
\begin{itemize}
\item \zh{发财} (fācái, s'enrichir/devenir riche) plus naturel que \zh{有钱} pour « s'enrichir vite ».
\item \zh{监督} (jiāndū, supervision/contrôle) vocabulaire d'extension précis.
\item \zh{这种情况} (zhè zhǒng qíngkuàng, cette situation) pour « la situation ».
\item \zh{法治} (fǎzhì, État de droit / règle de droit) terme institutionnel précis.
\item Ponctuation : virgules chinoises ，, point final 。.
\end{itemize}
\end{corrige}

\coursec{Version : chinois → français}

\courssub{Méthode de traduction (version)}

\begin{methode}[4 étapes pour une version réussie]
\begin{enumerate}
\item \textbf{Segmenter} : repérer les limites de phrases (。), les groupes nominaux (的), les structures (因为…所以…, 为了…).
\item \textbf{Glosser} : écrire le sens littéral mot à mot au-dessus du texte chinois.
\item \textbf{Restructurer} : réordonner selon la syntaxe française (SVO, subordination, prépositions).
\item \textbf{Polir} : choisir le registre adapté (langage soutenu pour un texte civique), vérifier l'orthographe et la ponctuation françaises.
\end{enumerate}
\end{methode}

\courssub{Exercices de version progressifs}

\begin{exercice}[Phrases isolées]
Traduis en français.
\begin{enumerate}
\item 腐败现象很严重，影响国家发展。
\item 因为监督不力，所以有人敢贪污。
\item 为了建设廉洁社会，我们必须加强法治教育。
\item 如果人人都守法，社会就会公平正义。
\end{enumerate}
\end{exercice}

\begin{corrige}[Exercice Version 1]
\begin{enumerate}
\item \textbf{La corruption est un phénomène grave, qui affecte le développement du pays.} (Structure : Sujet 腐败现象 + Verbe 很严重 ; 影响无主语, sujet implicite = la corruption).
\item \textbf{Parce que la supervision est inefficace, certains osent être corrompus.} (\zh{监督不力} jiāndū bùlì = supervision inefficace/faible ; \zh{敢} gǎn = oser).
\item \textbf{Pour construire une société intègre, nous devons renforcer l'éducation à l'État de droit.} (\zh{建设} jiànshè = construire/édifier ; \zh{廉洁} liánjié = intègre/probe ; \zh{必须} bìxū = devoir/falloir (obligation forte) ; \zh{法治教育} fǎzhì jiàoyù = éducation au droit/à l'État de droit).
\item \textbf{Si tout le monde respecte la loi, la société sera juste et équitable.} (\zh{人人} rénrén = tout le monde ; \zh{守法} shǒufǎ = respecter la loi ; \zh{公平正义} gōngpíng zhèngyì = justice et équité / juste et équitable).
\end{enumerate}
\end{corrige}

\begin{exercice}[Texte suivi — Niveau HSK 3]
Traduis le texte suivant en français soigné.

\begin{textebox}[Texte à traduire]
反腐倡廉是长期的任务。\\
Fǎn fǔ chàng lián shì chángqí de rènwù.\\
首先，要完善法律制度，让贪污者无处藏身。\\
Shǒuxiān, yào wánshàn fǎlǜ zhìdù, ràng tānwū zhě wú chù cángshēn.\\
其次，要加强道德教育，从小培养诚实守信的品质。\\
Qícì, yào jiāqiáng dàodé jiàoyù, cóng xiǎo péiyǎng chéngshí shǒuxìn de pǐnzhì.\\
最后，每个公民都要监督政府，不敢腐，不能腐，不想腐。\\
Zuìhòu, měi ge gōngmín dōu yào jiāndū zhèngfǔ, bù gǎn fǔ, bù néng fǔ, bù xiǎng fǔ.\\
只有这样，国家才能长治久安。\\
Zhèyàng zhǐ, guójiā cái néng chángzhìjiǔ'ān.
\end{textebox}
\end{exercice}

\begin{corrige}[Exercice Version 2]
\textbf{Traduction proposée :}

« Lutter contre la corruption et promouvoir l'intégrité est une tâche de long terme.
D'abord, il faut perfectionner le système juridique, pour que les corrompus n'aient nulle part où se cacher.
Ensuite, il faut renforcer l'éducation morale, cultiver dès le plus jeune âge la qualité d'honnêteté et de fiabilité.
Enfin, chaque citoyen doit superviser le gouvernement, [faire en sorte qu'on] n'ose pas être corrompu, qu'on ne puisse pas l'être, qu'on ne veuille pas l'être.
Seulement ainsi le pays pourra connaître une stabilité durable. »

\textbf{Notes linguistiques :}
\begin{itemize}
\item \zh{长期的任务} (chángqí de rènwù) : tâche à long terme / de longue haleine.
\item \zh{完善} (wánshàn) : perfectionner, améliorer, compléter.
\item \zh{法律制度} (fǎlǜ zhìdù) : système juridique / institutionnel.
\item \zh{无处藏身} (wú chù cángshēn) : n'avoir nulle part où se cacher (idiome).
\item \zh{贪污者} (tānwū zhě) : les corrompus (suffixe \zh{者} = celui qui...).
\item \zh{道德教育} (dàodé jiàoyù) : éducation morale / civique.
\item \zh{诚实守信} (chéngshí shǒuxìn) : honnête et tenant ses promesses (chengyu 4 caractères).
\item \zh{品质} (pǐnzhì) : qualité morale.
\item \zh{不敢腐，不能腐，不想腐} (bù gǎn fǔ, bù néng fǔ, bù xiǎng fǔ) : les « trois ne pas » (三不) de la stratégie anti-corruption chinoise (Xi Jinping) : ne pas oser (dissuasion), ne pas pouvoir (institutionnel), ne pas vouloir (idéologique). \zh{腐} abrège \zh{腐败} ici.
\item \zh{只有…才…} (zhǐyǒu… cái…) : seulement si… alors… (condition nécessaire).
\item \zh{长治久安} (chángzhìjiǔ'ān) : paix et stabilité durables (chengyu).
\end{itemize}
\end{corrige}

\coursec{Production guidée et grille d'évaluation}

\courssub{Tâche finale : Répondre à ton correspondant de Shanghai}

\begin{experience}[Production écrite : Courriel à un ami chinois]
\textbf{Consigne.} Tu réponds à l'email de ton correspondant (voir situation d'accroche). Rédige un message de \textbf{80 à 100 caractères chinois} (ponctuation comprise) structuré en 3 paragraphes :
\begin{enumerate}
\item \textbf{Problème} : La corruption est grave au Cameroun aussi.
\item \textbf{Causes} : Donne 2 causes (ex: argent facile, lois laxistes, manque d'éducation civique, pauvreté). Utilise 首先、其次.
\item \textbf{Solution + Espoir} : Propose une action (éducation, loi, contrôle) avec 为了…应该… Termine par 我相信… ou 如果…会…
\end{enumerate}
\textbf{Contraintes} : Écriture soignée (caractères vus en section 3), pinyin au brouillon, pas de traducteur automatique.
\end{experience}

\begin{exemplebox}[Modèle de production élève (≈ 90 caractères)]
\zh{你好！喀麦隆的腐败也很严重。首先，因为一些人想快点发财，所以贪污。其次，法律不够严格，监督也不够。为了解决这个问题，我们应该加强法治教育，做诚实的人。如果大家都诚实，喀麦隆会越来越好。我相信明天更好！祝好。}

\emph{Nǐ hǎo! Kāmàilóng de fǔbài yě hěn yánzhòng. Shǒuxiān, yīnwèi yìxiē rén xiǎng kuài diǎn fācái, suǒyǐ tānwū. Qícì, fǎlǜ bù gòu yángé, jiāndū yě bù gòu. Wèile jiějué zhège wèntí, wǒmen yīnggāi jiāqiáng fǎzhì jiàoyù, zuò chéngshí de rén. Rúguǒ dàjiā dōu chéngshí, Kāmàilóng huì yuè lái yuè hǎo. Wǒ xiāngxìn míngtiān gèng hǎo! Zhù hǎo.}
\end{exemplebox}

\courssub{Grille d'évaluation (sur 20 points) — Expression écrite}

\begin{center}
\begin{tabularx}{\linewidth}{|X|c|X|}
\hline
\rowcolor{popblueL} \textbf{Critère} & \textbf{Points} & \textbf{Indicateurs de réussite} \\
\hline
\textbf{Respect de la consigne \& Structure} & 4 & 3 parties visibles (Problème, Causes, Solution) ; connecteurs 首先/其次/最后 utilisés correctement ; longueur 80-100 caractères. \\
\hline
\textbf{Vocabulaire du thème} & 4 & Utilisation correcte d'au moins 5 mots-clés : 腐败、贪污、法律、严格、诚实、监督、法治、解决、加强、相信. Orthographe des caractères exacte (pas de confusion section 4). \\
\hline
\textbf{Grammaire \& Syntaxe} & 5 & Structures cibles maîtrisées : 因为…所以… ; 为了…应该/要… ; 如果…会… ; 不够 + adj. ; 越来越 + adj. ; ordre des mots correct (Sujet - Temps - Lieu - Verbe). Pas de \zh{是} + verbe. \\
\hline
\textbf{Cohérence \& Ponctuation} & 3 & Progression logique ; ponctuation chinoise (，。！) ; pas de phrases « telegram style » sans verbe. \\
\hline
\textbf{Qualité de l'écriture (Calligraphie/Soin)} & 4 & Caractères lisibles, proportions respectées, ordre des traits correct (radicaux bien placés), pas de traits manquants/superflus. \\
\hline
\textbf{Total} & \textbf{20} & \\
\hline
\end{tabularx}
\end{center}

\begin{aretenir}[Bilan de la leçon 30]
\begin{itemize}
\item \textbf{Savoir} : Vocabulaire corruption/intégrité (腐败, 贪污, 廉洁, 监督, 法治) ; Connecteurs argumentatifs (首先、因为…所以、为了、如果).
\item \textbf{Savoir-faire} : Structurer un texte 总—分—总 ; Écrire 7 caractères complexes (腐, 败, 贪, 污, 廉, 严, 格) sans confusion ; Traduire thème/version niveau HSK 3.
\item \textbf{Savoir-être} : Argumenter sur un sujet civique avec mesure ; Apprécier les valeurs partagées Cameroun/Chine (intégrité, État de droit).
\end{itemize}
\end{aretenir}

\coursec{Structure et connecteurs utiles}

Dans la section précédente, nous avons lu et commenté un texte modèle sur les facteurs de la corruption. Vous avez remarqué que ce texte n'est pas une simple suite de phrases isolées : chaque phrase est reliée à la suivante par de petits mots comme \zh{首先}, \zh{因为}, \zh{所以}, \zh{但是}, \zh{为了}. Ces mots s'appellent des \cle{connecteurs}. C'est eux qui donnent au texte sa logique et sa cohérence. Un texte sans connecteurs est comme une route sans panneaux : le lecteur se perd. Dans cette section, nous allons observer ces connecteurs, apprendre à les employer correctement, puis construire un plan de rédaction en quelques minutes.

\begin{definition}[Le connecteur : 连接词]
Un \cle{connecteur}, en chinois \zh{连接词} (\emph{liánjiēcí}, « mot qui relie »), est un mot ou un groupe de mots qui relie les phrases entre elles et donne de la cohérence au texte. Il indique la relation logique entre les idées : énumération, cause, conséquence, but, opposition.
\end{definition}

\courssub{Observer avant d'apprendre : un texte riche en connecteurs}

Lisons d'abord ce court texte original. Soulignez mentalement tous les connecteurs que vous reconnaissez. Notez la structure \zh{只有……才……} (seulement si… alors…) dans la conclusion, vue en Section 1.

\begin{textebox}[腐败的原因 (Les causes de la corruption)]
腐败是一个严重的问题。为什么有人贪污？首先，因为法律不严格，所以有人敢贪污。其次，因为工资太低，所以一些公务员想多拿钱。最后，因为没有人监督，所以贪污的人不怕。但是，腐败对国家和人民很不好。因此，我们应该想办法。为了解决这个问题，政府应该制定严格的法律，提高公务员的工资，还要加强监督。只有这样，我们的国家才能发展。\\
\emph{Fǔbài shì yí ge yánzhòng de wèntí. Wèishénme yǒu rén tānwū? Shǒuxiān, yīnwèi fǎlǜ bù yángé, suǒyǐ yǒu rén gǎn tānwū. Qícì, yīnwèi gōngzī tài dī, suǒyǐ yìxiē gōngwùyuán xiǎng duō ná qián. Zuìhòu, yīnwèi méiyǒu rén jiāndū, suǒyǐ tānwū de rén bù pà. Dànshì, fǔbài duì guójiā hé rénmín hěn bù hǎo. Yīncǐ, wǒmen yīnggāi xiǎng bànfǎ. Wèile jiějué zhège wèntí, zhèngfǔ yīnggāi zhìdìng yángé de fǎlǜ, tígāo gōngwùyuán de gōngzī, hái yào jiāqiáng jiāndū. Zhǐyǒu zhèyàng, wǒmen de guójiā cái néng fāzhǎn.}\\
« La corruption est un problème grave. Pourquoi certaines personnes pratiquent-elles la corruption ? D'abord, parce que la loi n'est pas stricte, certains osent être corrompus. Ensuite, parce que les salaires sont trop bas, certains fonctionnaires veulent prendre plus d'argent. Enfin, parce que personne ne contrôle, les corrompus n'ont pas peur. Mais la corruption est très mauvaise pour le pays et le peuple. C'est pourquoi nous devons trouver des solutions. Pour résoudre ce problème, le gouvernement doit établir des lois strictes, augmenter les salaires des fonctionnaires et renforcer le contrôle. Seulement ainsi, notre pays pourra se développer. »
\end{textebox}

\textbf{Glossaire}

\begin{center}
\begin{tabularx}{\linewidth}{|X|X|X|X|}
\hline
\rowcolor{popblueL} Caractères & Pinyin & Nature & Traduction \\
\hline
腐败 & fǔbài & nom / adj. & corruption ; corrompu \\
\hline
贪污 & tānwū & verbe & pratiquer la corruption, détourner \\
\hline
严格 & yángé & adj. & strict, rigoureux \\
\hline
工资 & gōngzī & nom & salaire \\
\hline
公务员 & gōngwùyuán & nom & fonctionnaire \\
\hline
监督 & jiāndū & verbe / nom & contrôler ; contrôle, supervision \\
\hline
解决 & jiějué & verbe & résoudre \\
\hline
制定 & zhìdìng & verbe & établir, élaborer (une loi) \\
\hline
提高 & tígāo & verbe & augmenter, élever \\
\hline
加强 & jiāqiáng & verbe & renforcer \\
\hline
\end{tabularx}
\end{center}

\textbf{Questions de compréhension}
\begin{enumerate}
\item Quelles sont les trois causes de la corruption mentionnées dans le texte ?
\item Quel connecteur introduit chaque cause ?
\item Quelles solutions le gouvernement doit-il prendre selon l'auteur ?
\end{enumerate}

\courssub{Les connecteurs d'énumération : 首先……其次……最后……}

\begin{propriete}[Énumérer les idées]
Pour présenter des idées dans un ordre clair (par exemple plusieurs causes), le chinois utilise la série :
\begin{itemize}
\item \zh{首先} (\emph{shǒuxiān}) = d'abord, premièrement ;
\item \zh{其次} (\emph{qícì}) = ensuite, deuxièmement ;
\item \zh{最后} (\emph{zuìhòu}) = enfin, finalement.
\end{itemize}
Structure : \textbf{首先，+ phrase 1。其次，+ phrase 2。最后，+ phrase 3。}\\
Chaque connecteur est placé en début de phrase et suivi d'une virgule chinoise \zh{，}.
\end{propriete}

\begin{exemplebox}[Énumération en action]
\zh{首先，我想谈腐败的原因；其次，我想谈解决的办法；最后，我想谈我们学生的责任。}\\
\emph{Shǒuxiān, wǒ xiǎng tán fǔbài de yuányīn ; qícì, wǒ xiǎng tán jiějué de bànfǎ ; zuìhòu, wǒ xiǎng tán wǒmen xuésheng de zérèn.}\\
« D'abord, je veux parler des causes de la corruption ; ensuite, je veux parler des solutions ; enfin, je veux parler de la responsabilité de nous, les élèves. »
\end{exemplebox}

\begin{attention}[Ne pas confondre 首先 et 第一]
\zh{第一} (\emph{dìyī}) signifie « premier » dans un classement ou une liste numérotée (\zh{第一名} = le premier de la classe). Pour énumérer des arguments dans un texte, utilisez \zh{首先}, pas \zh{第一}. De même, \zh{最后} (enfin, dans un raisonnement) ne se confond pas avec \zh{最近} (\emph{zuìjìn}, « récemment »).
\end{attention}

\courssub{Cause et conséquence : 因为……所以……, 由于, 因此}

\begin{propriete}[La paire 因为……所以……]
Structure : \textbf{因为 + cause，所以 + résultat。}\\
\zh{因为} (\emph{yīnwèi}) = parce que ; \zh{所以} (\emph{suǒyǐ}) = donc, c'est pourquoi.\\
Grande différence avec le français : en français, « parce que… donc… » dans la même phrase est incorrect ; en chinois, on peut — et on préfère souvent — garder les deux mots dans la même phrase. On peut aussi n'en garder qu'un seul : \zh{因为…} seul (en réponse à \zh{为什么}) ou \zh{所以}… seul (en réponse à une cause énoncée juste avant), mais jamais \zh{所以} en tout début de texte sans phrase précédente.
\end{propriete}

\begin{exemplebox}[Cause et conséquence]
\zh{因为法律不严格，所以有人贪污。}\\
\emph{Yīnwèi fǎlǜ bù yángé, suǒyǐ yǒu rén gǎn tānwū.}\\
« Parce que la loi n'est pas stricte, certaines personnes pratiquent la corruption. »\\[4pt]
\zh{因为工资太低，所以他离开了这个公司。}\\
\emph{Yīnwèi gōngzī tài dī, suǒyǐ tā líkāi le zhège gōngsī.}\\
« Comme le salaire était trop bas, il a quitté cette entreprise. »\\[4pt]
\zh{由于下大雨，比赛取消了。}\\
\emph{Yóuyú xià dà yǔ, bǐsài qǔxiāo le.}\\
« En raison de la forte pluie, le match a été annulé. »\\[4pt]
\zh{腐败对国家很不好，因此我们必须反对它。}\\
\emph{Fǔbài duì guójiā hěn bù hǎo, yīncǐ wǒmen bìxū fǎnduì tā.}\\
« La corruption est très mauvaise pour le pays, c'est pourquoi nous devons la combattre. »
\end{exemplebox}

\begin{definition}[由于 et 因此]
\zh{由于} (\emph{yóuyú}) = « en raison de, étant donné », plus formel que \zh{因为} ; il introduit la cause, souvent sans \zh{所以} derrière.\\
\zh{因此} (\emph{yīncǐ}) = « c'est pourquoi, par conséquent » ; il introduit la conséquence et, contrairement à \zh{所以}, il peut commencer une nouvelle phrase après un point.
\end{definition}

\begin{attention}[Le piège du « donc » français]
Erreur très fréquente des francophones : écrire une phrase, mettre un point, puis commencer la phrase suivante par \zh{所以} pour traduire « Donc, … ». En chinois, \zh{所以} doit être rattaché à sa cause dans la même phrase (après la virgule \zh{，}). Pour commencer une nouvelle phrase par « Donc / C'est pourquoi », utilisez \zh{因此}.\\
Incorrect : \zh{法律不严格。所以有人贪污。} (style lourd, à éviter)\\
Correct : \zh{法律不严格，所以有人贪污。} ou \zh{法律不严格。因此，有人贪污。}
\end{attention}

\begin{popfigure}
\begin{tikzpicture}[
box/.style={draw=popblue, fill=popblueL, rounded corners=3pt, align=center, font=\small, inner sep=5pt},
res/.style={draw=poporange, fill=poporangeL, rounded corners=3pt, align=center, font=\small, inner sep=5pt},
lbl/.style={font=\small\bfseries, text=popdark}]
\node[box, align=center] (cause) at (0,0){因为 + cause\\ \zh{因为法律不严格，}};
\node[res, align=center] (result) at (7.2,0){所以 + résultat\\ \zh{所以有人贪污。}};
\draw[-{Stealth[length=3mm]}, very thick, poppurple] (cause) -- (result) node[midway, above, lbl]{la cause entraîne le résultat};
\node[lbl, text=popblue] at (0,1.3) {CAUSE};
\node[lbl, text=poporange] at (7.2,1.3) {RÉSULTAT};
\end{tikzpicture}
\legende{Schéma de la phrase de cause : 因为 + cause，所以 + résultat — l'ordre est le même qu'en français.}
\end{popfigure}

\courssub{Le connecteur de but : 为了 + verbe}

\begin{propriete}[Exprimer le but avec 为了]
Structure : \textbf{为了 + verbe (+ complément)，sujet + verbe …}\\
\zh{为了} (\emph{wèile}) = « pour, afin de ». Il se place en début de phrase, suivi directement d'un verbe, puis d'une virgule, puis de l'action principale.
\end{propriete}

\begin{exemplebox}[Le but]
\zh{为了解决这个问题，政府应该制定严格的法律。}\\
\emph{Wèile jiějué zhège wèntí, zhèngfǔ yīnggāi zhìdìng yángé de fǎlǜ.}\\
« Pour résoudre ce problème, le gouvernement doit établir des lois strictes. »\\[4pt]
\zh{为了学好汉语，我每天练习写汉字。}\\
\emph{Wèile xuéhǎo Hànyǔ, wǒ měitiān liànxí xiě Hànzì.}\\
« Pour bien apprendre le chinois, je m'entraîne chaque jour à écrire les caractères. »
\end{exemplebox}

\begin{attention}[为了 ou 为 ?]
\zh{为了} est suivi d'un verbe ou d'un groupe verbal (\zh{为了学习} = pour étudier). Devant un nom de personne ou un bénéficiaire, on utilise plutôt \zh{为} : \zh{为人民服务} (\emph{wèi rénmín fúwù}) = « servir le peuple ». Ne dites pas \zh{为了人民服务} dans ce sens figé.
\end{attention}

\courssub{L'opposition : 但是 et 可是}

\begin{propriete}[Exprimer l'opposition]
\zh{但是} (\emph{dànshì}) et \zh{可是} (\emph{kěshì}) signifient tous les deux « mais ». \zh{但是} est plus soutenu et fréquent à l'écrit ; \zh{可是} est plus oral. Structure : \textbf{phrase 1，但是 + phrase 2。}\\
Contrairement au français, le chinois peut combiner \zh{虽然} (bien que) et \zh{但是} (mais) dans la même phrase : \zh{虽然……但是……}.
\end{propriete}

\begin{exemplebox}[L'opposition]
\zh{腐败对国家和人民很不好，但是有些人还是贪污。}\\
\emph{Fǔbài duì guójiā hé rénmín hěn bù hǎo, dànshì yǒuxiē rén háishi tānwū.}\\
« La corruption est très mauvaise pour le pays et le peuple, mais certaines personnes continuent quand même. »\\[4pt]
\zh{虽然这个问题很难，但是我们必须解决它。}\\
\emph{Suīrán zhège wèntí hěn nán, dànshì wǒmen bìxū jiějué tā.}\\
« Bien que ce problème soit difficile, nous devons le résoudre. »
\end{exemplebox}

\courssub{Ponctuation chinoise et majuscules : les règles d'or}

\begin{propriete}[，et 。]
\begin{itemize}
\item La virgule chinoise \zh{，} sépare les parties d'une même phrase (cause + résultat, but + action, énumération). Elle est pleine chasse : elle occupe la place d'un caractère.
\item Le point chinois \zh{。} (un petit cercle, pas un point) termine une phrase complète.
\item Il n'y a \textbf{pas de majuscules} en chinois : ni en début de phrase, ni pour les noms propres. \zh{喀麦隆} (Cameroun) s'écrit comme n'importe quel mot.
\item Le point-virgule chinois \zh{；} sépare des idées parallèles dans une énumération longue.
\end{itemize}
\end{propriete}

\courssub{Tableau récapitulatif des connecteurs}

\begin{center}
\begin{tabularx}{\linewidth}{|l|X|X|X|}
\hline
\rowcolor{popblueL} Fonction & Connecteur & Exemple & Traduction \\
\hline
Énumération & 首先 / 其次 / 最后 & 首先，我谈原因。 & D'abord, je parle des causes. \\
\hline
Cause & 因为 / 由于 & 因为工资低，所以他走了。 & Parce que le salaire était bas, il est parti. \\
\hline
Conséquence & 所以 / 因此 & 法律不严，因此有人贪污。 & La loi est laxiste, donc certains sont corrompus. \\
\hline
But & 为了 + verbe & 为了发展，我们要努力。 & Pour nous développer, nous devons travailler dur. \\
\hline
Opposition & 但是 / 可是 & 问题很大，但是有办法。 & Le problème est grand, mais il y a des solutions. \\
\hline
\end{tabularx}
\end{center}

\begin{popfigure}
\begin{tikzpicture}[mindmap, grow cyclic,
every node/.style={concept, font=\small},
concept color=popblue,
level 1/.append style={level distance=3.2cm, sibling angle=90},
level 2/.append style={level distance=2.4cm, sibling angle=50, font=\footnotesize}]
\node[concept, text=white, align=center]{连接词\\ connecteurs}
child[concept color=popgreen] { node[concept, align=center]{énumération}
  child { node[concept]{首先} }
  child { node[concept]{其次} }
  child { node[concept]{最后} }
}
child[concept color=poporange] { node[concept, align=center]{cause / conséquence}
  child { node[concept]{因为} }
  child { node[concept]{所以} }
  child { node[concept]{因此} }
}
child[concept color=poppurple] { node[concept, align=center]{but}
  child { node[concept]{为了} }
}
child[concept color=poppink] { node[concept, align=center]{opposition}
  child { node[concept]{但是} }
  child { node[concept]{可是} }
};
\end{tikzpicture}
\legende{Carte mentale des connecteurs classés par fonction logique.}
\end{popfigure}

\courssub{Construire son plan avant de rédiger}

Un bon texte commence toujours par un bon plan. Voici la démarche à suivre en deux minutes, avant d'écrire la première phrase.

\begin{methode}[Le plan en 2 minutes : 问题 → 原因 → 办法]
\begin{enumerate}
\item Écrivez trois mots-clés sur votre brouillon : \zh{问题} (\emph{wèntí}, problème), \zh{原因} (\emph{yuányīn}, causes), \zh{办法} (\emph{bànfǎ}, solutions).
\item Sous \zh{问题}, notez une phrase qui présente le sujet : \zh{腐败是一个严重的问题。}
\item Sous \zh{原因}, notez deux ou trois causes, chacune avec \zh{因为…所以…} ; annoncez-les avec \zh{首先、其次、最后}.
\item Sous \zh{办法}, notez une ou deux solutions introduites par \zh{为了…，应该…}.
\item Terminez par une phrase de conclusion avec \zh{因此} ou \zh{只有这样…才…} : \zh{只有这样，我们的国家才能发展。}
\end{enumerate}
Vous obtenez le squelette complet du texte : problème → 2-3 causes → solution → conclusion. Il ne reste plus qu'à rédiger en reliant les phrases avec les connecteurs.
\end{methode}

\begin{exoresolu}[Relier les phrases avec le bon connecteur]
Énoncé : reliez chaque paire de phrases avec le connecteur qui convient (因为…所以…, 为了, 但是).\\
1. 法律不严格。有人敢贪污。\\
2. 解决这个问题。政府应该加强监督。\\
3. 腐败是一个严重的问题。我们有办法解决它。
\tcblower
Solution détaillée :\\
1. Relation cause → résultat : \zh{因为法律不严格，所以有人敢贪污。} (\emph{Yīnwèi fǎlǜ bù yángé, suǒyǐ yǒu rén gǎn tānwū.}) « Parce que la loi n'est pas stricte, certains osent être corrompus. » On garde les deux connecteurs dans la même phrase, séparés par la virgule \zh{，}.\\
2. Relation but → action : \zh{为了解决这个问题，政府应该加强监督。} (\emph{Wèile jiějué zhège wèntí, zhèngfǔ yīnggāi jiāqiáng jiāndū.}) « Pour résoudre ce problème, le gouvernement doit renforcer le contrôle. » \zh{为了} se place en tête, suivi du verbe.\\
3. Relation d'opposition : \zh{腐败是一个严重的问题，但是我们有办法解决它。} (\emph{Fǔbài shì yí ge yánzhòng de wèntí, dànshì wǒmen yǒu bànfǎ jiějué tā.}) « La corruption est un problème grave, mais nous avons des moyens de la résoudre. »
\end{exoresolu}

\begin{exercice}[À vous de jouer]
1. Traduisez en chinois : « D'abord, je parle des causes ; ensuite, je parle des solutions. »\\
2. Complétez avec 所以 ou 因此 : \zh{工资太低，……一些公务员想多拿钱。}\\
3. Construisez une phrase de but avec \zh{为了} sur le thème de la lutte contre la corruption au Cameroun.
\end{exercice}

\begin{corrige}[Exercice « À vous de jouer »]
1. \zh{首先，我谈原因；其次，我谈办法。} (\emph{Shǒuxiān, wǒ tán yuányīn ; qícì, wǒ tán bànfǎ.})\\
2. \zh{工资太低，所以一些公务员想多拿钱。} — \zh{所以} convient car il est rattaché à la cause dans la même phrase ; \zh{因此} serait aussi possible dans un style plus formel.\\
3. Exemple : \zh{为了反对腐败，喀麦隆政府应该制定严格的法律。} (\emph{Wèile fǎnduì fǔbài, Kāmàilóng zhèngfǔ yīnggāi zhìdìng yángé de fǎlǜ.}) « Pour lutter contre la corruption, le gouvernement camerounais doit établir des lois strictes. »
\end{corrige}

\begin{aretenir}[L'essentiel de la section]
\begin{itemize}
\item Énumération : \zh{首先…其次…最后…} (d'abord, ensuite, enfin).
\item Cause : \zh{因为 + cause，所以 + résultat} — les deux mots peuvent coexister dans la même phrase, contrairement au français.
\item \zh{所以} ne commence jamais un texte et évite de commencer une phrase après un point ; dans ce cas, utilisez \zh{因此}.
\item But : \zh{为了 + verbe，…} ; opposition : \zh{但是 / 可是}.
\item Ponctuation : virgule \zh{，}, point \zh{。}, pas de majuscules.
\item Plan express : \zh{问题 → 原因 → 办法}, puis conclusion.
\end{itemize}
\end{aretenir}

\coursec{Écriture des caractères complexes du thème}

Dans la section précédente, nous avons vu comment organiser un texte sur les facteurs de la corruption grâce à une structure claire et à des connecteurs comme \zh{首先} (d'abord), \zh{其次} (ensuite) et \zh{所以} (donc). Mais un texte bien structuré ne suffit pas : à l'écrit, l'examinateur juge aussi la \cle{qualité de l'écriture des caractères}. Un caractère mal tracé peut changer de sens ou devenir illisible. Cette section vous apprend donc à écrire correctement, trait par trait, les caractères complexes du thème : \zh{腐}、\zh{败}、\zh{贪}、\zh{污}、\zh{廉}、\zh{诚}、\zh{严}、\zh{格}.

\courssub{Observer d'abord : les caractères du thème dans leur contexte}

Avant d'apprendre à tracer, observons ces caractères dans des phrases réelles. Lisez les exemples suivants et repérez les caractères qui se ressemblent.

\begin{exemplebox}[Les caractères du thème en contexte]
\zh{贪污是严重的社会问题。} \emph{Tānwū shì yánzhòng de shèhuì wèntí.} « La corruption est un grave problème social. »\\
\zh{腐败的官员会受到惩罚。} \emph{Fǔbài de guānyuán huì shòudào chéngfá.} « Les fonctionnaires corrompus seront punis. »\\
\zh{他是一个诚实的人。} \emph{Tā shì yí ge chéngshí de rén.} « C'est un homme honnête. »\\
\zh{政府对贪污行为采取严格的措施。} \emph{Zhèngfǔ duì tānwū xíngwéi cǎiqǔ yángé de cuòshī.} « Le gouvernement prend des mesures strictes contre les actes de corruption. »
\end{exemplebox}

Observez : \zh{贪} et \zh{败} contiennent tous les deux le même élément, en bas pour l'un, à gauche pour l'autre : \zh{贝}. Ce n'est pas un hasard. C'est ce que nous allons comprendre maintenant.

\courssub{Le radical 贝 : la clé de l'argent}

\begin{propriete}[Le radical 贝 (bèi), le coquillage-monnaie]
Le caractère \zh{贝} (bèi) signifie « coquillage ». Dans la Chine ancienne, les coquillages servaient de \cle{monnaie}. C'est pourquoi le radical \zh{贝} apparaît dans presque tous les caractères liés à l'argent, à la richesse ou aux échanges :
\begin{itemize}
\item \zh{贪} tān : avidité, cupidité ;
\item \zh{败} bài : défaite, corrompre ;
\item \zh{贩} fàn : vendre, commercer ;
\item \zh{贫} pín : pauvreté ;
\item \zh{贵} guì : cher, précieux ;
\item \zh{财} cái : richesse, fortune.
\end{itemize}
\textbf{Règle pratique :} quand vous rencontrez un caractère inconnu contenant \zh{贝}, supposez qu'il a un rapport avec l'argent. Cet indice sémantique aide énormément à deviner le sens.
\end{propriete}

\begin{savaistu}[Les coquillages, première monnaie de Chine]
Il y a plus de trois mille ans, sous la dynastie Shang, les Chinois utilisaient de petits coquillages marins comme monnaie d'échange. Les coquillages étant rares à l'intérieur des terres, ils avaient une grande valeur. Plus tard, on a fabriqué des imitations en bronze, puis des pièces de monnaie. Mais la trace est restée dans l'écriture : le radical \zh{贝} marque encore aujourd'hui les caractères liés à la richesse… et, ironie de l'histoire, à la corruption (\zh{贪污}、\zh{腐败}) ! Au Cameroun, avant la monnaie moderne, on utilisait aussi des objets-valeurs comme le sel ou les cauris dans les échanges : les deux civilisations ont connu une « monnaie-objet ».
\end{savaistu}

\courssub{Rappel : l'ordre général des traits}

Tout caractère chinois s'écrit selon des règles d'ordre strictes. Les respecter rend l'écriture plus rapide, plus belle et plus facile à mémoriser.

\begin{aretenir}[Les cinq règles d'or de l'ordre des traits]
\begin{enumerate}
\item De \cle{haut en bas} : \zh{三} (sān, trois) se trace du trait supérieur vers le trait inférieur.
\item De \cle{gauche à droite} : \zh{你} (nǐ) se trace d'abord à gauche (\zh{亻}), puis à droite.
\item L'\cle{extérieur avant l'intérieur} : dans \zh{腐}, on trace d'abord l'enveloppe \zh{广}, puis l'intérieur \zh{府}.
\item La \cle{fermeture en dernier} : dans \zh{国} (guó, pays), le trait horizontal du bas qui « ferme la porte » se trace en dernier.
\item L'horizontal \cle{avant} le vertical quand ils se croisent : \zh{十} (shí, dix) : d'abord 一, puis 丨.
\end{enumerate}
\end{aretenir}

\courssub{Étude caractère par caractère}

\textbf{1. 腐 fǔ — pourri, corrompu (14 traits)}\par
Décomposition : enveloppe \zh{广} (abri, bâtiment) + \zh{府} (palais, administration) ; à l'intérieur de \zh{府}, on trouve encore \zh{付} et \zh{肉} (la viande, en bas). Image mentale : « de la viande (\zh{肉}) laissée dans un bâtiment (\zh{广}) finit par pourrir ». Ordre : on trace d'abord l'enveloppe \zh{广} (point, horizontal, puis longue descente à gauche), puis l'intérieur de haut en bas. C'est le caractère de \zh{腐败} (fǔbài, corruption) et de \zh{腐化} (fǔhuà, se corrompre).

\textbf{2. 败 bài — défaite, corrompre (8 traits)}\par
Décomposition : \zh{贝} (argent) à gauche + \zh{攵} (frapper, action) à droite. Image mentale : « frapper l'argent → le perdre → la défaite ». Ordre : gauche d'abord (\zh{贝} : vertical, horizontal-crochet, puis la descente à gauche et le point final), puis \zh{攵} de haut en bas (4 traits). On le trouve dans \zh{腐败} (fǔbài), \zh{失败} (shībài, échec) et \zh{打败} (dǎbài, vaincre).

\textbf{3. 贪 tān — avidité (8 traits)}\par
Décomposition : \zh{今} (jīn, aujourd'hui) au-dessus + \zh{贝} (argent) en dessous. Image mentale : « vouloir l'argent \emph{aujourd'hui}, tout de suite → la cupidité ». Ordre : de haut en bas, d'abord \zh{今} (4 traits), puis \zh{贝} (4 traits). Mot-clé : \zh{贪污} (tānwū, détournement de fonds, corruption).

\textbf{4. 污 wū — sale, souillé (6 traits)}\par
Décomposition : radical de l'eau \zh{氵} (trois gouttes, à gauche) + \zh{亏} (à droite). Image mentale : « de l'eau souillée ». Ordre : les trois gouttes de gauche d'abord (de haut en bas), puis la partie droite. Mots : \zh{贪污} (tānwū), \zh{污染} (wūrǎn, pollution), \zh{污水} (wūshuǐ, eaux usées).

\textbf{5. 廉 lián — intègre, honnête (13 traits)}\par
Décomposition : même enveloppe \zh{广} que \zh{腐} + \zh{兼} (jiān, simultané) à l'intérieur. Ordre : enveloppe d'abord, puis l'intérieur de haut en bas. Mots : \zh{廉洁} (liánjié, intègre, incorruptible), \zh{廉政} (liánzhèng, gouvernement intègre). Remarquez le contraste utile pour vos rédactions : \zh{腐} (corrompu) et \zh{廉} (intègre) partagent la même enveloppe \zh{广} — deux destins opposés pour un même « bâtiment » administratif.

\textbf{6. 诚 chéng — honnête, sincère (8 traits)}\par
Décomposition : radical de la parole \zh{讠} (à gauche) + \zh{成} (chéng, accomplir) à droite. Image mentale : « des paroles (\zh{讠}) qui s'accomplissent (\zh{成}) → la sincérité ». Le radical \zh{讠} indique toujours un lien avec la parole : \zh{说} (parler), \zh{话} (parole), \zh{语} (langue). Mots : \zh{诚实} (chéngshí, honnête), \zh{诚信} (chéngxìn, bonne foi, intégrité).

\textbf{7. 严 yán — strict (7 traits) et 格 gé — norme (10 traits)}\par
\zh{严} : caractère compact, à tracer de haut en bas ; attention au dernier trait, une longue descente à gauche. \zh{格} : radical du bois \zh{木} à gauche + \zh{各} (gè, chaque) à droite, de gauche à droite. Ensemble : \zh{严格} (yángé, strict, rigoureux), comme dans \zh{严格的法律} (yángé de fǎlǜ, des lois strictes).

\begin{attention}[Caractères proches : ne pas confondre]
\begin{itemize}
\item \zh{贪} (tān, avidité) ≠ \zh{贫} (pín, pauvreté) : même base \zh{贝}, mais au-dessus \zh{今} (aujourd'hui) ≠ \zh{分} (diviser). Erreur classique des francophones : écrire \zh{贫污} au lieu de \zh{贪污} !
\item \zh{败} (bài, défaite) ≠ \zh{贩} (fàn, vendre) : même gauche \zh{贝}, mais la droite diffère : \zh{攵} ≠ \zh{反}.
\item \zh{诚} (chéng, honnête) ≠ \zh{城} (chéng, ville) : même prononciation, mais \zh{讠} (parole) ≠ \zh{土} (terre). \zh{诚实} = honnête ; \zh{城市} = ville.
\item \zh{污} (wū, sale) : n'oubliez pas la troisième goutte du radical \zh{氵}.
\end{itemize}
\end{attention}

\begin{methode}[Mémoriser un caractère complexe en trois étapes]
\begin{enumerate}
\item \textbf{Décomposer} le caractère en 2 ou 3 composants connus (radical + élément phonétique). Exemple : \zh{腐} = \zh{广} + \zh{府}.
\item \textbf{Inventer une petite histoire} qui relie les composants au sens. Exemple : « la viande \zh{肉} oubliée dans le bâtiment \zh{广} pourrit ».
\item \textbf{Écrire 5 fois} le caractère dans une grille \zh{田字格}, en comptant les traits à voix haute et en respectant l'ordre. Vérifiez ensuite avec le texte de la leçon.
\end{enumerate}
\end{methode}

\begin{popfigure}
\begin{center}
\begin{tikzpicture}[font=\small]
\node[draw=popblue, fill=popblueL, rounded corners, minimum width=1.1cm, minimum height=1.1cm] (fu) at (0,0) {\Large 腐};
\node[draw=popgreen, fill=popgreenL, rounded corners, minimum width=0.9cm, minimum height=0.9cm] (guang) at (-1.6,-1.8) {\Large 广};
\node[draw=poporange, fill=poporangeL, rounded corners, minimum width=0.9cm, minimum height=0.9cm] (fu2) at (1.6,-1.8) {\Large 府};
\draw[-{Stealth}, thick, popblue] (fu) -- (guang);
\draw[-{Stealth}, thick, popblue] (fu) -- (fu2);
\node[below=0.05cm of guang, text=popdark] {\footnotesize enveloppe};
\node[below=0.05cm of fu2, text=popdark] {\footnotesize intérieur};

\node[draw=popblue, fill=popblueL, rounded corners, minimum width=1.1cm, minimum height=1.1cm] (tan) at (4.5,0) {\Large 贪};
\node[draw=popgreen, fill=popgreenL, rounded corners, minimum width=0.9cm, minimum height=0.9cm] (jin) at (3.4,-1.8) {\Large 今};
\node[draw=poporange, fill=poporangeL, rounded corners, minimum width=0.9cm, minimum height=0.9cm] (bei1) at (5.6,-1.8) {\Large 贝};
\draw[-{Stealth}, thick, popblue] (tan) -- (jin);
\draw[-{Stealth}, thick, popblue] (tan) -- (bei1);
\node[below=0.05cm of jin, text=popdark] {\footnotesize haut};
\node[below=0.05cm of bei1, text=popdark] {\footnotesize bas};

\node[draw=popblue, fill=popblueL, rounded corners, minimum width=1.1cm, minimum height=1.1cm] (bai) at (9,0) {\Large 败};
\node[draw=popgreen, fill=popgreenL, rounded corners, minimum width=0.9cm, minimum height=0.9cm] (bei2) at (7.9,-1.8) {\Large 贝};
\node[draw=poporange, fill=poporangeL, rounded corners, minimum width=0.9cm, minimum height=0.9cm] (pu) at (10.1,-1.8) {\Large 攵};
\draw[-{Stealth}, thick, popblue] (bai) -- (bei2);
\draw[-{Stealth}, thick, popblue] (bai) -- (pu);
\node[below=0.05cm of bei2, text=popdark] {\footnotesize gauche};
\node[below=0.05cm of pu, text=popdark] {\footnotesize droite};
\end{tikzpicture}
\end{center}
\legende{Décomposition des caractères : 腐 = 广 + 府 (extérieur → intérieur) ; 贪 = 今 + 贝 (haut → bas) ; 败 = 贝 + 攵 (gauche → droite).}
\end{popfigure}

\begin{popfigure}
\begin{center}
\begin{tikzpicture}[scale=1.6]
\draw[thick, popdark] (0,0) rectangle (2,2);
\draw[dashed, popmuted] (1,0) -- (1,2);
\draw[dashed, popmuted] (0,1) -- (2,1);
\draw[dashed, popmuted] (0,0) -- (2,2);
\draw[dashed, popmuted] (0,2) -- (2,0);
\draw[very thick, popink] (0.45,1.75) -- (0.45,0.55);
\node[popink, font=\footnotesize] at (0.28,1.7) {1};
\draw[very thick, popink] (0.45,1.75) -- (0.95,1.75) -- (0.95,0.55);
\node[popink, font=\footnotesize] at (0.75,1.9) {2};
\draw[very thick, popink] (0.55,1.05) -- (0.85,1.05);
\node[popink, font=\footnotesize] at (0.7,1.22) {3};
\draw[very thick, popink] (0.62,0.5) -- (0.35,0.15);
\node[popink, font=\footnotesize] at (0.32,0.42) {4};
\draw[very thick, poporange] (1.25,1.8) -- (1.6,1.35);
\node[poporange, font=\footnotesize] at (1.28,1.92) {5};
\draw[very thick, poporange] (1.6,1.75) -- (1.25,1.3);
\node[poporange, font=\footnotesize] at (1.72,1.7) {6};
\draw[very thick, poporange] (1.35,1.25) -- (1.75,0.6);
\node[poporange, font=\footnotesize] at (1.5,1.15) {7};
\draw[very thick, poporange] (1.45,1.0) -- (1.15,0.2);
\node[poporange, font=\footnotesize] at (1.12,0.65) {8};
\end{tikzpicture}
\end{center}
\legende{Grille d'écriture 田字格 : les 8 traits de 败. Traits 1 à 4 : 贝 à gauche (de haut en bas) ; traits 5 à 8 : 攵 à droite.}
\end{popfigure}

\courssub{Tableau récapitulatif des caractères du thème}

\begin{center}
\begin{tabularx}{\linewidth}{|X|X|X|X|}\hline
\rowcolor{popblueL} Caractères & Pinyin & Nature & Traduction \\ \hline
\zh{腐} & fǔ & verbe / adj. & pourrir, corrompu \\ \hline
\zh{败} & bài & verbe & perdre, corrompre \\ \hline
\zh{腐败} & fǔbài & adj. / nom & corrompu ; corruption \\ \hline
\zh{贪} & tān & verbe / adj. & convoiter, avide \\ \hline
\zh{污} & wū & adj. & sale, souillé \\ \hline
\zh{贪污} & tānwū & nom / verbe & détournement, corruption \\ \hline
\zh{廉} & lián & adj. & intègre \\ \hline
\zh{廉洁} & liánjié & adj. & intègre, incorruptible \\ \hline
\zh{诚} & chéng & adj. & sincère \\ \hline
\zh{诚实} & chéngshí & adj. & honnête \\ \hline
\zh{严格} & yángé & adj. & strict, rigoureux \\ \hline
\zh{惩罚} & chéngfá & nom / verbe & punition, punir \\ \hline
\end{tabularx}
\end{center}

\begin{exemplebox}[Les mots du thème en phrases]
\zh{贪污浪费是极大的犯罪。} \emph{Tānwū làngfèi shì jí dà de zuìxíng.} « Le détournement et le gaspillage sont des crimes très graves. »\\
\zh{做人要诚实。} \emph{Zuòrén yào chéngshí.} « Dans la vie, il faut être honnête. »\\
\zh{老师对我们要求很严格。} \emph{Lǎoshī duì wǒmen yāoqiú hěn yángé.} « Le professeur est très strict envers nous. »\\
\zh{我们要建设廉洁的社会。} \emph{Wǒmen yào jiànshè liánjié de shèhuì.} « Nous devons construire une société intègre. »
\end{exemplebox}

\begin{exoresolu}[Compléter les caractères et tracer]
\textbf{Énoncé.} Complétez les mots suivants avec le caractère qui convient (\zh{贪}、\zh{贫}、\zh{败}、\zh{贩}、\zh{诚}、\zh{城}), puis donnez le pinyin et la traduction.\\
1. \zh{腐\_\_} (corruption) \quad 2. \_\_\zh{污} (détournement) \quad 3. \_\_\zh{穷} (pauvre) \quad 4. \_\_\zh{实} (honnête) \quad 5. \_\_\zh{市} (ville)
\tcblower
\textbf{Solution détaillée.}\\
1. \zh{腐败} (fǔbài) : « corruption ». On choisit \zh{败} (défaite) et non \zh{贩} (vendre) : la corruption est une « défaite morale ».\\
2. \zh{贪污} (tānwū) : « détournement ». On choisit \zh{贪} (\zh{今} + \zh{贝} : vouloir l'argent tout de suite) et non \zh{贫} (\zh{分} + \zh{贝} : partager l'argent → pauvreté).\\
3. \zh{贫穷} (pínqióng) : « pauvre ». Le caractère demandé est \zh{贫}.\\
4. \zh{诚实} (chéngshí) : « honnête ». On choisit \zh{诚} (radical de la parole \zh{讠}) et non \zh{城} (radical de la terre \zh{土}).\\
5. \zh{城市} (chéngshì) : « ville ». Le caractère demandé est \zh{城}.\\
\textbf{Astuce de vérification :} relisez chaque mot en identifiant le radical : \zh{讠} → parole, \zh{土} → terre, \zh{贝} → argent. Le radical confirme le sens.
\end{exoresolu}

\begin{exercice}[À vous de tracer]
\begin{enumerate}
\item Écrivez 5 fois chacun des caractères \zh{腐}、\zh{贪}、\zh{诚} dans une grille \zh{田字格}, en comptant les traits à voix haute (14, 8 et 8 traits).
\item Pour chaque caractère du tableau récapitulatif, écrivez sa décomposition (exemple : \zh{廉} = \zh{广} + \zh{兼}).
\item Recopiez la phrase \zh{腐败的官员会受到严格的惩罚。} en soignant l'ordre des traits, puis traduisez-la en français.
\end{enumerate}
\end{exercice}

\begin{corrige}[Exercice « À vous de tracer »]
1. Vérifiez le nombre de traits : \zh{腐} = 14 traits (enveloppe \zh{广} d'abord, puis l'intérieur de haut en bas) ; \zh{贪} = 8 traits (\zh{今} puis \zh{贝}) ; \zh{诚} = 8 traits (\zh{讠} puis \zh{成}).\\
2. \zh{腐} = \zh{广} + \zh{府} ; \zh{贪} = \zh{今} + \zh{贝} ; \zh{诚} = \zh{讠} + \zh{成} ; \zh{败} = \zh{贝} + \zh{攵} ; \zh{污} = \zh{氵} + \zh{亏} ; \zh{廉} = \zh{广} + \zh{兼} ; \zh{格} = \zh{木} + \zh{各}.\\
3. Traduction : « Les fonctionnaires corrompus recevront une punition sévère. » Attention à l'ordre des mots : en chinois, le complément \zh{严格的} se place avant le nom \zh{惩罚}, relié par la particule \zh{的} ; le verbe \zh{受到} (recevoir, subir) est suivi directement de son complément d'objet.
\end{corrige}

\begin{attention}[Erreurs fréquentes des francophones à l'écrit]
\begin{itemize}
\item \textbf{Oublier l'enveloppe d'abord} : dans \zh{腐} et \zh{廉}, beaucoup d'élèves tracent l'intérieur avant l'enveloppe \zh{广}. Résultat : le caractère est déséquilibré. Toujours l'extérieur avant l'intérieur.
\item \textbf{Écrire le pinyin à la place du caractère} dans une rédaction : à l'examen, un mot en pinyin dans un texte en caractères est pénalisé. Apprenez les 12 entrées du tableau par cœur.
\item \textbf{Confondre les tons} : \zh{诚实} se dit chéngshí (2\textsuperscript{e} ton), pas chěngshí ; \zh{严格} se dit yángé, avec deux 2\textsuperscript{e} tons.
\item \textbf{Mal fermer} \zh{贝} : le dernier trait (le point en bas à droite) se trace en dernier ; sans lui, \zh{贝} ressemble à \zh{见} (jiàn, voir).
\end{itemize}
\end{attention}

Vous savez maintenant écrire et distinguer les caractères complexes du thème de la corruption. La section suivante approfondira ce travail en comparant systématiquement les caractères proches, avant de passer à la traduction (thème et version).

\coursec{Caractères proches : ne pas confondre}

Dans la section précédente, nous avons appris à écrire correctement les caractères complexes du thème de la corruption (\zh{腐败}, \zh{贪污}, \zh{法律}, \zh{诚实}…) en décomposant leurs traits et leurs radicaux. Mais connaître l'ordre des traits ne suffit pas : certains caractères se ressemblent tellement qu'un seul trait de différence change complètement le sens. C'est le piège classique de l'examen : écrire \zh{城} (ville) au lieu de \zh{诚} (honnête) dans une rédaction sur l'honnêteté ! Cette section vous donne une méthode infaillible pour ne plus jamais confondre ces jumeaux dangereux.

\courssub{Qu'est-ce qu'un caractère de forme proche ?}

\begin{definition}[Les 形近字 xíngjìnzì]
Les \cle{形近字} (\emph{xíngjìnzì}) sont les \textbf{caractères de forme proche} : des caractères qui se ressemblent visuellement (même radical, même structure, ou un seul composant différent) mais qui ont un sens, et souvent une prononciation, différents. Ils sont la \textbf{première source de fautes d'écriture} chez les apprenants francophones, et les correcteurs d'examen les sanctionnent systématiquement : un caractère mal choisi est une faute de vocabulaire, pas une simple « faute de frappe ».
\end{definition}

En français, vous distinguez facilement « corruption » et « coeuruption » parce que la seconde forme n'existe pas. En chinois, le problème est plus subtil : \zh{败} et \zh{贩} existent tous les deux, se prononcent différemment, et un seul coup d'œil rapide peut vous faire choisir le mauvais. La solution n'est pas de « regarder plus attentivement », mais d'appliquer une \textbf{méthode d'analyse en deux étapes}.

\begin{methode}[Distinguer deux caractères proches en deux questions]
Face à deux caractères qui se ressemblent, posez-vous toujours les deux questions suivantes, dans cet ordre :
\begin{enumerate}
\item \textbf{Quel est le radical (la clé de sens) ?} Le radical indique la famille de sens : \zh{讠} (parole) pour ce qui se dit, \zh{土} (terre) pour ce qui se construit, \zh{氵} (eau) pour ce qui est liquide, \zh{彳} (pas, rue) pour ce qui est en mouvement ou réglé. Le sens de votre phrase désigne le bon radical.
\item \textbf{Quel est le composant phonétique (l'indice de son) ?} L'autre partie du caractère donne souvent la prononciation : \zh{反} fǎn dans \zh{贩} fàn, \zh{成} chéng dans \zh{诚} et \zh{城} chéng.
\end{enumerate}
Résumé : \cle{le radical donne le sens, le composant donne le son}. Vérifiez les deux, et l'erreur devient impossible.
\end{methode}

\courssub{Paire 1 : 败 bài (corrompre, perdre) / 贩 fàn (vendre)}

Observez d'abord ces deux phrases tirées de notre thème :

\begin{exemplebox}[Observation]
\zh{腐败是一个严重的问题。} (\emph{Fǔbài shì yī ge yánzhòng de wèntí.}) « La corruption est un problème grave. »\\
\zh{这个小贩卖水果。} (\emph{Zhège xiǎofàn mài shuǐguǒ.}) « Ce petit vendeur vend des fruits. »
\end{exemplebox}

Les deux caractères partagent le radical \zh{贝} (coquillage, autrefois monnaie → argent). Toute la différence est à droite : \zh{攵} (frapper, action) dans \zh{败} bài — l'argent « frappé », abîmé → perdre, corrompre ; \zh{反} fǎn/fàn dans \zh{贩} fàn — le composant phonétique donne le son → vendre, revendre.

\begin{attention}[Piège]
Dans \zh{腐败} fǔbài (corruption), le second caractère s'écrit avec \zh{攵} à droite. Si vous écrivez \zh{贩}, vous écrivez « vendre » : votre phrase sur la corruption devient une phrase sur le commerce ! Retenez : \textbf{败 = perdre/échouer} (\zh{失败} shībài, échec), \textbf{贩 = vendre} (\zh{小贩} xiǎofàn, petit vendeur).
\end{attention}

\courssub{Paire 2 : 贪 tān (avide) / 贫 pín (pauvre)}

\begin{exemplebox}[Observation]
\zh{贪官把钱拿走了。} (\emph{Tānguān bǎ qián ná zǒu le.}) « Le fonctionnaire corrompu a pris l'argent. »\\
\zh{我们要帮助贫困的人。} (\emph{Wǒmen yào bāngzhù pínkùn de rén.}) « Nous devons aider les personnes dans la pauvreté. »
\end{exemplebox}

Même radical \zh{贝} (argent) en bas ; la différence est en haut : \zh{今} jīn (aujourd'hui) dans \zh{贪} tān — vouloir l'argent \emph{aujourd'hui}, tout de suite → avidité ; \zh{分} fēn (diviser) dans \zh{贫} pín — l'argent \emph{divisé} jusqu'à ne plus rien avoir → pauvreté. Ces deux petites histoires de sens sont un excellent moyen mnémotechnique : \textbf{今 + 贝 = 贪} (l'avide veut tout maintenant), \textbf{分 + 贝 = 贫} (le pauvre a tout partagé).

\courssub{Paire 3 : 律 lǜ (loi) / 津 jīn (Tianjin)}

\zh{律} porte le radical \zh{彳} (pas, rue, règle de conduite) : on le trouve dans \zh{法律} fǎlǜ (la loi) et \zh{纪律} jìlǜ (la discipline). \zh{津} porte le radical \zh{氵} (eau) : c'est le caractère de la ville de \zh{天津} Tiānjīn, et son sens ancien est « gué, passage d'eau ». Même composant phonétique \zh{聿} à droite, mais radicaux opposés : \textbf{la règle marche (彳), le gué est dans l'eau (氵)}. Dans une rédaction sur la lutte contre la corruption, c'est toujours \zh{法律} avec \zh{彳} qu'il faut écrire.

\courssub{Paire 4 : 诚 chéng (honnête) / 城 chéng (ville)}

Voici la paire la plus dangereuse, car les deux caractères se prononcent \emph{exactement pareil} : chéng, même ton. À l'oral, impossible de les distinguer ; à l'écrit, seul le radical vous sauve.

\begin{exemplebox}[Exemples traduits]
\zh{他是一个诚实的人。} (\emph{Tā shì yī ge chéngshí de rén.}) « C'est un homme honnête. »\\
\zh{这个城市很大。} (\emph{Zhège chéngshì hěn dà.}) « Cette ville est grande. »
\end{exemplebox}

\begin{attention}[诚 ou 城 ? Le radical décide]
\zh{诚} et \zh{城} se prononcent tous les deux \emph{chéng} — seul le radical permet de choisir le bon caractère à l'écrit : \zh{讠} (parole) pour \zh{诚} → ce qu'on dit est vrai → \textbf{honnête, sincère} (\zh{诚实} chéngshí) ; \zh{土} (terre) pour \zh{城} → les murs de terre → \textbf{ville, muraille} (\zh{城市} chéngshì, \zh{长城} Chángchéng, la Grande Muraille). Question à vous poser : « Est-ce que je parle de parole ou de terre ? »
\end{attention}

\courssub{Paire 5 : 严 yán (strict) / 产 chǎn (produire)}

Ces deux caractères n'ont ni même radical ni même son, mais leurs silhouettes se ressemblent (une ligne horizontale haute, des traits verticaux qui descendent). \zh{严} yán signifie strict, sévère : \zh{严格} yángé (strict), \zh{严重} yánzhòng (grave). \zh{产} chǎn signifie produire : \zh{生产} shēngchǎn (produire), \zh{产品} chǎnpǐn (produit). Astuce : \zh{严} a deux petits traits symétriques en haut, comme deux sourcils froncés d'un chef \emph{strict} ; \zh{产} n'en a qu'un.

\courssub{Tableau récapitulatif des cinq paires}

\begin{center}
\begin{tabularx}{\linewidth}{|X|X|X|X|}
\hline
\rowcolor{popblueL} Caractère & Pinyin / Sens & Différence & Mot du thème \\
\hline
败 & bài, corrompre, perdre & 贝 + 攵 (frapper) & 腐败 fǔbài, corruption \\
\hline
贩 & fàn, vendre & 贝 + 反 (phonétique) & 小贩 xiǎofàn, vendeur \\
\hline
贪 & tān, avide & 今 + 贝 & 贪污 tānwū, détourner \\
\hline
贫 & pín, pauvre & 分 + 贝 & 贫穷 pínqióng, pauvreté \\
\hline
律 & lǜ, loi, règle & 彳 (pas, règle) & 法律 fǎlǜ, loi \\
\hline
津 & jīn, gué ; Tianjin & 氵 (eau) & 天津 Tiānjīn \\
\hline
诚 & chéng, honnête & 讠(parole) + 成 & 诚实 chéngshí, honnête \\
\hline
城 & chéng, ville & 土 (terre) + 成 & 城市 chéngshì, ville \\
\hline
严 & yán, strict, grave & deux traits en haut & 严重 yánzhòng, grave \\
\hline
产 & chǎn, produire & un trait en haut & 生产 shēngchǎn, produire \\
\hline
\end{tabularx}
\end{center}

\begin{popfigure}
\begin{tikzpicture}[
  cell/.style={draw=popline, rounded corners=2pt, minimum width=2.6cm, minimum height=1.5cm, align=center, font=\large},
  head/.style={font=\bfseries, text=popdark}
]
\node[head] at (-1.6,3.3) {Caractère A};
\node[head] at (1.6,3.3) {Caractère B};
\node[cell, fill=popblueL, align=center] at (-1.6,2.3){败 bài\\ \textcolor{poporange}{攵} à droite};
\node[cell, fill=poporangeL, align=center] at (1.6,2.3){贩 fàn\\ \textcolor{popblue}{反} à droite};
\node[cell, fill=popblueL, align=center] at (-1.6,0.9){贪 tān\\ \textcolor{poporange}{今} en haut};
\node[cell, fill=poporangeL, align=center] at (1.6,0.9){贫 pín\\ \textcolor{popblue}{分} en haut};
\node[cell, fill=popblueL, align=center] at (-1.6,-0.5){诚 chéng\\ \textcolor{poporange}{讠} parole};
\node[cell, fill=poporangeL, align=center] at (1.6,-0.5){城 chéng\\ \textcolor{popblue}{土} terre};
\node[cell, fill=popblueL, align=center] at (-1.6,-1.9){律 lǜ\\ \textcolor{poporange}{彳} règle};
\node[cell, fill=poporangeL, align=center] at (1.6,-1.9){津 jīn\\ \textcolor{popblue}{氵} eau};
\node[cell, fill=popblueL, align=center] at (-1.6,-3.3){严 yán\\ \textcolor{poporange}{丷} deux traits};
\node[cell, fill=poporangeL, align=center] at (1.6,-3.3){产 chǎn\\ \textcolor{popblue}{一} un trait};
\end{tikzpicture}
\legende{Cinq paires de 形近字 : le composant différent est surligné en couleur. Comparez toujours les deux caractères côte à côte avant d'écrire.}
\end{popfigure}

\courssub{Stratégie d'entraînement : la phrase personnelle}

La meilleure façon de fixer une paire de 形近字 est d'écrire \textbf{une phrase avec chaque caractère}, si possible une phrase qui vous parle (votre ville, votre école, le Cameroun). Exemples :

\begin{exemplebox}[Une phrase par caractère]
\zh{我们要反对腐败。} (\emph{Wǒmen yào fǎnduì fǔbài.}) « Nous devons lutter contre la corruption. »\\
\zh{市场上的小贩很多。} (\emph{Shìchǎng shàng de xiǎofàn hěn duō.}) « Il y a beaucoup de petits vendeurs au marché. »\\
\zh{诚实的学生不说谎。} (\emph{Chéngshí de xuésheng bù shuōhuǎng.}) « Un élève honnête ne ment pas. »\\
\zh{雅温得是一个大城市。} (\emph{Yǎwēndé shì yī ge dà chéngshì.}) « Yaoundé est une grande ville. »\\
\zh{考试纪律非常严格。} (\emph{Kǎoshì jìlǜ fēicháng yángé.}) « La discipline d'examen est très stricte. »
\end{exemplebox}

\begin{exoresolu}[Choisir le bon caractère]
Énoncé — Complétez chaque phrase avec le bon caractère de la paire entre parenthèses, puis traduisez :
\begin{enumerate}
\item \zh{他是一个（诚 / 城）实的人。}
\item \zh{这个（诚 / 城）市非常漂亮。}
\item \zh{（贪 / 贫）官受到法律的惩罚。}
\item \zh{法（律 / 津）保护公民。}
\end{enumerate}
\tcblower
Solution détaillée :
\begin{enumerate}
\item \zh{诚} : le sens est « honnête » (\zh{诚实} chéngshí), il faut le radical de la parole \zh{讠}. Traduction : « C'est un homme honnête. »
\item \zh{城} : le sens est « ville » (\zh{城市} chéngshì), il faut le radical de la terre \zh{土}. Traduction : « Cette ville est très belle. »
\item \zh{贪} : \zh{贪官} tānguān, « fonctionnaire avide/corrompu » ; \zh{今} au-dessus de \zh{贝}. Traduction : « Les fonctionnaires corrompus sont punis par la loi. »
\item \zh{律} : \zh{法律} fǎlǜ, « la loi » ; radical \zh{彳}, pas \zh{氵}. Traduction : « La loi protège les citoyens. »
\end{enumerate}
Méthode appliquée : pour chaque item, on a d'abord identifié le \emph{sens} voulu (question 1 : quel radical ?), puis vérifié la forme complète du caractère avant d'écrire.
\end{exoresolu}

\begin{exercice}[À vous de jouer]
Écrivez une phrase en chinois avec chacun des caractères suivants, en soulignant le radical qui vous a guidé : \zh{败}, \zh{贩}, \zh{贫}, \zh{严}, \zh{产}. Puis traduisez vos cinq phrases en français.
\end{exercice}

\begin{aretenir}[L'essentiel de la section]
\begin{itemize}
\item Les \cle{形近字} (caractères de forme proche) sont la première source de fautes à l'examen : un trait différent = un mot différent.
\item Méthode en deux questions : \textbf{quel radical} (le sens) ? \textbf{quel composant phonétique} (le son) ?
\item Paires du thème à connaître par cœur : \zh{败}/\zh{贩}, \zh{贪}/\zh{贫}, \zh{律}/\zh{津}, \zh{诚}/\zh{城} (homophones !), \zh{严}/\zh{产}.
\item Entraînement efficace : comparer les paires côte à côte et écrire une phrase personnelle avec chaque caractère.
\end{itemize}
\end{aretenir}

\coursec{Thème : français → chinois}

Après avoir appris à distinguer les caractères proches du thème (\zh{贪} et \zh{贫}, \zh{法} et \zh{去}, \zh{肃} et \zh{萧}…), vous savez désormais écrire correctement les mots de la corruption. Il reste une étape décisive : les \textbf{assembler en phrases}. C'est exactement ce que demande l'épreuve de \cle{thème} : traduire du français vers le chinois. Cette section vous donne une méthode en quatre étapes, puis cinq phrases modèles commentées une à une, comme à l'examen.

\courssub{Rappel de méthode : les quatre étapes du thème}

\begin{methode}[Réussir le thème en quatre étapes]
Le thème ne se fait jamais mot à mot. Suivez toujours cet ordre :
\begin{enumerate}
\item \textbf{Lire} la phrase française entière, sans écrire. Repérez le sens global : qui fait quoi, quand, pourquoi ?
\item \textbf{Analyser} : soulignez le \cle{sujet}, le \cle{verbe} et le \cle{complément}. Repérez les connecteurs (parce que, donc, pour, mais…) et les expressions figées (pas assez, de plus en plus…).
\item \textbf{Traduire la structure d'abord} : posez le squelette chinois \zh{Sujet + Verbe + Complément}, puis ajoutez les connecteurs, les adverbes et les compléments de temps ou de lieu \emph{avant} le verbe.
\item \textbf{Vérifier} : ordre des mots, tons, radicaux des caractères, ponctuation chinoise pleine chasse (\zh{，} \zh{。}) .
\end{enumerate}
\end{methode}

\begin{popfigure}
\begin{tikzpicture}[
  etape/.style={rectangle, rounded corners=6pt, draw=popblue, fill=popblueL, thick, minimum width=3.1cm, minimum height=1.1cm, align=center, font=\small},
  fleche/.style={-{Stealth[length=3mm]}, very thick, poporange}
]
\node[etape, align=center] (e1) at (0,0){\textbf{1. Lire}\\ sens global};
\node[etape, right=1.1cm of e1, align=center] (e2){\textbf{2. Analyser}\\ S + V + C};
\node[etape, right=1.1cm of e2, align=center] (e3){\textbf{3. Traduire}\\ structure d'abord};
\node[etape, right=1.1cm of e3, align=center] (e4){\textbf{4. Vérifier}\\ tons, traits, 。};
\draw[fleche] (e1) -- (e2);
\draw[fleche] (e2) -- (e3);
\draw[fleche] (e3) -- (e4);
\draw[fleche, dashed, popgreen] (e4.north) to[bend left=35] node[above, font=\small, text=popgreen]{si doute, relire} (e1.north);
\end{tikzpicture}
\legende{La méthode de traduction en quatre étapes : lire, analyser, traduire, vérifier.}
\end{popfigure}

\begin{attention}[Le piège n° 1 : le mot à mot]
Ne traduisez \textbf{jamais} mot à mot. Le français dit « La corruption est un problème grave » ; le chinois dira littéralement « corruption est un [classificateur] grave de problème » : \zh{腐败是一个严重的问题。} Remarquez le classificateur \zh{个} et la particule \zh{的} entre l'adjectif et le nom. Traduisez la \emph{structure}, pas les mots isolés.
\end{attention}

\courssub{Phrase 1 guidée : phrase simple avec 是}

\textbf{Phrase française :} « La corruption est un problème grave. »

\textbf{Analyse (étape 2) :}
\begin{itemize}
\item Sujet : \emph{la corruption} → \zh{腐败} (fǔbài)
\item Verbe : \emph{est} → \zh{是} (shì)
\item Complément : \emph{un problème grave} → \zh{一个严重的问题} (yí gè yánzhòng de wèntí)
\end{itemize}

\textbf{Traduction (étape 3) :} on assemble dans le même ordre qu'en français — c'est une phrase déclarative avec \zh{是}, l'ordre est identique :

\begin{exemplebox}[Solution de la phrase 1]
\zh{腐败是一个严重的问题。}\\
\emph{Fǔbài shì yí gè yánzhòng de wèntí.}\\
« La corruption est un problème grave. »
\end{exemplebox}

\textbf{Points de vigilance :} n'oubliez ni le classificateur \zh{个}, ni \zh{的} après l'adjectif de deux syllabes \zh{严重}. Un adjectif dissyllabique devant un nom exige presque toujours \zh{的}.

\courssub{Phrase 2 : la cause avec 因为…所以…}

\textbf{Phrase française :} « Parce que certaines personnes veulent de l'argent rapidement, elles détournent des fonds. »

\textbf{Analyse :} deux propositions reliées par une relation de \cle{cause}. En français, « parce que » suffit seul ; en chinois, on utilise très souvent la paire \zh{因为…所以…} (\emph{parce que… donc…}), même si le français ne dit pas « donc ».

\begin{propriete}[La corrélation 因为…所以…]
Structure : \zh{因为} + cause, \zh{所以} + conséquence.\\
Les deux membres peuvent être utilisés ensemble (usage le plus fréquent et le plus sûr à l'examen) ou séparément, mais jamais \zh{因为} seul suivi d'une virgule sans suite logique claire. En français, on ne traduit qu'un seul des deux (« parce que » ou « donc »).
\end{propriete}

\textbf{Traduction :}

\begin{exemplebox}[Solution de la phrase 2]
\zh{因为一些人想很快地有钱，所以他们贪污。}\\
\emph{Yīnwèi yìxiē rén xiǎng hěn kuài de yǒu qián, suǒyǐ tāmen tānwū.}\\
« Parce que certaines personnes veulent de l'argent rapidement, elles détournent des fonds. »
\end{exemplebox}

\textbf{Commentaire :} \zh{一些人} (yìxiē rén) = « certaines personnes » ; \zh{想} (xiǎng) = « vouloir » ; \zh{很快地有钱} = « avoir de l'argent rapidement » (l'adverbe de manière se place avant le verbe, avec \zh{地}) ; \zh{贪污} (tānwū) = « détourner des fonds, être corrompu ».

\courssub{Phrase 3 : « pas assez » avec 不够}

\textbf{Phrase française :} « La loi n'est pas assez stricte. »

\textbf{Analyse :} le mot-clé est \emph{pas assez}. En chinois, « assez + adjectif » se dit \zh{够} + adjectif ; la négation « pas assez » est donc \zh{不够} + adjectif.

\begin{propriete}[不够 + adjectif]
Structure : Sujet + \zh{不够} + adjectif.\\
\zh{够} (gòu) signifie « suffire, être assez ». Placé devant un adjectif, il forme « assez… » ; avec la négation \zh{不}, « pas assez… ». L'adjectif chinois fait office de verbe : on n'ajoute \textbf{pas} \zh{是} devant.
\end{propriete}

\begin{exemplebox}[Solution de la phrase 3]
\zh{法律不够严格。}\\
\emph{Fǎlǜ bù gòu yángé.}\\
« La loi n'est pas assez stricte. »
\end{exemplebox}

\begin{attention}[Ne dites jamais 不很严格]
Les francophones traduisent souvent « pas assez » par \zh{不很} sous l'influence du français « pas très ». C'est une erreur : \zh{不很严格} signifie « pas très strict », ce qui est plus faible et différent. « Pas assez » = \zh{不够} : \zh{不够严格}. De même : « pas assez clair » → \zh{不够清楚}, « pas assez rapide » → \zh{不够快}.
\end{attention}

\courssub{Phrase 4 : le but avec 为了 et le devoir avec 应该}

\textbf{Phrase française :} « Pour résoudre ce problème, nous devons renforcer l'éducation. »

\textbf{Analyse :} deux idées : le \cle{but} (« pour résoudre… ») et le \cle{devoir} (« nous devons… »).

\begin{propriete}[应该 + verbe : le devoir, la recommandation]
\zh{应该} (yīnggāi) exprime le devoir ou la recommandation morale, comme « devoir » ou « il faut » en français. Il se place \textbf{avant le verbe principal} : Sujet + \zh{应该} + verbe + complément. Négation : \zh{不应该} (« ne devoir pas »). Ne pas confondre avec \zh{必须} (bìxū, obligation forte) ni \zh{得} (děi, obligation pratique).
\end{propriete}

\begin{propriete}[为了 + verbe : le but]
\zh{为了} (wèile) introduit un but : « pour, afin de ». Structure : \zh{为了} + verbe + complément, + proposition principale. Il se place en tête de phrase, suivi d'une virgule chinoise \zh{，}.
\end{propriete}

\begin{exemplebox}[Solution de la phrase 4]
\zh{为了解决这个问题，我们应该加强教育。}\\
\emph{Wèile jiějué zhège wèntí, wǒmen yīnggāi jiāqiáng jiàoyù.}\\
« Pour résoudre ce problème, nous devons renforcer l'éducation. »\\
(\zh{应该} = devoir, conseil moral ; \zh{加强} = renforcer ; \zh{解决} = résoudre.)
\end{exemplebox}

\courssub{Phrase 5 (complément examen) : 做 + nom de rôle}

\textbf{Phrase française :} « Nous devons être des citoyens honnêtes. »

\textbf{Analyse :} « être » + une qualité de personne se rend souvent par \zh{做} (zuò, « faire, agir comme ») devant un nom de rôle ou de statut : \zh{做诚实的公民} = « être / se comporter en citoyen honnête ».

\begin{exemplebox}[Solution de la phrase 5]
\zh{我们应该做诚实的公民。}\\
\emph{Wǒmen yīnggāi zuò chéngshí de gōngmín.}\\
« Nous devons être des citoyens honnêtes. »
\end{exemplebox}

\textbf{Remarque :} \zh{是诚实的公民} serait grammaticalement possible, mais \zh{做} exprime mieux l'effort de conduite (« se comporter en… »), ce qui correspond au sens moral de la phrase.

\courssub{Tableau récapitulatif des structures du thème}

\begin{center}
\begin{tabularx}{\linewidth}{|X|X|X|}
\hline
\rowcolor{popblueL} Structure & Exemple en caractères + pinyin & Traduction \\
\hline
S + 是 + 数量 + adj. + 的 + N & \zh{腐败是一个严重的问题。} \emph{Fǔbài shì yí gè yánzhòng de wèntí.} & La corruption est un problème grave. \\
\hline
因为 + cause，所以 + résultat & \zh{因为一些人想很快地有钱，所以他们贪污。} \emph{Yīnwèi yìxiē rén xiǎng hěn kuài de yǒu qián, suǒyǐ tāmen tānwū.} & Parce que certaines personnes veulent de l'argent vite, elles détournent des fonds. \\
\hline
S + 不够 + adjectif & \zh{法律不够严格。} \emph{Fǎlǜ bù gòu yángé.} & La loi n'est pas assez stricte. \\
\hline
为了 + V，S + 应该 + V & \zh{为了解决这个问题，我们应该加强教育。} \emph{Wèile jiějué zhège wèntí, wǒmen yīnggāi jiāqiáng jiàoyù.} & Pour résoudre ce problème, nous devons renforcer l'éducation. \\
\hline
S + 应该 + 做 + adj. + 的 + N & \zh{我们应该做诚实的公民。} \emph{Wǒmen yīnggāi zuò chéngshí de gōngmín.} & Nous devons être des citoyens honnêtes. \\
\hline
\end{tabularx}
\end{center}

\begin{center}
\begin{tabularx}{\linewidth}{|X|X|X|X|}
\hline
\rowcolor{popblueL} Caractères & Pinyin & Nature & Traduction \\
\hline
\zh{腐败} & fǔbài & nom / adj. & corruption ; corrompu \\
\hline
\zh{严重} & yánzhòng & adjectif & grave, sérieux \\
\hline
\zh{贪污} & tānwū & verbe & détourner des fonds \\
\hline
\zh{法律} & fǎlǜ & nom & loi, législation \\
\hline
\zh{严格} & yángé & adjectif & strict, rigoureux \\
\hline
\zh{够} & gòu & verbe / adv. & suffire, assez \\
\hline
\zh{解决} & jiějué & verbe & résoudre \\
\hline
\zh{应该} & yīnggāi & auxiliaire & devoir (recommandation) \\
\hline
\zh{加强} & jiāqiáng & verbe & renforcer \\
\hline
\zh{诚实} & chéngshí & adjectif & honnête \\
\hline
\zh{公民} & gōngmín & nom & citoyen \\
\hline
\zh{为了} & wèile & préposition & pour, afin de \\
\hline
\end{tabularx}
\end{center}

\courssub{Exercice résolu d'entraînement au thème}

\begin{exoresolu}[Thème guidé : trois phrases]
Traduisez en chinois : 1) « Parce que la loi n'est pas assez stricte, la corruption est un problème grave. » 2) « Nous devons être honnêtes. » 3) « Pour résoudre ce problème, nous devons renforcer la loi. »
\tcblower
\textbf{Étape 1-2 (analyse) :} la phrase 1 combine la cause (\zh{因为…所以…}) et « pas assez » (\zh{不够}) ; la phrase 2 utilise \zh{应该} + adjectif (sans \zh{是} !) ; la phrase 3 reprend \zh{为了} + \zh{应该}.\\
\textbf{Étape 3 (traduction) :}\\
1) \zh{因为法律不够严格，所以腐败是一个严重的问题。} \emph{Yīnwèi fǎlǜ bù gòu yángé, suǒyǐ fǔbài shì yí gè yánzhòng de wèntí.}\\
2) \zh{我们应该诚实。} \emph{Wǒmen yīnggāi chéngshí.} — l'adjectif \zh{诚实} fait office de verbe : pas de \zh{是}.\\
3) \zh{为了解决这个问题，我们应该加强法律。} \emph{Wèile jiějué zhège wèntí, wǒmen yīnggāi jiāqiáng fǎlǜ.}\\
\textbf{Étape 4 (vérification) :} tons de \zh{严格} (yángé, 2\textsuperscript{e} + 2\textsuperscript{e} ton), radical 讠dans \zh{诚}, virgule chinoise \zh{，} et point \zh{。} bien placés.
\end{exoresolu}

\courssub{Auto-correction : la grille finale avant de rendre la copie}

\begin{methode}[Vérifier sa traduction en trois points]
\begin{enumerate}
\item \textbf{Les tons} : relisez chaque mot à voix basse. Confondre \zh{mǎi} (买, acheter) et \zh{mài} (卖, vendre), ou écrire \zh{yángé} sans tons, coûte des points.
\item \textbf{Les radicaux} : vérifiez les caractères du thème : \zh{贪} (radical 贝, la coquille → l'argent), \zh{法} (radical 氵, l'eau), \zh{诚} (radical 讠, la parole). Un radical faux = un caractère faux.
\item \textbf{La ponctuation chinoise} : virgule \zh{，}, point \zh{。}, jamais de point français « . » ni de virgule « , » dans une phrase chinoise. Pas d'espace entre les caractères.
\end{enumerate}
\end{methode}

\begin{attention}[Bilan des erreurs fréquentes au thème]
\begin{itemize}
\item Oublier \zh{的} entre un adjectif dissyllabique et le nom : écrire \zh{严重问题} au lieu de \zh{严重的问题}.
\item Traduire « pas assez » par \zh{不很} au lieu de \zh{不够}.
\item Mettre \zh{是} devant un adjectif-verbe : \zh{法律是不严格} est faux ; dites \zh{法律不严格} ou \zh{法律不够严格}.
\item Placer l'adverbe après le verbe comme en français : « vouloir rapidement » se dit \zh{想很快地有钱}, l'adverbe \zh{很快} passe \emph{avant} le verbe.
\end{itemize}
\end{attention}

Vous disposez maintenant des cinq phrases modèles et d'une méthode fiable. La section suivante inversera le sens du travail : la \textbf{version}, du chinois vers le français.

\coursec{Modèle de texte commenté : Les facteurs de la corruption}

Pour produire un texte argumentatif simple en chinois (HSK 3-4), il faut maîtriser la structure « Constat $\to$ Causes énumérées $\to$ Solution prospective ». Le texte ci-dessous sert de modèle pour l'expression écrite (thème) et de source pour la version.

\begin{textebox}[Texte modèle : Les facteurs de la corruption]
腐败有很多原因。首先，一些人太爱钱了。其次，有的法律不够严格。最后，教育也很重要。为了建设一个诚实的社会，我们应该从小学习诚实。
\end{textebox}

\begin{popfigure}
\begin{tikzpicture}[
    mindmap,
    concept color=popblue,
    level 1 concept/.append style={level distance=3.5cm, sibling angle=120, font=\small},
    level 2 concept/.append style={level distance=2.8cm, sibling angle=60, font=\tiny},
    every node/.style={concept, execute at begin node=\hskip0pt},
    text=white
]
\node[align=center]{Texte argumentatif\\HSK 3-4}
    child[concept color=poporange] { node[align=center]{Constat initial\\腐败有很多原因} }
    child[concept color=popgreen] { node[align=center]{Énumération des causes\\连词 : 首先 / 其次 / 最后}
        child { node[align=center]{Cause 1 : Cupidité\\一些人太爱钱了} }
        child { node[align=center]{Cause 2 : Droit\\有的法律不够严格} }
        child { node[align=center]{Cause 3 : Éducation\\教育也很重要} }
    }
    child[concept color=poppurple] { node[align=center]{Proposition / But\\为了建设诚实的社会\\我们应该从小学习诚实} };
\end{tikzpicture}
\legende{Structure logique du texte modèle : du constat à la recommandation.}
\end{popfigure}

\begin{definition}[Analyse de la structure du modèle]
Le texte suit un plan en \textbf{5 segments} (5 phrases), typique de l'exposé clair :
\begin{enumerate}
    \item \textbf{Phrase d'amorce (Thèse) :} \zh{腐败有很多原因。} (Sujet + \zh{有} + Complément). Pose le sujet.
    \item \textbf{Argument 1 (Individuel) :} \zh{首先，一些人太爱钱了。} (Connecteur + Sujet + Structure d'excès \zh{太\ldots 了}).
    \item \textbf{Argument 2 (Institutionnel) :} \zh{其次，有的法律不够严格。} (Connecteur + Sujet \zh{有的} + Négation de degré \zh{不够}).
    \item \textbf{Argument 3 (Éducatif) :} \zh{最后，教育也很重要。} (Connecteur de clôture + \zh{也} pour l'addition).
    \item \textbf{Conclusion prospective (But + Devoir) :} \zh{为了建设一个诚实的社会，我们应该从小学习诚实。} (Préposition de but \zh{为了} antéposée + Modal \zh{应该} + Temps \zh{从小}).
\end{enumerate}
\end{definition}

\coursec{Structures et connecteurs utiles}

Maîtriser ces « briques » syntaxiques est indispensable pour réussir le thème et la version sur ce thème.

\begin{propriete}[Connecteurs logiques d'énumération (Ordre du discours)]
\begin{center}
\begin{tabularx}{\linewidth}{|X|X|X|X|}
\hline
\rowcolor{popblueL} \textbf{Chinois} & \textbf{Pinyin} & \textbf{Français} & \textbf{Registre / Emploi} \\
\hline
首先 & shǒuxiān & Premièrement, en premier lieu & Formel, écrit, ouvre l'énumération. \\
\hline
其次 & qícì & Deuxièmement, en second lieu & Formel, suit \zh{首先}. \\
\hline
然后 & ránhòu & Ensuite, puis & Moins formel, narration ou suite logique. \\
\hline
最后 & zuìhòu & Finalement, enfin, pour finir & Marque la fin de l'énumération. \\
\hline
另外 & lìngwài & Par ailleurs, de plus & Ajoute un point non prévu dans l'ordre strict. \\
\hline
\end{tabularx}
\end{center}
\end{propriete}

\begin{propriete}[Structures grammaticales clés du thème]
\begin{center}
\begin{tabularx}{\linewidth}{|X|X|X|}
\hline
\rowcolor{poporangeL} \textbf{Structure} & \textbf{Schéma / Règle} & \textbf{Exemple (Caractères + Pinyin + Traduction)} \\
\hline
\zh{有} (Existence) & Sujet + \zh{有} + Objet (Nombre + Nom) & \zh{腐败有很多原因。} (Fǔbài yǒu hěn duō yuányīn.) La corruption a de nombreuses causes. \\
\hline
\zh{太\ldots 了} (Excès) & Sujet + \zh{太} + Adj/V + \zh{了} $\to$ « Trop \ldots » (Péjoratif) & \zh{这道菜太辣了。} (Zhè dào cài tài là le.) Ce plat est \textbf{trop} épicé. \\
\hline
\zh{不够} (Insuffisance) & Sujet + \zh{不够} + Adj $\to$ « Pas assez \ldots » & \zh{时间不够多。} (Shíjiān bùgòu duō.) Le temps n'est pas suffisant. \\
\hline
\zh{有的} (Pronom indéfini) & \zh{有的} + Nom $\to$ « Certains (parmi d'autres) » & \zh{有的学生不喜欢写字。} (Yǒude xuésheng bù xǐhuan xiězì.) Certains élèves n'aiment pas écrire. \\
\hline
\zh{从小} (Temporel) & Sujet + \zh{从小} + Verbe $\to$ « Dès l'enfance / Depuis petit » & \zh{他从小喜欢足球。} (Tā cóng xiǎo xǐhuan zúqiú.) Il aime le foot depuis petit. \\
\hline
\zh{为了\ldots} (But) & \zh{为了} + Verbe/Objectif , Sujet + \zh{应该/要/必须} + Verbe & \zh{为了健康，我们要运动。} (Wèile jiànkāng, wǒmen yào yùndòng.) Pour la santé, nous devons faire du sport. \\
\hline
\zh{应该} (Devoir moral) & Sujet + \zh{应该} + Verbe $\to$ Conditionnel \emph{devrions} / \emph{il faut que} & \zh{我们应该帮助别人。} (Wǒmen yīnggāi bāngzhù biérén.) Nous devrions aider les autres. \\
\hline
\zh{建设} (Verbe abstrait) & \zh{建设} + Pays / Société / Avenir / Parti / Armée & \zh{建设美丽中国。} (Jiànshè měilì Zhōngguó.) Construire une belle Chine. \\
\hline
\end{tabularx}
\end{center}
\end{propriete}

\coursec{Écriture des caractères complexes du thème}

Les caractères ci-dessous sont « complexes » (nombre de traits élevé, structure emboîtée ou clé phonétique non évidente). Maîtrisez l'ordre des traits (règle générale : haut $\to$ bas, gauche $\to$ droite, horizontal $\to$ vertical, trait central $\to$ côtés, fermer en dernier).

\begin{center}
\begin{tabularx}{\linewidth}{|c|X|X|c|X|}
\hline
\rowcolor{popgreenL} \textbf{Car.} & \textbf{Pinyin / Sens} & \textbf{Radical (Clé)} & \textbf{Nb traits} & \textbf{Ordre des traits \& Décomposition} \\
\hline
腐 & fǔ (pourri, corrompre) & \zh{月} (viande/charme, 4 tr.) & 11 & 1. \zh{月} (haut-bas) 2. \zh{府} (fu) : \zh{广} (large) + \zh{付} (fu phonétique : \zh{人} + \zh{寸}). \textbf{Total 11}. \\
\hline
败 & bài (échec, corruption) & \zh{攵} (taper/action, 4 tr.) & 11 & 1. Partie sup. \zh{卜} (bu) + \zh{十} (shi) + \zh{一} (yi) 2. \zh{攵} (quatre traits : point, horiz, crochet, point). \textbf{Total 11}. \\
\hline
原 & yuán (origine, cause) & \zh{厂} (falaise/usine, 2 tr.) & 10 & 1. \zh{厂} (haut-gauche $\to$ bas) 2. \zh{泉} (quan, source) : \zh{白} (bai) + \zh{水} (shui). \textbf{Total 10}. \\
\hline
严 & yán (strict, sévère) & \zh{土} (terre, 3 tr.) & 7 & \textbf{Simplifié} : \zh{土} + \zh{工} (gong). \textit{Traditionnel : \zh{嚴} (bouche + phonétique)}. \textbf{Total 7}. \\
\hline
格 & gé (règle, standard) & \zh{木} (bois, 4 tr.) & 10 & 1. \zh{木} 2. \zh{各} (ge phonétique) : \zh{口} + \zh{夂} (pied qui traîne). \textbf{Total 10}. \\
\hline
建 & jiàn (construire, bâtir) & \zh{廴} (marche/chemin, 3 tr.) & 9 & 1. \zh{聿} (yu, pinceau) : traits 1-6 2. \zh{廴} (traits 7-9 : point, horiz, crochet final). \textbf{Total 9}. \\
\hline
设 & shè (établir, installer) & \zh{讠} (parole, 2 tr.) & 6 & \textbf{Simplifié} : \zh{讠} + \zh{殳} (shu, arme). \textit{Traditionnel : \zh{設} (parole + phonétique complex)}. \textbf{Total 6}. \\
\hline
诚 & chéng (honnête, sincère) & \zh{讠} (parole, 2 tr.) & 8 & \textbf{Simplifié} : \zh{讠} + \zh{成} (cheng, réussir). \textit{Traditionnel : \zh{誠} (parole + 成)}. \textbf{Total 8}. \\
\hline
实 & shí (réel, honnête) & \zh{宀} (toit, 3 tr.) & 8 & \textbf{Simplifié} : \zh{宀} + \zh{贞} (zhen) simplifié en \zh{贝} + \zh{工} ? Non, \zh{实} = \zh{宀} + \zh{贞} (simplifié). \textit{Traditionnel : \zh{實} (toit + 宀 + 貝 + ...)}. \textbf{Total 8}. \\
\hline
\end{tabularx}
\end{center}

\begin{experience}[Atelier écriture : « Dictée muette »]
\begin{enumerate}
    \item L'enseignant dicte le pinyin (ex: \emph{fǔbài}), les élèves écrivent les caractères sur ardoise.
    \item Vérification collective : ordre des traits, proportion, position du radical.
    \item Focus sur \zh{建} (clé \zh{廴} en bas à gauche, trait final long) et \zh{败} (clé \zh{攵} à droite, 4 traits précis).
    \item Exercice de « caractères à trous » : compléter le radical manquant (\zh{扌} vs \zh{攵} vs \zh{讠}).
\end{enumerate}
\end{experience}

\coursec{Caractères proches : ne pas confondre}

Ces paires partagent composants, prononciation proche ou sens voisin. La confusion est fréquente en rédaction.

\begin{center}
\begin{tabularx}{\linewidth}{|X|X|X|X|}
\hline
\rowcolor{poppurpleL} \textbf{Caractère cible (Leçon)} & \textbf{Caractère proche (Piège)} & \textbf{Différence clé (Graphie / Sens)} & \textbf{Exemple distinctif} \\
\hline
\zh{败} (bài) \newline Échec, corruption & \zh{费} (fèi) \newline Frais, dépenser & \zh{败} : clé \zh{攵} (action/taper). \zh{费} : clé \zh{贝} (argent/coquillage). & \zh{腐败} (corruption) vs \zh{费用} (frais/dépenses). \\
\hline
\zh{严} (yán) \newline Strict, sévère & \zh{研} (yán) \newline Rechercher, étudier & \zh{严} : haut \zh{土} (terre). \zh{研} : haut \zh{竹} (bambou) + bas \zh{石} (pierre). & \zh{严格} (strict) vs \zh{研究} (recherche/étudier). \\
\hline
\zh{设} (shè) \newline Établir, supposer & \zh{订} (dìng) \newline Fixer, réserver, commander & \zh{设} : clé \zh{讠} + \zh{殳} (arme). \zh{订} : clé \zh{讠} + \zh{丁} (clou/homme). & \zh{设计} (concevoir) vs \zh{订票} (réserver un billet). \\
\hline
\zh{诚} (chéng) \newline Honnête, sincère & \zh{成} (chéng) \newline Réussir, devenir, achevé & Même phonétique \zh{成}. \zh{诚} a clé \zh{讠} (parole). \zh{成} = \zh{戊} + \zh{丁}. & \zh{诚实} (honnête) vs \zh{成功} (réussite). \\
\hline
\zh{实} (shí) \newline Réel, solide, honnête & \zh{时} (shí) \newline Temps, heure, moment & \zh{实} : clé \zh{宀} (toit) + \zh{贞}. \zh{时} : clé \zh{日} (soleil) + \zh{寺} (temple). & \zh{诚实} (honnête) vs \zh{时间} (temps). \\
\hline
\zh{法} (fǎ) \newline Loi, méthode, France & \zh{去} (qù) \newline Aller, partir & \zh{法} : \zh{氵} (eau) + \zh{去} (phonétique). \zh{去} : \zh{土} + \zh{厶} (privé). & \zh{法律} (loi) vs \zh{去年} (l'an dernier). \\
\hline
\end{tabularx}
\end{center}

\begin{attention}[Pièges graphiques fréquents en examen]
\begin{itemize}
    \item \textbf{\zh{建} vs \zh{健} (jiàn) :} \zh{建} (construire, clé \zh{廴}) vs \zh{健} (sain, clé \zh{亻}). \emph{Ne pas écrire \zh{健设}}.
    \item \textbf{\zh{设} vs \zh{施} (shī) :} \zh{设} (établir, clé \zh{讠}) vs \zh{施} (appliquer, clé \zh{方}). \emph{Ne pas écrire \zh{施计}} pour \zh{设计}.
    \item \textbf{Clé \zh{讠} (parole) :} \zh{诚}, \zh{设}, \zh{课}, \zh{读}, \zh{语}, \zh{说}, \zh{订}, \zh{计}, \zh{论}, \zh{识}. Toujours à gauche, 2 traits (point, courbe horizontale).
    \item \textbf{Clé \zh{攵} (action) :} \zh{败}, \zh{教}, \zh{数}, \zh{敬}, \zh{政}, \zh{放}, \zh{攻}, \zh{敲}. Toujours à droite, 4 traits (点, 横, 撇, 点).
\end{itemize}
\end{attention}

\coursec{Thème : français → chinois}

Le thème exige de mobiliser le vocabulaire, les structures (connecteurs, 太...了, 不够, 为了...应该, 从小) et l'écriture correcte des caractères.

\begin{methode}[Traduire du français vers le chinois : la méthode en 4 étapes]
\begin{enumerate}
    \item \textbf{Analyser la phrase française} : Identifier le sujet, le verbe, les compléments, le temps, le registre. Repérer les connecteurs logiques.
    \item \textbf{Transposer en structure chinoise} : Appliquer l'ordre chinois (Sujet + Temps + Lieu + Manière + Verbe + Objet). Placer les connecteurs en tête de phrase. Choisir la structure grammaticale adéquate (\zh{太\ldots 了}, \zh{不够}, \zh{为了\ldots}).
    \item \textbf{Sélectionner le vocabulaire précis} : Choisir le bon mot (ex: \zh{建设} pas \zh{盖} pour société ; \zh{诚实} adj/nom ; \zh{从小} pas \zh{小时候} pour habitude passée/présente).
    \item \textbf{Écrire les caractères} : Respecter l'ordre des traits, les radicaux, la ponctuation chinoise (。 ， ；). Relire pour éviter les caractères proches.
\end{enumerate}
\end{methode}

\begin{exercice}[Entraînement au thème : Phrases à traduire]
Traduisez en chinois (caractères + pinyin) les phrases suivantes. Utilisez le vocabulaire et les structures de la leçon.
\begin{enumerate}
    \item La corruption a de nombreuses causes.
    \item Premièrement, certains fonctionnaires sont trop cupides (aiment trop l'argent).
    \item Deuxièmement, certaines lois ne sont pas assez strictes.
    \item Enfin, l'éducation morale est aussi très importante.
    \item Afin de bâtir une société honnête, nous devrions apprendre l'honnêteté dès l'enfance.
\end{enumerate}
\end{exercice}

\begin{corrige}[Exercice Thème]
\begin{enumerate}
    \item \zh{腐败有很多原因。} (Fǔbài yǒu hěn duō yuányīn.)
    \item \zh{首先，一些官员太爱钱了。} (Shǒuxiān, yīxiē guānyuán tài ài qián le.) \textit{Note : \zh{官员} pour fonctionnaire/cadre. \zh{太爱钱了} structure d'excès.}
    \item \zh{其次，有的法律不够严格。} (Qícì, yǒude fǎlǜ bùgòu yángé.) \textit{Note : \zh{有的} pour « certaines ». \zh{不够严格} pour « pas assez strictes ».}
    \item \zh{最后，道德教育也很重要。} (Zuìhòu, dàodé jiàoyù yě hěn zhòngyào.) \textit{Note : \zh{道德教育} (éducation morale). \zh{也} placé avant \zh{很}.}
    \item \zh{为了建设一个诚实的社会，我们应该从小学习诚实。} (Wèile jiànshè yí gè chéngshí de shèhuì, wǒmen yīnggāi cóng xiǎo xuéxí chéngshí.) \textit{Note : \zh{为了} en tête. \zh{建设} + \zh{诚实的社会}. \zh{应该} $\to$ devoir moral. \zh{从小} avant verbe. \zh{学习诚实} (VO).}
\end{enumerate}
\end{corrige}

% --- DEBUT DU BLOC VERSION (REVU ET AMELIORE) ---

\coursec{Version : chinois → français}

Après avoir travaillé le thème (français → chinois) qui consistait à exprimer les causes de la corruption en mobilisant le vocabulaire et les structures du module, nous abordons maintenant l'exercice inverse : la \cle{version}. Il ne s'agit pas de traduire mot à mot, mais de \textbf{restituer le sens} en français correct, naturel et adapté au registre (ici, un exposé clair et structuré). La version révèle votre compréhension fine de la syntaxe chinoise (ordre des mots, particules, connecteurs) et votre capacité à lever les ambiguïtés lexicales.

\courssub{Texte source et première approche}

Voici le texte original à traduire. Lisez-le une première fois à haute voix pour en saisir le rythme et la structure globale (énoncé de fait $\to$ énumération des causes $\to$ conclusion prospective).

\begin{textebox}[Texte à traduire : Les facteurs de la corruption]
腐败有很多原因。首先，一些人太爱钱了。其次，有的法律不够严格。最后，教育也很重要。为了建设一个诚实的社会，我们应该从小学习诚实。
\end{textebox}

\begin{center}
\begin{tabularx}{\linewidth}{|X|X|X|X|}
\hline
\rowcolor{popblueL} \textbf{Caractères} & \textbf{Pinyin} & \textbf{Nature} & \textbf{Traduction} \\
\hline
腐败 & fǔbài & n. / v. & corruption ; corrompre \\
原因 & yuányīn & n. & cause, raison \\
首先 & shǒuxiān & adv. & premièrement, en premier lieu \\
一些 & yīxiē & det. & certains, quelques-uns \\
太 \ldots 了 & tài \ldots le & struct. & trop \ldots (excès) \\
爱钱 & ài qián & v. + n. & aimer l'argent, être cupide \\
其次 & qícì & adv. & deuxièmement, en second lieu \\
有的 & yǒude & pron. & certains (parmi d'autres) \\
法律 & fǎlǜ & n. & loi, législation \\
不够 & bùgòu & adv. & ne pas être assez, insuffisant \\
严格 & yángé & adj. & strict, rigoureux \\
最后 & zuìhòu & adv. & finalement, en dernier lieu \\
教育 & jiàoyù & n. & éducation \\
也 & yě & adv. & aussi \\
重要 & zhòngyào & adj. & important \\
为了 & wèile & prep. & pour, afin de \\
建设 & jiànshè & v. & construire, édifier, bâtir \\
诚实 & chéngshí & adj. & honnête, sincère \\
社会 & shèhuì & n. & société \\
应该 & yīnggāi & v. aux. & devoir, falloir \\
从小 & cóng xiǎo & loc. adv. & dès l'enfance, dès le plus jeune âge \\
学习 & xuéxí & v. & apprendre, étudier \\
\hline
\end{tabularx}
\end{center}

\courssub{Méthodologie de la version : trois étapes}

La version suit un processus rigoureux en trois temps. Ne sautez pas l'étape 1 : les connecteurs logiques (首先, 其次, 最后, 为了) sont le squelette du texte.

\begin{methode}[Traduire un texte chinois vers le français : la méthode en 3 étapes]
\begin{enumerate}
    \item \textbf{Découper en segments de sens} : Repérez les connecteurs logiques (首先, 其次, 最后, 为了\ldots) et la ponctuation chinoise (。). Chaque proposition principale ou subordonnée = un segment. Numérotez-les.
    \item \textbf{Traduire segment par segment (brouillon)} : Traduisez \emph{littéralement} d'abord pour verrouiller le sens grammatical (sujet, verbe, compléments), puis proposez une \emph{équivalence française} correcte. Attention aux faux-amis (\zh{有的} $\neq$ « avoir de »), aux structures sans équivalent direct (\zh{太\ldots 了}, \zh{从小}) et à l'ordre des compléments (temps/lieu avant le verbe en chinois, souvent après en français).
    \item \textbf{Relire, lisser et accorder} : Assemblez les segments. Vérifiez les accords (genre, nombre), les temps (présent de vérité générale, conditionnel pour \zh{应该}), les articles (définis/indéfinis) et la fluidité (liaison des phrases par des connecteurs français : \emph{premièment / deuxièmement / enfin ; afin de / pour}). Relisez le texte final à voix haute : est-il naturel en français ?
\end{enumerate}
\end{methode}

\begin{popfigure}
\begin{tikzpicture}[
    node distance=1.2cm and 2cm,
    box/.style={draw=popblue, thick, rounded corners, fill=popblueL, text width=4.5cm, align=center, font=\small},
    arrow/.style={-Stealth, thick, popdark}
]
\node[box, align=center] (e1){\textbf{ÉTAPE 1}\\\textbf{Découper}\\(Segments + Connecteurs)};
\node[box, right=of e1, align=center] (e2){\textbf{ÉTAPE 2}\\\textbf{Traduire}\\(Littéral $\to$ Français correct)};
\node[box, right=of e2, align=center] (e3){\textbf{ÉTAPE 3}\\\textbf{Relire / Lisser}\\(Accords, Temps, Fluidité)};
\draw[arrow] (e1.east) -- (e2.west) node[midway, above, font=\tiny, popdark] {Repérage structurel};
\draw[arrow] (e2.east) -- (e3.west) node[midway, above, font=\tiny, popdark] {Reformulation};
\node[below=0.3cm of e1, font=\tiny, popmuted] {首先, 其次, 最后, 为了};
\node[below=0.3cm of e2, font=\tiny, popmuted] {有的, 太...了, 从小, 建设};
\node[below=0.3cm of e3, font=\tiny, popmuted] {Articles, Conditionnel, Connecteurs fr.};
\end{tikzpicture}
\legende{Frise des 3 étapes de la version : du texte chinois au français abouti.}
\end{popfigure}

\courssub{Analyse segment par segment et traduction commentée}

Appliquons la méthode au texte. Nous identifions \textbf{5 segments} (5 phrases chinoises).

\begin{definition}[Segment 1 : Constat initial]
\textbf{Chinois :} \zh{腐败有很多原因。} \\
\textbf{Pinyin :} Fǔbài yǒu hěn duō yuányīn. \\
\textbf{Structure :} Sujet (\zh{腐败}) + Verbe \zh{有} (avoir / il y a) + Complément d'objet (\zh{很多原因}). \\
\textbf{Analyse :} \zh{有} exprime l'existence ici (« il y a »). \zh{原因} est le nom tête. Pas de mot pour « il y a » en chinois, le verbe \zh{有} suffit. \\
\textbf{Traduction littérale :} La corruption a beaucoup de causes. \\
\textbf{Français lissé :} \textbf{La corruption a de nombreuses causes.} (Ou : \emph{Il y a de nombreuses causes à la corruption.})
\end{definition}

\begin{definition}[Segment 2 : Cause 1 — Cupidité individuelle]
\textbf{Chinois :} \zh{首先，一些人太爱钱了。} \\
\textbf{Pinyin :} Shǒuxiān, yīxiē rén tài ài qián le. \\
\textbf{Structure :} Connecteur \zh{首先} + Sujet \zh{一些人} + Structure d'excès \zh{太\ldots 了} + Verbe \zh{爱} + Objet \zh{钱}. \\
\textbf{Analyse :} \zh{首先} = « Premièrement / En premier lieu ». \zh{一些人} = « Certaines personnes / Quelques-uns ». \textbf{Piège majeur :} \zh{太\ldots 了} n'exprime pas l'intensité simple (« très ») mais l'\textbf{excès} (« trop »). \zh{爱钱} = « aimer l'argent » (souvent péjoratif : cupidité). La particule finale \zh{了} marque l'état réalisé ou le constat. \\
\textbf{Traduction littérale :} Premièrement, certaines personnes trop aimer argent [excès]. \\
\textbf{Français lissé :} \textbf{Premièrement, certaines personnes aiment trop l'argent.} (Registre soutenu : \emph{sont trop cupides}).
\end{definition}

\begin{attention}[Piège : \zh{太\ldots 了} = « Trop » (excès) et non « Très » (intensité)]
Ne traduisez \textbf{JAMAIS} \zh{太爱钱了} par « aiment \emph{très} l'argent » ou « aiment \emph{beaucoup} l'argent ».
\begin{itemize}
    \item \zh{很爱钱} (hěn ài qián) = aiment \emph{beaucoup} l'argent (intensité forte).
    \item \zh{太爱钱了} (tài ài qián le) = aiment \emph{trop} l'argent (excès, connotation négative, critique).
\end{itemize}
\emph{Exemple :} \zh{这道菜太辣了。} (Zhè dào cài tài là le.) $\to$ Ce plat est \textbf{trop} épicé (je ne peux pas le manger). Pas « très épicé ».
\end{attention}

\begin{definition}[Segment 3 : Cause 2 — Insuffisance légale]
\textbf{Chinois :} \zh{其次，有的法律不够严格。} \\
\textbf{Pinyin :} Qícì, yǒude fǎlǜ bùgòu yángé. \\
\textbf{Structure :} Connecteur \zh{其次} + Sujet \zh{有的法律} + Négation \zh{不} + Degré \zh{够} + Adjectif \zh{严格}. \\
\textbf{Analyse :} \zh{其次} = « Deuxièmement / En second lieu ». \textbf{Piège lexical :} \zh{有的} (yǒude) est un pronom indéfini signifiant « certains (parmi un ensemble) », « quelques-unes ». Il ne se traduit pas par « avoir de ». \zh{不够} + adj. = « ne pas être assez \ldots » / « insuffisamment \ldots ». \zh{严格} = strict. \\
\textbf{Traduction littérale :} Deuxièmement, certaines lois pas assez strictes. \\
\textbf{Français lissé :} \textbf{En second lieu, certaines lois ne sont pas assez strictes.} (Ou : \emph{la législation n'est pas suffisamment rigoureuse}).
\end{definition}

\begin{definition}[Segment 4 : Cause 3 — Éducation]
\textbf{Chinois :} \zh{最后，教育也很重要。} \\
\textbf{Pinyin :} Zuìhòu, jiàoyù yě hěn zhòngyào. \\
\textbf{Structure :} Connecteur \zh{最后} + Sujet \zh{教育} + Aussi \zh{也} + Degré \zh{很} + Adjectif \zh{重要}. \\
\textbf{Analyse :} \zh{最后} = « Finalement / En dernier lieu / Pour finir ». \zh{也} place l'éducation au même niveau d'importance que les causes précédentes. \zh{很} lie le sujet à l'adjectif (pas de verbe « être » en chinois). \\
\textbf{Traduction littérale :} Finalement, éducation aussi très importante. \\
\textbf{Français lissé :} \textbf{Enfin, l'éducation est aussi très importante.}
\end{definition}

\begin{definition}[Segment 5 : Finalité et Recommandation (Proposition conditionnelle)]
\textbf{Chinois :} \zh{为了建设一个诚实的社会，我们应该从小学习诚实。} \\
\textbf{Pinyin :} Wèile jiànshè yí gè chéngshí de shèhuì, wǒmen yīnggāi cóng xiǎo xuéxí chéngshí. \\
\textbf{Structure :} But \zh{为了} + Verbe \zh{建设} + Objet \zh{一个诚实的社会} (avec \zh{的} attributif), Sujet \zh{我们} + Modal \zh{应该} + Temps \zh{从小} + Verbe \zh{学习} + Objet \zh{诚实}. \\
\textbf{Analyse :} \\
\begin{itemize}
    \item \zh{为了\ldots} (Wèile\ldots) : Proposition de but antéposée. « Afin de / Pour \ldots ». Le sujet de la principale (\zh{我们}) est l'agent du but.
    \item \zh{建设} (jiànshè) : « Construire, bâtir, édifier » (souvent pour des choses abstraites : société, pays, avenir). \textbf{Piège :} Ne pas traduire par « construire une maison » seulement.
    \item \zh{诚实的社会} : \zh{的} marque l'attribut (société honnête).
    \item \zh{应该} (yīnggāi) : Devoir moral / conseil fort $\to$ \textbf{Conditionnel} en français (\emph{devrions / devons}).
    \item \zh{从小} (cóng xiǎo) : Locution adverbiale temporelle = « dès l'enfance », « dès le plus jeune âge », « depuis petit ». Place \textbf{avant} le verbe en chinois ; en français, souvent en fin de phrase ou après le sujet.
    \item \zh{学习诚实} : Apprendre l'honnêteté (nom) / apprendre à être honnête.
\end{itemize}
\textbf{Traduction littérale :} Pour construire un honnête société, nous devrions dès-petit apprendre honnêteté. \\
\textbf{Français lissé :} \textbf{Afin de bâtir une société honnête, nous devrions apprendre l'honnêteté dès le plus jeune âge.}
\end{definition}

\courssub{Traduction finale assemblée et commentée}

Voici le résultat de l'étape 3 (Relire / Lisser). Notez l'harmonisation des connecteurs (\emph{Premièrement / En second lieu / Enfin}) et l'usage du conditionnel pour la recommandation.

\begin{exoresolu}[Version complète commentée]
\textbf{Énoncé :} Traduire en français le texte suivant :
\zh{腐败有很多原因。首先，一些人太爱钱了。其次，有的法律不够严格。最后，教育也很重要。为了建设一个诚实的社会，我们应该从小学习诚实。}

\tcblower

\textbf{Solution détaillée :}

\textbf{La corruption a de nombreuses causes.} \\
\textbf{Premièrement, certaines personnes aiment trop l'argent.} \\
\textbf{En second lieu, certaines lois ne sont pas assez strictes.} \\
\textbf{Enfin, l'éducation est aussi très importante.} \\
\textbf{Afin de bâtir une société honnête, nous devrions apprendre l'honnêteté dès le plus jeune âge.}

\textbf{Commentaire du professeur :}
\begin{itemize}
    \item \textbf{Phrase 1 :} \emph{De nombreuses causes} (plus soutenu que \emph{beaucoup de causes}). \zh{有} $\to$ \emph{il y a / avoir}.
    \item \textbf{Phrase 2 :} \zh{首先} $\to$ \emph{Premièrement} (enumération formelle). \zh{太\ldots 了} $\to$ \emph{trop} (excès). \zh{一些人} $\to$ \emph{certaines personnes} (indéterminé).
    \item \textbf{Phrase 3 :} \zh{其次} $\to$ \emph{En second lieu} (variante de \emph{Deuxièmement}). \zh{有的} $\to$ \emph{certaines} (pronom/dét.). \zh{不够严格} $\to$ \emph{ne sont pas assez strictes} (accord f. pl. avec \emph{lois}).
    \item \textbf{Phrase 4 :} \zh{最后} $\to$ \emph{Enfin} (marque la fin de l'énumération). \zh{也} $\to$ \emph{aussi} (ajout). \zh{很} $\to$ \emph{très} (liaison adjectif).
    \item \textbf{Phrase 5 :} \zh{为了} $\to$ \emph{Afin de} (but, registre écrit). \zh{建设} $\to$ \emph{bâtir / édifier} (abstract). \zh{应该} $\to$ \emph{devrions} (conditionnel de conseil). \zh{从小} $\to$ \emph{dès le plus jeune âge} (standard) ou \emph{dès l'enfance}. \zh{学习诚实} $\to$ \emph{apprendre l'honnêteté} (nominalisation) ou \emph{apprendre à être honnête}.
\end{itemize}
\end{exoresolu}

\courssub{Tableau récapitulatif : Points de vigilance chinois $\to$ français}

\begin{center}
\begin{tabularx}{\linewidth}{|X|X|X|}
\hline
\rowcolor{poporangeL} \textbf{Élément chinois} & \textbf{Piège / Difficulté} & \textbf{Rendu français correct} \\
\hline
\zh{有的} (yǒude) & Ne pas confondre avec \zh{有} (avoir) + \zh{的} (particule). C'est un pronom : « certains d'entre eux ». & \emph{Certaines (lois), certains (gens), quelques-uns.} \\
\hline
\zh{太\ldots 了} (tài\ldots le) & Structure d'\textbf{excès}, pas d'intensité. Le \zh{了} final est obligatoire pour cette signification. & \emph{Trop \ldots} (ex: \emph{trop cher, trop loin}). Jamais « très ». \\
\hline
\zh{从小} (cóng xiǎo) & Complément de temps antéposé (avant le verbe). Signifie « depuis l'âge jeune ». & \emph{Dès le plus jeune âge, dès l'enfance, depuis petit.} \\
\hline
\zh{建设} (jiànshè) & Verbe « construire » pour des entités abstraites/collectives (pays, société, parti, armée). & \emph{Bâtir, édifier, construire (une société, un pays, l'avenir).} \\
\hline
\zh{应该} (yīnggāi) & Verbe modal : obligation morale, conseil, attente sociale. & \emph{Devoir / Falloir} $\to$ \textbf{Conditionnel présent} (\emph{devrions}) ou \emph{il faut que} + subjonctif. \\
\hline
\zh{为了\ldots} (wèile\ldots) & Préposition de but, toujours en tête de phrase (sujet partagé avec la principale). & \emph{Afin de / Pour / Dans le but de} + infinitif (si même sujet) / \emph{afin que} + subjonctif (si sujet diff.). \\
\hline
\end{tabularx}
\end{center}

\courssub{Exercices d'application : Version guidée}

Appliquez la méthode en 3 étapes (Découper $\to$ Traduire $\to$ Lisser) aux phrases ci-dessous. Elles réemploient le vocabulaire et les structures du thème (corruption, éducation, droit, connecteurs logiques).

\begin{exercice}[Entraînement à la version : Phrases isolées]
Traduisez en français les phrases suivantes. Respectez le registre (exposé clair) et les accords.
\begin{enumerate}
    \item \zh{为了减少腐败，政府必须加强法律的执行。} (Wèile jiǎnshǎo fǔbài, zhèngfǔ bìxū jiāqiáng fǎlǜ de zhíxíng.)
    \item \zh{有些官员利用职权谋取私利，这太不公平了。} (Yǒuxiē guānyuán lìyòng zhíquán móuqǔ sīlì, zhè tài bù gōngpíng le.)
    \item \zh{首先，我们要树立正确的价值观；其次，学校要加强道德教育。} (Shǒuxiān, wǒmen yào shùlì zhèngquè de jiàzhíguān; qícì, xuéxiào yào jiāqiáng dàodé jiàoyù.)
\end{enumerate}
\textit{Vocabulaire nouveau : \zh{减少} (jiǎnshǎo, réduire), \zh{政府} (zhèngfǔ, gouvernement), \zh{必须} (bìxū, devoir absolument), \zh{加强} (jiāqiáng, renforcer), \zh{执行} (zhíxíng, appliquer/exécuter), \zh{官员} (guānyuán, fonctionnaire/cadre), \zh{利用} (lìyòng, utiliser/abuser de), \zh{职权} (zhíquán, pouvoir/fonction), \zh{谋取} (móuqǔ, chercher à obtenir), \zh{私利} (sīlì, intérêt personnel), \zh{公平} (gōngpíng, juste/équitable), \zh{树立} (shùlì, établir/bâtir - abstrait), \zh{价值观} (jiàzhíguān, valeurs), \zh{道德} (dàodé, moralité/éthique).}
\end{exercice}

\begin{corrige}[Exercice Version]
\begin{enumerate}
    \item \textbf{Analyse :} But (\zh{为了减少腐败}) + Sujet (\zh{政府}) + Modal fort (\zh{必须}) + Verbe (\zh{加强}) + Objet (\zh{法律的执行}).
    \textbf{Traduction :} \textbf{Afin de réduire la corruption, le gouvernement doit impérativement renforcer l'application des lois.} (Ou : \emph{Pour diminuer la corruption, l'État doit absolument renforcer l'exécution de la loi.})
    
    \item \textbf{Analyse :} Sujet \zh{有些官员} (Certains fonctionnaires) + Verbe \zh{利用} + Objet \zh{职权} + Verbe \zh{谋取} + Objet \zh{私利} (construction en série V1 O1, V2 O2) ; Commentaire \zh{这太不公平了} (Sujet \zh{这} + Excès \zh{太\ldots 了} + Nég \zh{不} + Adj \zh{公平}).
    \textbf{Traduction :} \textbf{Certains fonctionnaires abusent de leur pouvoir pour servir leurs intérêts personnels, c'est trop injuste.} (Ou : \emph{...c'est vraiment trop injuste.}) \textit{Note : \zh{利用\ldots 谋取} $\to$ « abuser de \ldots pour \ldots » ; \zh{太不公平了} $\to$ « c'est trop injuste / vraiment injuste ».}
    
    \item \textbf{Analyse :} Deux propositions coordonnées par \zh{;} (point-virgule chinois). Prop 1 : \zh{首先} + \zh{我们要树立\ldots} (Volonté/Objectif). Prop 2 : \zh{其次} + \zh{学校要加强\ldots}.
    \textbf{Traduction :} \textbf{Premièrement, nous devons établir des valeurs correctes ; en second lieu, les écoles doivent renforcer l'éducation morale.} (Ou : \emph{Premièrement, il faut instaurer de justes valeurs ; deuxièmement, l'école doit renforcer l'enseignement de l'éthique.})
\end{enumerate}
\end{corrige}

\coursec{Production guidée et grille d'évaluation}

Cette activité finale mobilise l'ensemble des compétences de la leçon : vocabulaire, structures, connecteurs, écriture des caractères, cohérence textuelle.

\begin{experience}[Production écrite : « L'affiche de la semaine de l'intégrité »]
\textbf{Consigne :} Votre lycée organise une « Semaine de l'Intégrité ». Vous devez rédiger un court texte (environ \textbf{80 à 100 caractères chinois}) pour une affiche. Le texte doit :
\begin{enumerate}
    \item Présenter le problème : la corruption a des causes multiples.
    \item Citer \textbf{au moins 2 causes} (individuelle, légale, éducative) avec connecteurs (\zh{首先}, \zh{其次}, \zh{最后}).
    \item Proposer \textbf{une solution concrète} avec \zh{为了\ldots, 我们应该\ldots} ou \zh{必须\ldots}.
    \item Utiliser \textbf{au moins 5 caractères complexes} étudiés (\zh{腐, 败, 严, 格, 建, 设, 诚, 实}).
    \item Soigner l'écriture (caractères carrés, proportions, ordre des traits) et la ponctuation chinoise.
\end{enumerate}

\textbf{Plan suggéré :}
\begin{itemize}
    \item Phrase 1 : \zh{腐败有很多原因。}
    \item Phrase 2 : \zh{首先，\ldots 太\ldots 了。} (Cause 1)
    \item Phrase 3 : \zh{其次，有的\ldots 不够\ldots{}。} (Cause 2)
    \item Phrase 4 : \zh{为了\ldots{}，我们应该/必须\ldots} (Solution)
\end{itemize}

\textbf{Exemple de production élève (modèle) :}
\zh{腐败有很多原因。首先，一些人太贪心了。其次，有的法律不够严格。为了建设诚实社会，我们必须从小加强道德教育。}
(96 caractères - Niveau visé : HSK 3/4)
\end{experience}

\begin{aretenir}[Grille d'évaluation — Expression écrite (20 points)]
\begin{center}
\begin{tabularx}{\linewidth}{|X|c|X|}
\hline
\rowcolor{popgoldL} \textbf{Critère} & \textbf{Points} & \textbf{Indicateurs de réussite (Niveau Première A)} \\
\hline
\textbf{Contenu \& Structure} & 5 & Sujet respecté (causes + solution). Plan logique (Constat $\to$ Causes $\to$ Solution). Connecteurs utilisés correctement (\zh{首先/其次/最后/为了}). \\
\hline
\textbf{Vocabulaire} & 5 & Vocabulaire du module employé (\zh{腐败, 原因, 爱钱, 法律, 严格, 教育, 诚实, 社会, 建设}). Précision lexicale (pas de \zh{盖} pour \zh{建设}). Variété (au moins 8 items distincts). \\
\hline
\textbf{Grammaire \& Syntaxe} & 5 & Structures cibles maîtrisées : \zh{有} (existence), \zh{太\ldots 了} (excès), \zh{不够} (insuffisance), \zh{有的} (pronom), \zh{从小} (place avant verbe), \zh{应该/必须} (modal), \zh{为了} (but antéposé). Phrases complètes, ordre SOV correct. \\
\hline
\textbf{Caractères \& Écriture} & 3 & Caractères complexes écrits correctement (traits, radicaux, proportions). Pas de confusion caractères proches (\zh{败/费, 严/研, 设/订, 诚/成, 实/时}). Ponctuation chinoise correcte (。 ， ？). \\
\hline
\textbf{Cohérence \& Registre} & 2 & Texte fluide, ton approprié (exposé formel). Liaisons logiques claires. Pas de traduction mot à mot du français. \\
\hline
\textbf{TOTAL} & \textbf{20} &  \\
\hline
\end{tabularx}
\end{center}
\end{aretenir}

\courssub{Conseils pour l'examen (Épreuve écrite)}

\begin{aretenir}[Check-list Version (Épreuve de 2h / Bac)]
\begin{itemize}
    \item \textbf{Temps de lecture :} 5 min pour lire le texte chinois 2 fois (une fois global, une fois en segmentant).
    \item \textbf{Brouillon obligatoire :} Faites un tableau 3 colonnes : \emph{Segment chinois / Traduction littérale / Français lissé}.
    \item \textbf{Surveillez :} 
    \begin{itemize}
        \item Les \textbf{connecteurs logiques} (variété : \emph{Premièrement / D'abord / En premier lieu ; Deuxièmement / Ensuite / Par ailleurs ; Enfin / Pour terminer}).
        \item L'\textbf{accord des participes passés} et adjectifs (ex: \zh{法律} f. sg/pl ? $\to$ \emph{lois} f. pl $\to$ \emph{strictes}).
        \item Les \textbf{modaux} : \zh{应该/必须/要} $\to$ \emph{devoir / falloir / avoir à} (souvent conditionnel ou présent de nécessité).
        \item La place des \textbf{compléments de temps} (\zh{从小}, \zh{每天}, \zh{以前}) : en français, souvent en fin de phrase ou après le sujet, \textbf{jamais} entre le verbe et l'objet direct.
    \end{itemize}
    \item \textbf{Relisez votre français} sans regarder le chinois : est-ce que ça sonne « français » ?
\end{itemize}
\end{aretenir}

\coursec{Production guidée et grille d'évaluation}

Après avoir traduit du chinois vers le français (version) et du français vers le chinois (thème), vous êtes maintenant prêts pour l'étape la plus importante de cette leçon : \cle{produire votre propre texte} en chinois. C'est exactement ce qui vous sera demandé à l'examen : rédiger un court texte structuré, avec des caractères correctement écrits, sur un sujet de société. Nous allons procéder méthodiquement : d'abord observer un modèle complet, puis comprendre le plan imposé, puis rédiger, puis évaluer avec une grille officielle.

\courssub{Le sujet et le plan imposé}

\begin{definition}[Le sujet de production]
Rédigez un texte de \textbf{5 à 8 phrases} en chinois sur le thème : \zh{腐败的原因和解决办法} (\emph{Fǔbài de yuányīn hé jiějué bànfǎ}) « Les causes de la corruption et les solutions ». Vous devez utiliser les mots obligatoires suivants : \zh{腐败} (corruption), \zh{原因} (cause, raison), \zh{贪污} (détournement, corruption active), \zh{法律} (loi), \zh{严格} (strict), \zh{éducation} (éducation), \zh{诚实} (honnête), \zh{应该} (devoir, il faut).
\end{definition}

\begin{center}
\begin{tabularx}{\linewidth}{|X|X|X|X|}
\hline
\rowcolor{popblueL} Caractères & Pinyin & Nature & Traduction \\
\hline
腐败 & fǔbài & nom & corruption \\
\hline
原因 & yuányīn & nom & cause, raison \\
\hline
贪污 & tānwū & verbe / nom & détournement, corruption active \\
\hline
法律 & fǎlǜ & nom & loi \\
\hline
严格 & yángé & adjectif & strict, rigoureux \\
\hline
教育 & jiàoyù & nom / verbe & éducation, éduquer \\
\hline
诚实 & chéngshí & adjectif & honnête, sincère \\
\hline
应该 & yīnggāi & verbe modal & devoir, il faut \\
\hline
\end{tabularx}
\end{center}

Le plan est \textbf{imposé} : il comporte quatre blocs, chacun signalé par des connecteurs précis. Un correcteur doit pouvoir voir votre plan \emph{sans effort}, rien qu'en repérant les connecteurs.

\begin{center}
\begin{tabularx}{\linewidth}{|X|X|X|}
\hline
\rowcolor{popblueL} Bloc & Contenu & Connecteurs obligatoires \\
\hline
\zh{引言} (introduction) & 1 phrase : présenter le problème & \zh{现在} (aujourd'hui) \\
\hline
\zh{原因} (causes) & 2 à 3 phrases : les facteurs & \zh{首先} (d'abord), \zh{其次} (ensuite) \\
\hline
\zh{办法} (solutions) & 1 à 2 phrases : les remèdes & \zh{为了} (afin de), \zh{应该} (il faut) \\
\hline
\zh{结论} (conclusion) & 1 phrase : espoir ou bilan & \zh{我相信} (je crois que) \\
\hline
\end{tabularx}
\end{center}

\begin{popfigure}
\begin{tikzpicture}[
  bloc/.style={rectangle, rounded corners=4pt, draw=popblue, fill=popblueL, text width=4.6cm, align=center, minimum height=1.1cm, font=\small},
  fleche/.style={-{Stealth[length=2.5mm]}, thick, poporange}
]
\node[bloc, align=center] (b1) at (0,0){\textbf{引言} yǐnyán\\ \zh{现在，腐败是一个大问题。}};
\node[bloc, fill=poporangeL, draw=poporange, align=center] (b2) at (6.2,0){\textbf{原因} yuányīn\\ \zh{首先… 其次…}};
\node[bloc, fill=popgreenL, draw=popgreen, align=center] (b3) at (0,-2.4){\textbf{办法} bànfǎ\\ \zh{为了… 我们应该…}};
\node[bloc, fill=poppurpleL, draw=poppurple, align=center] (b4) at (6.2,-2.4){\textbf{结论} jiélùn\\ \zh{我相信… 越来越好。}};
\draw[fleche] (b1) -- (b2);
\draw[fleche] (b2) -- (b3);
\draw[fleche] (b3) -- (b4);
\end{tikzpicture}
\legende{Organigramme du plan imposé en quatre blocs : introduction, causes, solutions, conclusion.}
\end{popfigure}

\courssub{Observation : un modèle de texte complet}

Avant de rédiger, observons un texte modèle qui respecte exactement le plan et la banque de mots. Repérez les connecteurs et comptez les phrases.

\begin{textebox}[Modèle : 腐败的原因和解决办法]
现在，腐败是很多国家的大问题，我们国家也有这个问题。\\
\emph{Xiànzài, fǔbài shì hěn duō guójiā de dà wèntí, wǒmen guójiā yě yǒu zhège wèntí.}\\
« Aujourd'hui, la corruption est un grand problème dans beaucoup de pays, et notre pays connaît aussi ce problème. »\\[4pt]
首先，一些人贪污，因为他们想要很快有很多钱。其次，法律不够严格，所以这些人不害怕。\\
\emph{Shǒuxiān, yìxiē rén tānwū, yīnwèi tāmen xiǎng yào hěn kuài yǒu hěn duō qián. Qícì, fǎlǜ bù gòu yángé, suǒyǐ zhèxiē rén bù hàipà.}\\
« D'abord, certaines personnes se livrent à la corruption parce qu'elles veulent avoir beaucoup d'argent très vite. Ensuite, les lois ne sont pas assez strictes, donc ces personnes n'ont pas peur. »\\[4pt]
为了解决这个问题，我们应该有更严格的法律，也应该在学校教育年轻人要诚实。\\
\emph{Wèile jiějué zhège wèntí, wǒmen yīnggāi yǒu gèng yángé de fǎlǜ, yě yīnggāi zài xuéxiào jiàoyù niánqīngrén yào chéngshí.}\\
« Pour résoudre ce problème, nous devrions avoir des lois plus strictes et aussi éduquer les jeunes à l'école pour qu'ils soient honnêtes. »\\[4pt]
我相信我们的社会会越来越好。\\
\emph{Wǒ xiāngxìn wǒmen de shèhuì huì yuèlái yuè hǎo.}\\
« Je crois que notre société deviendra de mieux en mieux. »
\end{textebox}

\textbf{Analyse du modèle.} Le texte compte 6 phrases : 1 phrase d'introduction, 2 phrases de causes (avec \zh{首先} et \zh{其次}), 1 phrase de solutions (avec \zh{为了} et \zh{应该}), 1 phrase de conclusion. Les huit mots obligatoires sont tous présents. Chaque phrase est courte (10 à 20 caractères) et grammaticalement simple : c'est volontaire, et c'est ce qui rapporte des points.

\begin{exemplebox}[Phrases de conclusion réutilisables]
\zh{我相信我们的社会会越来越好。} (\emph{Wǒ xiāngxìn wǒmen de shèhuì huì yuèlái yuè hǎo.}) « Je crois que notre société deviendra de mieux en mieux. »\\
\zh{如果每个人都诚实，我们的国家会越来越强大。} (\emph{Rúguǒ měige rén dōu chéngshí, wǒmen de guójiā huì yuèlái yuè qiángdà.}) « Si chacun est honnête, notre pays deviendra de plus en plus fort. »\\
\zh{让我们一起努力！} (\emph{Ràng wǒmen yìqǐ nǔlì!}) « Efforçons-nous ensemble ! »\\
Remarquez la structure \zh{越来越} + adjectif (\emph{yuèlái yuè} + adj.) : « de plus en plus… ». C'est une structure de niveau HSK 3 très appréciée des correcteurs.
\end{exemplebox}

\courssub{La méthode de rédaction en 20 minutes}

\begin{methode}[Gagner des points à l'examen : la stratégie du texte simple et propre]
\begin{enumerate}
\item \textbf{2 minutes — Préparer.} Soulignez les mots obligatoires du sujet. Écrivez au brouillon le squelette : \zh{现在…} / \zh{首先…} / \zh{其次…} / \zh{为了…应该…} / \zh{我相信…}.
\item \textbf{12 minutes — Rédiger.} Écrivez des phrases \textbf{courtes et justes} plutôt que longues et fautives. Une phrase Sujet + Verbe + Complément bien construite vaut plus qu'une phrase complexe pleine d'erreurs. N'utilisez que des structures que vous maîtrisez (\zh{因为…所以…}, \zh{为了…应该…}, \zh{越来越}).
\item \textbf{3 minutes — Relecture grammaticale.} Vérifiez l'ordre des mots, la place des compléments de temps et de lieu avant le verbe, l'emploi de \zh{的}.
\item \textbf{3 minutes — Relecture graphique.} Vérifiez chaque caractère : trait manquant, mauvais radical, confusion avec un caractère proche (\zh{贪} / \zh{贫}, \zh{律} / \zh{津}). Écrivez dans les cases, un caractère par case, sans déborder.
\end{enumerate}
Règle d'or : un \textbf{plan visible} (connecteurs), des \textbf{caractères propres}, des \textbf{phrases courtes et justes} — voilà les trois sources de points.
\end{methode}

\begin{attention}[Chaque caractère mal écrit coûte des points]
À l'examen, la grille prévoit 4 points sur 20 pour l'écriture des caractères. Un trait manquant (\zh{诚} sans son trait de \zh{成}), un mauvais radical (\zh{腐} écrit avec la clé de la viande \zh{月} oubliée), un caractère confondu avec son voisin (\zh{贪污} écrit \zh{贫污}) : chaque faute est pénalisée. Réservez \textbf{toujours 3 minutes} à la relecture graphique, uniquement consacrée aux caractères. Rappel des pièges du thème : \zh{贪} (tān, clé \zh{贝} « l'argent ») ne s'écrit pas comme \zh{贫} (pín, pauvre) ; \zh{律} (lǜ, clé \zh{彳} « pas ») ne s'écrit pas comme \zh{津} (jīn, clé \zh{氵} « eau ») ; \zh{腐} comporte la clé \zh{月} (viande) à gauche, pas la clé \zh{疒} (maladie).
\end{attention}

\courssub{La grille d'évaluation sur 20 points}

Votre texte sera évalué, par votre professeur puis par un camarade, selon la grille officielle suivante. Apprenez-la par cœur : elle vous dit exactement où sont les points.

\begin{center}
\begin{tabularx}{\linewidth}{|X|p{1.5cm}|X|}
\hline
\rowcolor{popblueL} \textbf{Critère} & \textbf{Points} & \textbf{Ce que le correcteur regarde} \\
\hline
Contenu et respect du plan & /5 & Les 4 blocs présents, les 8 mots obligatoires utilisés, les idées logiques \\
\hline
Connecteurs & /4 & \zh{首先}, \zh{其次}, \zh{为了}, \zh{应该}, \zh{我相信} bien placés \\
\hline
Grammaire & /4 & Ordre des mots, \zh{因为…所以…}, \zh{的}, \zh{越来越} corrects \\
\hline
Caractères (traits et radicaux) & /4 & Traits complets, bons radicaux, pas de confusion \\
\hline
Ponctuation et présentation & /3 & 。 ， corrects, un caractère par case, texte lisible \\
\hline
\end{tabularx}
\end{center}

\begin{exoresolu}[Appliquer la grille : corriger un texte d'élève]
Énoncé. Voici le texte d'un élève : \zh{腐败是大问题。首先，一些人贪污，因为想有钱。其次，法律不严挌。为了问题，我们应该教育年轻人诚实。我相信社会越来越好。} Attribuez une note sur 20 avec la grille.
\tcblower
Solution détaillée. \textbf{Contenu et plan : 4/5} — les 4 blocs sont présents, mais le mot obligatoire \zh{严格} est mal écrit et \zh{原因} est absent. \textbf{Connecteurs : 4/4} — \zh{首先}, \zh{其次}, \zh{为了}, \zh{应该}, \zh{我相信} sont tous là et bien placés. \textbf{Grammaire : 3/4} — \zh{为了问题} est incorrect : \zh{为了} exige un verbe ou un objectif, il faut \zh{为了解决这个问题} ; de plus \zh{因为想有钱} est trop elliptique, mieux vaut \zh{因为他们想要有很多钱}. \textbf{Caractères : 2/4} — \zh{严挌} est une faute grave : le caractère correct est \zh{格} (clé \zh{木}, le bois), pas \zh{挌} ; \zh{解决} est omis. \textbf{Ponctuation et présentation : 3/3} — ponctuation correcte. \textbf{Total : 16/20.} Leçon à retenir : une seule faute de caractère sur un mot obligatoire coûte cher ; la relecture graphique aurait sauvé 2 points.
\end{exoresolu}

\courssub{Activité de classe : rédaction et relecture croisée}

\begin{experience}[Rédaction individuelle puis évaluation entre pairs]
\textbf{Étape 1 (20 minutes) — Rédaction individuelle.} Sur votre feuille de cases d'écriture (田字格), rédigez votre texte « \zh{腐败的原因和解决办法} » en 5 à 8 phrases, en suivant le plan imposé et en utilisant les 8 mots obligatoires. Un caractère par case, ponctuation pleine chasse dans sa propre case.\\
\textbf{Étape 2 (10 minutes) — Relecture croisée.} Échangez votre feuille avec votre voisin. Notez son texte avec la grille sur 20 : entourez en rouge les caractères fautifs, soulignez les connecteurs trouvés, cochez les mots obligatoires présents. Écrivez un commentaire bienveillant : un point fort et un conseil.\\
\textbf{Étape 3 (5 minutes) — Retour et correction.} Récupérez votre feuille, lisez les remarques, et réécrivez au propre les deux phrases les plus fautives.
\end{experience}

\begin{exercice}[À la maison : deuxième production]
Rédigez une deuxième version améliorée de votre texte, en intégrant les corrections de votre camarade. Contrainte supplémentaire : ajoutez une phrase avec la structure \zh{如果…就…} (si… alors…), par exemple : \zh{如果法律更严格，贪污的人就会更少。} (\emph{Rúguǒ fǎlǜ gèng yángé, tānwū de rén jiù huì gèng shǎo.}) « Si les lois sont plus strictes, il y aura moins de personnes corrompues. »
\end{exercice}

\begin{corrige}[Exercice « À la maison »]
Exemple de texte complet attendu (6 phrases) : \zh{现在，腐败的原因很多，这是一个大问题。首先，一些人贪污，因为他们想很快有很多钱。其次，法律不够严格，教育也不够。为了解决这个问题，我们应该有更严格的法律。如果每个人都诚实，我们的国家就会越来越好。我相信我们的社会会越来越好。} Vérifiez : les 8 mots obligatoires (\zh{腐败、原因、贪污、法律、严格、教育、诚实、应该}) sont tous présents, les connecteurs \zh{首先、其次、为了、应该、如果…就…、我相信} structurent le texte, et chaque phrase reste courte et correcte.
\end{corrige}

\begin{savaistu}[Comment le Bac juge l'expression écrite chinoise]
Au baccalauréat, l'expression écrite en chinois est jugée sur la \textbf{communication réelle}, pas sur l'ambition. Les correcteurs appliquent un principe simple : un texte simple, correct et bien structuré obtient toujours une meilleure note qu'un texte ambitieux mais incorrect. Un élève qui écrit 6 phrases courtes sans faute avec un plan clair décroche la mention ; un élève qui tente des phrases compliquées avec des caractères approximatifs perd des points sur presque tous les critères à la fois (grammaire, caractères, présentation). En Chine, les professeurs répètent aux élèves : \zh{写对的，不写难的。} (\emph{Xiě duì de, bù xiě nán de.}) « Écris ce qui est juste, pas ce qui est difficile. » C'est aussi la meilleure stratégie pour vous à Yaoundé, Douala ou Buea.
\end{savaistu}

\begin{aretenir}[L'essentiel de la production guidée]
\begin{itemize}
\item Plan imposé en 4 blocs : \zh{引言} → \zh{原因} (\zh{首先、其次}) → \zh{办法} (\zh{为了、应该}) → \zh{结论} (\zh{我相信}).
\item 8 mots obligatoires : \zh{腐败、原因、贪污、法律、严格、教育、诚实、应该}.
\item Gestion du temps : 2 min de préparation, 12 min de rédaction, 3 min de relecture grammaticale, 3 min de relecture graphique.
\item Grille sur 20 : contenu et plan /5, connecteurs /4, grammaire /4, caractères /4, ponctuation et présentation /3.
\item Stratégie gagnante : phrases courtes et justes, plan visible, caractères propres — \zh{写对的，不写难的} !
\end{itemize}
\end{aretenir}

\newpage

\coursec{Activité d'intégration}

\coursec{Leçon 30 — Expression écrite : Facteurs de la corruption ; Caractères complexes}

\courssub{Objectifs de l'activité}
\noindent À l'issue de cette activité d'intégration, l'élève doit être capable de :
\begin{itemize}
    \item \textbf{Mobiliser} le lexique spécifique de la corruption (\cle{腐败 fǔbài}, \cle{贿赂 huìlù}, \cle{贪污 tānwū}, \cle{监督 jiāndū}, \cle{透明度 tòumíngdù}, \cle{公务员 gōngwùyuán}, \cle{制度 zhìdù}, \cle{加强 jiāqiáng}, \cle{增加 zēngjiā}, \cle{未来 wèilái}) dans un texte cohérent.
    \item \textbf{Structurer} un texte argumentatif simple en utilisant les connecteurs logiques étudiés : \cle{首先 shǒuxiān} (premier), \cle{其次 qícì} (deuxièmement), \cle{为了 wèile} (afin de / dans le but de), \cle{应该 yīnggāi} (devoir / il faut que).
    \item \textbf{Écrire correctement} les caractères complexes du thème (ordre des traits, radicaux, distinction des caractères proches) sur des cases d'écriture standardisées (\emph{Tiangezhi} 田字格).
    \item \textbf{S'auto-évaluer et évaluer un pair} à l'aide d'une grille de critères (contenu, structure, caractères, grammaire, ponctuation).
    \item \textbf{Présenter} son travail à l'oral avec une prononciation correcte (tons, liaison) lors de la lecture des meilleures productions.
\end{itemize}

\courssub{Situation}
\noindent Tu t'appelles \textbf{Ayuk} (chinois : \zh{阿育 Āyù}), élève en classe de Première A au \textbf{Lycée Bilingue d'Etoug-Ebe} à \textbf{Yaoundé}. Tu participes à un programme d'échange épistolaire avec le \textbf{Lycée n°4 de Nankin} (南京第四中学 \zh{Nánjīng Dì Sì Zhōngxué}). Ton correspondant, \textbf{Lǐ Wěi} (李伟), t'a écrit pour te demander ton avis sur les grands défis sociaux du Cameroun. Tu as décidé de lui répondre en rédigeant une lettre manuscrite sur le thème de la lutte contre la corruption, un sujet d'actualité brûlant au Cameroun comme en Chine. Tu veux lui montrer que la jeunesse camerounaise est lucide et engagée.

\begin{textebox}[Extrait de l'e-mail reçu de Li Wei (version simplifiée pour niveau HSK 2-3)]
\shortstack[l]{\zh{亲爱的阿育：}\\ \zh{你好！}\\ \zh{我是李伟。我听说喀麦隆足球很厉害，雅温得也很美。}\\ \zh{请问：你们国家现在有什么社会问题？}\\ \zh{我想了解真实的喀麦隆。}\\ \zh{祝学习进步！}\\ \zh{李伟}\\ \zh{2024年10月15日}}
\par\medskip
\emph{Qīn'ài de Āyù: Nǐ hǎo! Wǒ shì Lǐ Wěi. Wǒ tīngshuō Kāmàilóng zúqiú hěn lìhai, Yǎwēndé yě hěn měi. Qǐngwèn: nǐmen guójiā xiànzài yǒu shénme shèhuì wèntí? Wǒ xiǎng liǎojiě zhēnshí de Kāmàilóng. Zhù xuéxí jìnbù! Lǐ Wěi 2024 nián 10 yuè 15 rì.}
\par\medskip
\textbf{Traduction :} « Cher Ayuk, Bonjour ! Je suis Li Wei. J'ai entendu dire que le football camerounais est très fort et que Yaoundé est très belle. Question : quel est le problème social actuel dans votre pays ? Je veux connaître le vrai Cameroun. Bonne progression dans vos études ! Li Wei, 15 octobre 2024. »
\end{textebox}

\noindent Ta mission : rédiger la lettre de réponse manuscrite (sur papier à cases) que tu enverras à Li Wei.

\courssub{Consignes}
\noindent Le travail se déroule en \textbf{quatre phases} sur 2 heures.

\begin{enumerate}
    \item \textbf{Phase 1 — Préparation et rédaction (45 min)} \\
    \begin{itemize}
        \item Rédige ton brouillon au crayon sur ton cahier. La lettre doit contenir \textbf{6 à 8 phrases} complètes.
        \item \textbf{Structure imposée :}
        \begin{enumerate}
            \item Formules d'appel et de politesse d'ouverture (\zh{亲爱的李伟：}, \zh{你好！}).
            \item Présentation du problème : la corruption est un problème grave au Cameroun.
            \item Cause 1 introduite par \cle{首先} (ex. : les salaires sont bas).
            \item Cause 2 introduite par \cle{其次} (ex. : la surveillance est faible).
            \item Solution introduite par \cle{为了……，应该……} (ex. : augmenter les salaires, renforcer les lois).
            \item Vœu pour l'avenir (ex. : j'espère que le Cameroun sera plus propre/prospère).
            \item Formules de clôture et signature (\zh{此致}, \zh{敬礼}, \zh{阿育}, \zh{日期}).
        \end{enumerate}
        \item Recopie ta version finale \textbf{au stylo noir} sur la feuille à cases d'écriture fournie (format \emph{Tiangezhi} 田字格). Soigne l'ordre des traits, la proportion des composants et la ponctuation pleine chasse (\zh{， 。 ！ ：}).
    \end{itemize}

    \item \textbf{Phase 2 — Relecture croisée et évaluation par les pairs (30 min)} \\
    \begin{itemize}
        \item Échange ta lettre avec ton voisin de table.
        \item Utilise la \textbf{Grille d'évaluation} (fournie ci-après) pour noter la lettre de ton camarade. Coche les cases, écris un commentaire constructif en français ou en chinois simple au dos de la feuille.
        \item Rends la lettre annotée à son auteur.
    \end{itemize}

    \item \textbf{Phase 3 — Correction finale et affichage (20 min)} \\
    \begin{itemize}
        \item Prends connaissance des remarques de ton pair. Corrige tes erreurs (caractères mal formés, ton manquant au pinyin si tu l'as noté, connecteur mal utilisé) sur une \textbf{nouvelle feuille propre} (version définitive pour affichage).
        \item Remets la version finale au professeur.
    \end{itemize}

    \item \textbf{Phase 4 — Lecture à voix haute et débat court (25 min)} \\
    \begin{itemize}
        \item Le professeur sélectionne 3 à 4 lettres variées. Les auteurs les lisent devant la classe.
        \item La classe note les arguments entendus. Mini-débat dirigé par le prof : « Selon vous, quelle solution est la plus réaliste ? Pourquoi ? » (en français ou chinois simple).
        \item Affichage des meilleures lettres au « Mur de la Chine » de la classe.
    \end{itemize}
\end{enumerate}

\courssub{Ressources à mobiliser}
\noindent Voici les outils linguistiques et graphiques à garder sous les yeux pendant la rédaction.

\begin{definition}[Lexique clé du thème « Corruption » (HSK 3-4)]
\begin{center}
\begin{tabularx}{\linewidth}{|X|X|X|X|}
\hline
\rowcolor{popblueL} \textbf{Caractères} & \textbf{Pinyin} & \textbf{Nature} & \textbf{Traduction} \\ \hline
腐败 & fǔbài & V / N & Corruption (pourrir + perdre) ; corrompu \\ \hline
贿赂 & huìlù & V / N & Corrompre (pots-de-vin) / pots-de-vin \\ \hline
贪污 & tānwū & V & Détourner des fonds publics ; malversation \\ \hline
贪官 & tānguān & N & Fonctionnaire corrompu (lit. officiel avide) \\ \hline
监督 & jiāndū & V / N & Surveiller, superviser / surveillance, contrôle \\ \hline
透明度 & tòumíngdù & N & Transparence (degré de transparence) \\ \hline
薪水 & xīnshuǐ & N & Salaire, paye \\ \hline
法律 & fǎlǜ & N & Loi, législation \\ \hline
严格 & yángé & Adj & Strict, rigoureux \\ \hline
干净 & gānjìng & Adj & Propre, honnête (fig.) \\ \hline
社会 & shèhuì & N & Société \\ \hline
问题 & wèntí & N & Problème \\ \hline
解决 & jiějué & V & Résoudre (un problème) \\ \hline
希望 & xīwàng & V / N & Espérer / espoir \\ \hline
公务员 & gōngwùyuán & N & Fonctionnaire \\ \hline
制度 & zhìdù & N & Système, régime, institution \\ \hline
加强 & jiāqiáng & V & Renforcer, consolider \\ \hline
增加 & zēngjiā & V & Augmenter, accroître \\ \hline
未来 & wèilái & N & Avenir, futur \\ \hline
\end{tabularx}
\end{center}
\end{definition}

\begin{propriete}[Connecteurs logiques pour l'argumentation (Niveau HSK 3)]
\begin{center}
\begin{tabularx}{\linewidth}{|X|X|X|}
\hline
\rowcolor{popgreenL} \textbf{Structure / Connecteur} & \textbf{Exemple en caractères + Pinyin} & \textbf{Traduction} \\ \hline
\cle{首先} \zh{，} Sujet + Verbe/Adj & \zh{首先，薪水太低。} (Shǒuxiān, xīnshuǐ tài dī.) & Premièrement, les salaires sont trop bas. \\ \hline
\cle{其次} \zh{，} Sujet + Verbe/Adj & \zh{其次，监督不严格。} (Qícì, jiāndū bù yángé.) & Deuxièmement, la surveillance n'est pas stricte. \\ \hline
\cle{为了} + But \zh{，} Sujet + \cle{应该} / \cle{要} / \cle{必须} + Verbe & \zh{为了解决腐败，政府应该增加薪水。} (Wèile jiějué fǔbài, zhèngfǔ yīnggāi zēngjiā xīnshuǐ.) & Pour résoudre la corruption, le gouvernement devrait augmenter les salaires. \\ \hline
\cle{另外} \zh{，} ... & \zh{另外，法律必须严格。} (Lìngwài, fǎlǜ bìxū yángé.) & De plus, la loi doit être stricte. \\ \hline
\cle{总之} \zh{，} ... & \zh{总之，我们要反对腐败。} (Zǒngzhí, wǒmen yào fǎnduì fǔbài.) & En résumé, nous devons nous opposer à la corruption. \\ \hline
\end{tabularx}
\end{center}
\end{propriete}

\begin{attention}[Pièges francophones sur les connecteurs]
\begin{itemize}
    \item \textbf{Ordre des mots :} En chinois, \cle{为了} (but) se place \textbf{AVANT} le sujet (\zh{为了目的，主语应该...}). En français, « Pour + infinitif » est souvent en début de phrase mais la structure chinoise est rigide : \zh{为了 + [But], [Sujet] + 应该/必须 + [Action]}.
    \item \textbf{首先 / 其次} sont des adverbes de séquence : ils occupent la position \textbf{Thème} (début de phrase), suivis d'une virgule \zh{，}. Ne pas dire \zh{我首先认为...} (moi d'abord je pense) sauf si le sujet est contrasté. Préférer \zh{首先，我认为...}.
    \item \textbf{应该} (devoir moral/conseil) vs \textbf{必须} (obligation forte/contrainte) vs \textbf{要} (volonté/obligation future). Pour une proposition de solution, \cle{应该} est le plus adapté.
    \item \textbf{Sandhi tonal (rappel) :} \zh{一} (yī) devient \emph{yí} devant un 4e ton (\zh{一个} \rightarrow \emph{yí gè}) et \emph{yì} devant 1er, 2e, 3e ton. \zh{不} (bù) devient \emph{bú} devant un 4e ton (\zh{不够} \rightarrow \emph{bú gòu}). En pinyin standard, on écrit la forme de citation (yī, bù) ; à l'oral, on applique le sandhi.
\end{itemize}
\end{attention}

\begin{propriete}[Écriture des caractères complexes du thème : Analyse graphémique]
\noindent Maîtriser l'écriture, c'est connaître le \textbf{radical} (clé de dictionnaire), l'\textbf{ordre des traits} et la \textbf{structure} (gauche-droite, haut-bas, enveloppement).
\begin{center}
\begin{tabularx}{\linewidth}{|X|X|X|X|}
\hline
\rowcolor{poporangeL} \textbf{Caractère} & \textbf{Radical (Clé)} & \textbf{Nb traits} & \textbf{Structure \& Astuce mnémotechnique} \\ \hline
\zh{腐} fǔ & \zh{月} (viande/charcuterie) & 13 & Gauche-Droite. \textbf{月} + \zh{府} (fǔ, bureau/gouvernement). \emph{La viande pourrit dans les bureaux $\rightarrow$ corruption.} \\ \hline
\zh{败} bài & \zh{攵} (pū, frapper/action) & 11 & Gauche-Droite. \zh{貝} (coquillage/monnaie) + \zh{攵}. \emph{Perte de l'argent $\rightarrow$ échec, pourriture.} \\ \hline
\zh{监} jiān & \zh{皿} (mǐn, récipient) & 13 & Haut-Bas. \zh{臥} (wò, coucher) au-dessus de \zh{皿}. \emph{Regarder de haut dans un récipient $\rightarrow$ surveiller.} \\ \hline
\zh{督} dū & \zh{目} (mù, œil) & 13 & Gauche-Droite. \zh{目} + \zh{叔} (shū, oncle). \emph{L'œil de l'oncle surveille.} \\ \hline
\zh{贿} huì & \zh{贝} (bèi, coquillage/argent) & 13 & Haut-Bas. \zh{穴} (xué, trou/caverne) au-dessus de \zh{贝}. \emph{L'argent dans un trou caché $\rightarrow$ pot-de-vin.} \\ \hline
\zh{赂} lù & \zh{贝} (bèi) & 13 & Gauche-Droite. \zh{贝} + \zh{各} (gè). \emph{Chacun (各) donne de l'argent (贝) $\rightarrow$ pots-de-vin.} \\ \hline
\zh{贪} tān & \zh{贝} (bèi) & 11 & Haut-Bas. \zh{今} (jīn, maintenant) au-dessus de \zh{贝}. \emph{Vouloir l'argent maintenant $\rightarrow$ avide.} \\ \hline
\zh{污} wū & \zh{氵} (shuǐ, eau) & 6 & Gauche-Droite. \zh{氵} + \zh{于} (yú). \emph{Eau sale.} \\ \hline
\zh{严} yán & \zh{山} (shān, montagne) & 17 & Haut-Bas. \zh{山} + \zh{彐} + \zh{口} + \zh{工} + \zh{土}. \emph{Montagne haute et escarpée $\rightarrow$ strict.} \\ \hline
\zh{格} gé & \zh{木} (mù, bois) & 10 & Haut-Bas. \zh{木} + \zh{各}. \emph{Le bois a des normes (grilles).} \\ \hline
\end{tabularx}
\end{center}
\end{propriete}

\begin{aretenir}[Ordre des traits : Règles d'or pour ces caractères]
\begin{enumerate}
    \item \textbf{Haut $\rightarrow$ Bas} : \zh{监} (écrire le haut \zh{臥} puis \zh{皿}), \zh{贿} (\zh{穴} puis \zh{贝}), \zh{贪} (\zh{今} puis \zh{贝}), \zh{严} (\zh{山} puis \zh{彐} puis \zh{口} puis \zh{工} puis \zh{土}), \zh{格} (\zh{木} puis \zh{各}).
    \item \textbf{Gauche $\rightarrow$ Droite} : \zh{腐} (\zh{月} puis \zh{府}), \zh{败} (\zh{貝} puis \zh{攵}), \zh{督} (\zh{目} puis \zh{叔}), \zh{赂} (\zh{贝} puis \zh{各}), \zh{污} (\zh{氵} puis \zh{于}).
    \item \textbf{Horizontal avant Vertical} : Crucial pour \zh{月} (traits horizontaux avant le trait vertical fermant), \zh{木}, \zh{土}, \zh{工}, \zh{贝}.
    \item \textbf{Enveloppement} : \zh{穴} (haut et côtés avant le trait final de fermeture bas), \zh{臥} (cadre avant contenu), \zh{广} (dans \zh{府} : haut et côté gauche avant le contenu).
\end{enumerate}
\end{aretenir}

\begin{attention}[Caractères proches : Ne pas confondre ! (Exercices de discrimination visuelle)]
\begin{center}
\begin{tabularx}{\linewidth}{|X|X|X|}
\hline
\rowcolor{poppinkL} \textbf{Paire} & \textbf{Différence graphique} & \textbf{Astuce} \\ \hline
\zh{腐} (fǔ, pourrir) \textbf{vs} \zh{府} (fǔ, bureau/gouvernement) & \zh{腐} a le radical \zh{月} (viande) à gauche. \zh{府} n'a pas de radical à gauche (c'est le phonétique de \zh{腐}). & \emph{La viande (\zh{月}) pourrit dans le bureau (\zh{府}).} \\ \hline
\zh{监} (jiān, surveiller) \textbf{vs} \zh{鉴} (jiàn, miroir/réfléchir) & \zh{监} : haut = \zh{臥} (coucher). \zh{鉴} : haut = \zh{金} (métal) + \zh{白} (blanc). & \emph{Surveiller (\zh{监}) depuis un lit (\zh{臥}) ; Miroir (\zh{鉴}) en métal (\zh{金}) blanc.} \\ \hline
\zh{督} (dū, superviser) \textbf{vs} \zh{叔} (shū, oncle) & \zh{督} = \zh{目} (œil) + \zh{叔}. \zh{叔} est le phonétique. & \emph{L'œil (\zh{目}) de l'oncle (\zh{叔}) surveille.} \\ \hline
\zh{贿} (huì, pot-de-vin) \textbf{vs} \zh{穴} (xué, trou) & \zh{贿} = \zh{穴} (haut) + \zh{贝} (bas). \zh{穴} est seul. & \emph{L'argent (\zh{贝}) caché dans un trou (\zh{穴}).} \\ \hline
\zh{贪} (tān, avide) \textbf{vs} \zh{今} (jīn, maintenant) & \zh{贪} = \zh{今} (haut) + \zh{贝} (bas). & \emph{Vouloir l'argent (\zh{贝}) maintenant (\zh{今}).} \\ \hline
\zh{严} (yán, strict) \textbf{vs} \zh{岩} (yán, rocher/falaise) & \zh{严} : radical \zh{山} en haut (structure H-B). \zh{岩} : radical \zh{山} à gauche (structure G-D) + \zh{石}. & \emph{Strict (\zh{严}) = Montagne haute (haut). Rocher (\zh{岩}) = Montagne + Pierre (gauche).} \\ \hline
\end{tabularx}
\end{center}
\end{attention}

\begin{methode}[Comment rédiger une lettre formelle simple en chinois (Format standard)]
\begin{enumerate}
    \item \textbf{En-tête (aligné à gauche, haut de page) :} \zh{亲爱的李伟：} (Cher Li Wei, avec deux-points pleins \zh{：}). \textit{Ne pas mettre de virgule.}
    \item \textbf{Corps du texte :} Commencer sur la ligne suivante, \textbf{rentrant de deux cases} (deux caractères) pour le premier paragraphe. Phrases séparées par \zh{。}. Virgules \zh{，} pour les clauses.
    \item \textbf{Formules de politesse de fin (alignées à gauche, bas de page) :}
    \begin{itemize}
        \item \zh{此致} (Cǐ zhì — « Je m'arrête ici ») : aligné à gauche, ligne nouvelle.
        \item \zh{敬礼} (Jìng lǐ — « Salutations respectueuses ») : sur la ligne d'après, \textbf{rentrant de deux cases} par rapport à \zh{此致}.
    \end{itemize}
    \item \textbf{Signature et date (alignés à droite) :}
    \begin{itemize}
        \item \zh{阿育} (Ton nom chinois).
        \item \zh{2024年11月20日} (Date : Année + 月 + Jour + 日).
    \end{itemize}
\end{enumerate}
\end{methode}

\courssub{Proposition de production et corrigé commenté}
\noindent Voici une production modèle respectant toutes les contraintes (longueur, structure, caractères, connecteurs). Elle sert de référence pour l'enseignant et de support de correction collective.

\begin{exoresolu}[Lettre modèle : « Non à la corruption ! »]
\textbf{Énoncé :} Rédige la lettre de réponse à Li Wei (6-8 phrases, structure imposée, caractères soignés).
\tcblower
\textbf{Production modèle (Version finale pour affichage) :}

\begin{textebox}[Lettre manuscrite d'Ayuk à Li Wei — Version calligraphiée sur cases]
\shortstack[l]{\zh{亲爱的李伟：}\\ \zh{你好！}\\ \zh{喀麦隆的腐败问题是一个严重的社会问题。}\\ \zh{首先，公务员的薪水太低。}\\ \zh{其次，监督制度不够严格。}\\ \zh{为了解决这个问题，政府应该增加薪水，加强法律监督。}\\ \zh{我希望喀麦隆的未来更干净、更美好。}\\ \zh{此致}\\ \zh{敬礼}\\ \zh{阿育}\\ \zh{2024年11月20日}}
\par\medskip
\textbf{Pinyin :} \\
Qīn'ài de Lǐ Wěi: \\
Nǐ hǎo! \\
Kāmàilóng de fǔbài wèntí shì yī gè yánzhòng de shèhuì wèntí. \\
Shǒuxiān, gōngwùyuán de xīnshuǐ tài dī. \\
Qícì, jiāndū zhìdù bù gòu yángé. \\
Wèile jiějué zhège wèntí, zhèngfǔ yīnggāi zēngjiā xīnshuǐ, jiāqiáng fǎlǜ jiāndū. \\
Wǒ xīwàng Kāmàilóng de wèilái gèng gānjìng, gèng měihǎo. \\
Cǐ zhì \\
Jìng lǐ \\
Ā yù \\
2024 nián 11 yuè 20 rì
\par\medskip
\textbf{Traduction française :} \\
Cher Li Wei, \\
Bonjour ! \\
Le problème de la corruption au Cameroun est un problème social grave. \\
Premièrement, les salaires des fonctionnaires sont trop bas. \\
Deuxièmement, le système de surveillance n'est pas assez strict. \\
Pour résoudre ce problème, le gouvernement devrait augmenter les salaires et renforcer la surveillance légale. \\
J'espère que l'avenir du Cameroun sera plus propre (honnête) et plus beau. \\
Salutations respectueuses, \\
Ayuk \\
20 novembre 2024
\end{textebox}

\textbf{Commentaire linguistique et pédagogique (pour l'enseignant / auto-correction élève) :}

\begin{enumerate}
    \item \textbf{Format épistolaire respecté :} Appel \zh{亲爱的李伟：} (deux-points pleins). Retrait de deux cases pour le premier paragraphe (\zh{喀麦隆...}). Formules de clôture \zh{此致 / 敬礼} alignées correctement (Cǐzhì à gauche, Jìnglǐ rentré de deux cases). Signature et date à droite.
    \item \textbf{Vocabulaire ciblé utilisé :} \zh{腐败} (fǔbài), \zh{社会问题} (shèhuì wèntí), \zh{公务员} (gōngwùyuán, fonctionnaire — HSK 3), \zh{薪水} (xīnshuǐ), \zh{监督制度} (jiāndū zhìdù, système de surveillance), \zh{严格} (yángé), \zh{政府} (zhèngfǔ), \zh{增加} (zēngjiā), \zh{加强} (jiāqiáng, renforcer), \zh{法律} (fǎlǜ), \zh{干净} (gānjìng, sens figuré « propre/honnête »), \zh{美好} (měihǎo), \zh{未来} (wèilái).
    \item \textbf{Grammaire / Connecteurs :}
    \begin{itemize}
        \item \zh{首先}... \zh{其次}... : Séquence logique claire, position initiale + virgule.
        \item \zh{为了解决这个问题，政府应该...} : Structure \cle{为了 + But, Sujet + 应该 + Verbe} parfaitement appliquée. Notez l'enchaînement de deux verbes \zh{增加... 加强...} partageant le même sujet \zh{政府}.
        \item \zh{不够严格} (bù gòu yángé) : Structure « pas assez + adj » (HSK 3), plus naturel que \zh{不严格} seul pour nuancer.
        \item \zh{更... 更...} (gèng... gèng...) : Comparatif de supériorité pour le vœu final.
    \end{itemize}
    \item \textbf{Points d'attention graphique (Caractères complexes) :}
    \begin{itemize}
        \item \zh{腐} : Vérifier le radical \zh{月} (pas \zh{⺼} mal formé) et la partie droite \zh{府} (le haut \zh{广} + \zh{付}).
        \item \zh{败} : Radical \zh{攵} (4 traits : point, horizontale, pliée, point) bien distinct à droite. Gauche \zh{貝} (yeux + traits).
        \item \zh{监} : Structure Haut-Bas. Le haut \zh{臥} s'écrit : \zh{人} + \zh{日} + \zh{彐} + \zh{人} (bas). Attention à ne pas confondre avec \zh{鉴} (métal \zh{金} en haut).
        \item \zh{督} : Radical \zh{目} (œil) à gauche, bien carré. Droite \zh{叔} : \zh{又} + \zh{小} (3 traits).
        \item \zh{严} : 17 traits. Ordre standard : \zh{山} (3) $\rightarrow$ \zh{彐} (3) $\rightarrow$ \zh{口} (3) $\rightarrow$ \zh{工} (3) $\rightarrow$ \zh{土} (3). Insister sur \textbf{Haut $\rightarrow$ Bas} et \textbf{Horizontal $\rightarrow$ Vertical} pour chaque composant.
        \item \zh{格} : Haut \zh{木} (4 traits), Bas \zh{各} (6 traits). Structure Haut-Bas. Ne pas écrire \zh{木} trop large.
    \end{itemize}
    \item \textbf{Ponctuation :} Utilisation exclusive de la ponctuation pleine chasse \zh{， 。 ！ ：}. Pas d'espaces entre caractères. Retour à la ligne physique sur la feuille à cases pour chaque phrase ou segment long.
\end{enumerate}
\end{exoresolu}

\courssub{Bilan de l'activité}
\noindent Cette activité d'intégration a permis de mettre en œuvre l'ensemble des compétences du Module 3 « Citoyenneté et ouverture au monde » dans une tâche finale authentique et motivante.

\begin{aretenir}[Bilan des acquis — Grille de compétences transversales]
\begin{center}
\begin{tabularx}{\linewidth}{|X|X|}
\hline
\rowcolor{poppurpleL} \textbf{Compétence / Savoir-faire} & \textbf{Preuve dans l'activité} \\ \hline
\textbf{Expression écrite} & Production d'un texte argumentatif structuré (6-8 phrases), cohérent, format lettre. \\ \hline
\textbf{Vocabulaire thématique} & Emploi correct de 10+ mots du thème (腐败, 贿赂, 薪水, 监督, 严格, 透明度...). \\ \hline
\textbf{Grammaire / Syntaxe} & Maîtrise des connecteurs \zh{首先/其次/为了/应该}, ordre des mots SOV/Thème-Rhème, structure \zh{不够+Adj}. \\ \hline
\textbf{Écriture des caractères} & Caractères complexes (腐, 败, 监, 督, 贿, 赂, 贪, 污, 严, 格) écrits dans le respect de l'ordre des traits, radicaux, proportions sur cases \emph{Tiangezhi}. \\ \hline
\textbf{Traduction (Thème/Version)} & Passage du français (brouillon mental) au chinois (texte final) et relecture via pinyin/traduction. \\ \hline
\textbf{Expression orale} & Lecture à voix haute devant la classe (prononciation, tons, débit). \\ \hline
\textbf{Compétence socioculturelle} & Contexte réel Cameroun-Chine (Yaoundé/Nankin), comparaison des enjeux de gouvernance, citoyenneté active. \\ \hline
\textbf{Autonomie / Évaluation} & Relecture par les pairs avec grille critériée, correction autonome, méta-cognition. \\ \hline
\end{tabularx}
\end{center}

\end{aretenir}

\begin{experience}[Prolongements possibles (Différenciation / Approfondissement)]
\begin{itemize}
    \item \textbf{Niveau A2+ (Renforcement) :} Écrire une version simplifiée (4-5 phrases) en remplaçant \zh{监督制度} par \zh{监督} seul, \zh{加强法律监督} par \zh{管得严一点} (guǎn de yán yīdiǎn - gérer plus strictement). Focus sur l'écriture des 5 caractères « stars » : \zh{腐败 监督 严格}.
    \item \textbf{Niveau B1 (Approfondissement) :} Ajouter un paragraphe sur la \textbf{Convention de l'Union Africaine de 2003} (签署公约 \zh{qiānyǔ gōngyuē}) ou la \textbf{Chine} (opération \zh{打虎拍蝇} dǎhǔ pāiyíng « frapper les tigres et écraser les mouches »). Recherche documentaire en chinois simplifié.
    \item \textbf{Numérique :} Saisir la lettre finale sur \textbf{OnBuch+} ou traitement de texte chinois (Sogou/Google Pinyin) pour valider la reconnaissance des caractères (pinyin $\rightarrow$ caractère), puis impression pour affichage.
    \item \textbf{Interdisciplinarité (EMC / Histoire-Géo) :} Étude de cas comparative : \zh{新加坡} (Singapour) / \zh{香港} (Hong Kong) / \zh{博茨瓦纳} (Botswana) — modèles de lutte anti-corruption réussis. Débat : « La corruption est-elle une fatalité culturelle ? » (Non, c'est un problème de système/institutions).
\end{itemize}
\end{experience}

\begin{savaistu}[Regard croisé Cameroun — Chine : La lutte anti-corruption]
\noindent \textbf{Cameroun :} Création de la \textbf{CONAC} (Commission Nationale Anti-Corruption) en 2006. Opération \textbf{« Épervier »} (depuis 2006) : arrestation de hauts fonctionnaires. Défis : indépendance de la justice, protection des lanceurs d'alerte, faibles salaires de base.
\par\medskip
\noindent \textbf{Chine :} Campagne \textbf{« Frapper les tigres et écraser les mouches »} (打虎拍蝇 \zh{dǎhǔ pāiyíng}) lancée par Xi Jinping en 2012. \zh{老虎} (lǎohǔ) = hauts dirigeants ; \zh{苍蝇} (cāngying) = petits fonctionnaires locaux. Création de la \textbf{Commission nationale de supervision} (国家监察委员会 \zh{Guójiā Jiānchá Wěiyuánhuì}) en 2018, fusionnant les fonctions de contrôle du Parti et de l'État. Résultat : des millions de fonctionnaires sanctionnés.
\par\medskip
\noindent \textbf{Point commun :} Volonté politique affichée au plus haut niveau. \textbf{Différence :} En Chine, le système est centralisé, le Parti contrôle l'appareil d'État (监察 \zh{jiānchá} + 纪律 \zh{jìlǜ}). Au Cameroun, séparation des pouvoirs constitutionnelle, mais défis d'effectivité judiciaire. \textbf{Leçon pour l'élève :} La corruption n'est pas une « culture », c'est un délit réprimé par la loi dans les deux pays. La jeunesse a un rôle de \zh{监督者} (superviseur/citoyen vigilant).
\end{savaistu}

\newpage

\coursec{Méthodes et exercices résolus}

\coursec{Leçon 30 — Expression écrite : facteurs de la corruption ; caractères complexes}

\courssub{Modèle de texte commenté}

\begin{textebox}[Texte modèle : Les causes de la corruption]
\zh{腐败是一个严重的社会问题。} \\
\zh{首先，有些人因为贫穷，所以想用不正当的办法得到钱。} \\
\zh{其次，法律不够完善，或者有人不遵守法律。} \\
\zh{最后，贪婪也是一个重要原因。} \\
\zh{因此，政府应该加强教育，严格执法，让每个人都做诚实的人。} \\

(Fǔbài shì yí ge yánzhòng de shèhuì wèntí. Shǒuxiān, yǒuxiē rén yīnwèi pínqióng, suǒyǐ xiǎng yòng bù zhèngdāng de bànfǎ dédào qián. Qícì, fǎlǜ bùgòu wánshàn, huòzhě yǒuxiē rén bù zūnshǒu fǎlǜ. Zuìhòu, tānlán yě shì yí ge zhòngyào yuányīn. Yīncǐ, zhèngfǔ yīnggāi jiāqiáng jiàoyù, yángé zhífǎ, ràng gè gèrén dōu zuò chéngshí de rén.)

« La corruption est un problème social grave. D'abord, certaines personnes, à cause de la pauvreté, veulent obtenir de l'argent par des moyens illégitimes. Ensuite, la loi n'est pas assez parfaite, ou bien certaines personnes ne respectent pas la loi. Enfin, la cupidité est aussi une cause importante. Par conséquent, le gouvernement devrait renforcer l'éducation, appliquer la loi strictement, pour que chacun devienne une personne honnête. »
\end{textebox}

\begin{propriete}[Analyse du texte modèle]
\begin{itemize}
\item \textbf{Structure globale} : Problème $\rightarrow$ Facteurs (3) $\rightarrow$ Conséquence/Solution. Adapté au niveau HSK 3 (phrases courtes, vocabulaire fréquent).
\item \textbf{Connecteurs logiques} : \zh{首先} (premier), \zh{其次} (second), \zh{最后} (dernier), \zh{因此} (par conséquent). Ils structurent l'argumentation.
\item \textbf{Structures grammaticales clés} :
  \begin{itemize}
  \item \zh{因为...所以...} (parce que... donc...) : lie cause et conséquence \emph{à l'intérieur de la phrase}.
  \item \zh{不够 + adjectif/verbe} (pas assez...) : \zh{不够完善} « pas assez parfaite ».
  \item \zh{应该 + verbe} (devrait) : modalité déontique pour la recommandation.
  \item \zh{让 + quelqu'un + verbe} (faire en sorte que / laisser) : construction causative.
  \end{itemize}
\item \textbf{Vocabulaire thème} : \zh{腐败} (corruption), \zh{贫穷} (pauvreté), \zh{不正当} (illégitime), \zh{完善} (parfait/complet), \zh{遵守} (respecter), \zh{贪婪} (cupidité), \zh{加强} (renforcer), \zh{执法} (appliquer la loi), \zh{诚实} (honnête).
\end{itemize}
\end{propriete}

\courssub{Structure et connecteurs utiles}

\begin{propriete}[Connecteurs logiques pour l'argumentation (niveau HSK 2-3)]
\begin{center}\begin{tabularx}{\linewidth}{|X|X|X|X|}\hline
\rowcolor{popblueL} Chinois & Pinyin & Français & Emploi / Position \\ \hline
\zh{首先} & shǒuxiān & d'abord, en premier & Début de phrase, suivi de \zh{，} \\ \hline
\zh{其次} & qícì & ensuite, en second & Début de phrase, suivi de \zh{，} \\ \hline
\zh{最后} & zuìhòu & enfin, en dernier & Début de phrase, suivi de \zh{，} \\ \hline
\zh{因为}... \zh{所以}... & yīnwèi... suǒyǐ... & parce que... donc... & \zh{因为} clause 1, \zh{所以} clause 2 (sujet souvent répété) \\ \hline
\zh{因此} & yīncǐ & par conséquent, c'est pourquoi & Début de phrase (conséquence globale) \\ \hline
\zh{但是} & dànshì & mais & Début de phrase ou après sujet (adversatif) \\ \hline
\zh{而且} & érqiě & de plus, en outre & Liaison deux verbes/adjectifs (progressif) \\ \hline
\end{tabularx}\end{center}
\end{propriete}

\begin{aretenir}[Structure type de paragraphe argumentatif (5-6 phrases)]
\zh{腐败是一个严重的问题。} (Énoncé du problème) \\
\zh{首先，[Facteur 1]，所以[Conséquence 1]。} \\
\zh{其次，[Facteur 2]。} \\
\zh{最后，[Facteur 3]。} \\
\zh{因此，我们/政府应该[Solution]。} (Conclusion / Proposition)

\textbf{Rappel ordre des mots (règle d'or)} : Sujet + \textbf{Temps} + \textbf{Lieu} + \textbf{Manière} + Verbe + Complément d'objet.
\emph{Ex.} \zh{他 昨天 在办公室 偷偷地 给了 官员 钱。}
\end{aretenir}

\courssub{Écriture des caractères complexes du thème}

\begin{methode}[Maîtriser les traits, radicaux et structure des caractères clés]
\begin{enumerate}
\item \textbf{Identifier le radical (clé)} : il indique le champ sémantique. Ex. \zh{肉} (viande/corps) pour \zh{腐} ; \zh{贝} (coquillage/valeur) pour \zh{败} ; \zh{讠} (parole) pour \zh{诚} ; \zh{钅} (métal) pour \zh{钱} ; \zh{氵} (eau) pour \zh{法}.
\item \textbf{Respecter l'ordre des traits universel} : haut $\rightarrow$ bas ; gauche $\rightarrow$ droite ; horizontal avant vertical ; extérieur avant intérieur ; fermer la boîte en dernier ; trait central avant ailes.
\item \textbf{Décomposer en composants fonctionnels} (phonétique / sémantique) pour mémoriser.
\item \textbf{S'entrainer} : écrire 5 fois chaque caractère en prononçant le pinyin (ton inclus) à voix haute.
\end{enumerate}
\end{methode}

\begin{exoresolu}[Analyse détaillée des caractères du thème]
Énoncé : Pour chaque caractère, donne le radical, le nombre de traits (approximatif), la structure (ex. haut-bas, gauche-droite) et la logique de composition.
\tcblower
\begin{center}\begin{tabularx}{\linewidth}{|X|X|X|X|X|}\hline
\rowcolor{popblueL} Caractère & Pinyin & Radical (Clé) & Structure & Logique / Décomposition \\ \hline
\zh{腐} & fǔ & \zh{肉} (bas, 6 tr.) & Haut-bas (\zh{广} + \zh{付} + \zh{肉}) & \zh{广} abri + \zh{付} (phonétique : donner/transmettre) + \zh{viande} $\rightarrow$ viande qui change/transmet $\rightarrow$ pourrir, corrompre. \\ \hline
\zh{败} & bài & \zh{贝} (gauche, 4 tr.) & Gauche-droite & \zh{贝} (coquillage/argent) + \zh{攵} (frapper/action) $\rightarrow$ perdre son argent/une bataille $\rightarrow$ échouer, corruption. \\ \hline
\zh{贪} & tān & \zh{贝} (bas, 4 tr.) & Haut-bas & \zh{今} (maintenant) au-dessus de \zh{贝} (argent) $\rightarrow$ vouloir l'argent maintenant $\rightarrow$ convoiter, être avide. \\ \hline
\zh{钱} & qián & \zh{钅} (gauche, 5 tr.) & Gauche-droite & \zh{钅} (métal/or) + \zh{戈} (hallebard/arme) $\rightarrow$ armes en métal = monnaie ancienne. \\ \hline
\zh{法} & fǎ & \zh{氵} (gauche, 3 tr.) & Gauche-droite & \zh{氵} (eau = norme, standard) + \zh{去} (phonétique + sens : aller/ôter) $\rightarrow$ standard qui écarte le mal $\rightarrow$ loi, méthode. \\ \hline
\zh{律} & lǜ & \zh{彳} (gauche, 3 tr.) & Gauche-droite & \zh{彳} (pas, marche) + \zh{聿} (pinceau/loi écrite) $\rightarrow$ marcher selon la règle $\rightarrow$ loi, discipline. \\ \hline
\zh{诚} & chéng & \zh{讠} (gauche, 2 tr.) & Gauche-droite & \zh{讠} (parole) + \zh{成} (accomplir, phonétique) $\rightarrow$ parole accomplie $\rightarrow$ sincère, honnête. \\ \hline
\end{tabularx}\end{center}
\end{exoresolu}

\begin{popfigure}
\begin{tikzpicture}[
  mindmap, grow cyclic, every node/.style=concept, concept color=popblue,
  level 1/.append style={level distance=3.5cm, sibling angle=90, concept color=poporange},
  level 2/.append style={level distance=2.5cm, sibling angle=45, concept color=popgreen},
  text=white, font=\small
]
\node {腐败因素}
  child { node {贫穷 \zh{(贫穷)}}
    child { node {收入低} }
    child { node {生活压力} } }
  child { node {法律缺失 \zh{(法律不全)}}
    child { node {监管弱} }
    child { node {执法不严} } }
  child { node {心理因素 \zh{(贪婪/侥幸)}}
    child { node {贪图钱财} }
    child { node {侥幸心理} } }
  child { node {教育不足 \zh{(缺乏诚信教育)}}
    child { node {道德缺失} }
    child { node {榜样缺乏} } };
\end{tikzpicture}
\legende{Carte mentale : les quatre piliers des facteurs de corruption (vocabulaire HSK 3-4).}
\end{popfigure}

\courssub{Caractères proches : ne pas confondre}

\begin{methode}[Stratégie de distinction des caractères proches]
\begin{enumerate}
\item \textbf{Isoler le trait distinctif} : quel trait ou composant change ? (Radical, trait ajouté/supprimé, position).
\item \textbf{Ancrer par le mot} : ne jamais mémoriser le caractère seul, toujours dans un mot fréquent (\zh{钱} dans \zh{钱包}, \zh{线} dans \zh{电线}).
\item \textbf{Produire une phrase minimale} contrastive pour chaque paire.
\end{enumerate}
\end{methode}

\begin{definition}[Paires à risque haute fréquence — Thème Corruption]
\begin{center}\begin{tabularx}{\linewidth}{|X|X|X|X|}\hline
\rowcolor{popblueL} Caractère cible & Pinyin & Mot ancrage (HSK) & Piège visuel (Caractère / Pinyin / Sens) \\ \hline
\zh{钱} & qián & \zh{钱包} qiánbāo (porte-monnaie) & \zh{线} xiàn (fil, ligne) — radical \zh{纟} (soie) vs \zh{钅} (métal) \\ \hline
\zh{诚} & chéng & \zh{诚实} chéngshí (honnête) & \zh{城} chéng (ville, cité) — radical \zh{讠} (parole) vs \zh{土} (terre) \\ \hline
\zh{法} & fǎ & \zh{法律} fǎlǜ (loi) & \zh{去} qù (aller) — \zh{法} a \zh{氵} (3 points eau) à gauche \\ \hline
\zh{败} & bài & \zh{腐败} fǔbài (corruption) & \zh{贩} fàn (vendre, colporter) — \zh{败} a \zh{攵} (frapper) à droite ; \zh{贩} a \zh{贩} (phonétique) \\ \hline
\zh{贿} & huì & \zh{贿赂} huìlù (pot-de-vin) & \zh{货} huò (marchandise) — \zh{贿} a \zh{穴} (caverne) en haut ; \zh{货} a \zh{化} \\ \hline
\zh{职} & zhí & \zh{职务} zhíwù (fonction, poste) & \zh{植} zhí (planter) — \zh{职} a \zh{耳} (oreille) en bas ; \zh{植} a \zh{木} (bois) \\ \hline
\end{tabularx}\end{center}
\end{definition}

\begin{exoresolu}[Entraînement : choisir le bon caractère]
Énoncé : Complète les phrases avec le bon caractère de la paire entre parenthèses.
\tcblower
(1) 我们必须遵守法\_\_。 (\zh{律} / \zh{去}) \\
$\rightarrow$ \zh{我们应该遵守法律。} (Wǒmen yīnggāi zūnshǒu \textbf{fǎlǜ}.) « Nous devons respecter la loi. » \emph{\zh{律} a le radical \zh{彳} (pas/marche), \zh{去} n'en a pas.}

(2) 他是一个\_\_实的人。 (\zh{诚} / \zh{城}) \\
$\rightarrow$ \zh{他是一个诚实的人。} (Tā shì yí ge \textbf{chéng}shí de rén.) « C'est une personne honnête. » \emph{\zh{诚} a le radical \zh{讠} (parole), \zh{城} a \zh{土} (terre).}

(3) 他想得到很多\_\_。 (\zh{钱} / \zh{线}) \\
$\rightarrow$ \zh{他想得到很多钱。} (Tā xiǎng dédào hěn duō \textbf{qián}.) « Il veut obtenir beaucoup d'argent. » \emph{\zh{钱} a \zh{钅} (métal), \zh{线} a \zh{纟} (soie).}

(4) 那个\_\_员收受了贿\_\_。 (\zh{职} / \zh{植} ; \zh{贿} / \zh{货}) \\
$\rightarrow$ \zh{那个职员收受了贿赂。} (Nà ge \textbf{zhí}yuán shōushòu le \textbf{huì}lù.) « Ce fonctionnaire a accepté des pots-de-vin. » \emph{\zh{职} (poste) vs \zh{植} (planter) ; \zh{贿} (pot-de-vin, caverne \zh{穴}) vs \zh{货} (marchandise).}
\end{exoresolu}

\courssub{Thème : français → chinois}

\begin{methode}[Procédure rigoureuse pour le thème (Français $\rightarrow$ Chinois)]
\begin{enumerate}
\item \textbf{Analyser la phrase française} : Sujet / Verbe / Compléments (Temps, Lieu, Manière, Objet).
\item \textbf{Réordonner selon la syntaxe chinoise} : Sujet + \textbf{Temps} + \textbf{Lieu} + \textbf{Manière} + Verbe + Objet.
\item \textbf{Sélectionner le vocabulaire exact} (registre, collocations : \zh{遵守法律} pas *\zh{尊重法律}).
\item \textbf{Insérer les mots-outils obligatoires} :
  \begin{itemize}
  \item Classificateur : \zh{一个问题}, \zh{一位官员} (\zh{位} respectueux).
  \item \zh{的} (épithète adjectif 2+ syllabes : \zh{严重的问题} ; possesseur : \zh{我的钱}).
  \item \zh{地} (adverbe manière : \zh{偷偷地给}).
  \item \zh{了} (aspect accompli / changement d'état : \zh{给了钱}).
  \item \zh{过} (expérience : \zh{去过北京}).
  \end{itemize}
\item \textbf{Vérifier} : tons (pinyin), caractères (proches), ponctuation chinoise (。 ，).
\end{enumerate}
\end{methode}

\begin{exoresolu}[Thème guidé — Corrigé commenté]
Énoncé : Traduis en chinois.
\tcblower
\textbf{(1) « La corruption est un problème social grave. »} \\
\zh{腐败是一个严重的社会问题。} (Fǔbài shì yí ge yánzhòng de shèhuì wèntí.) \\
\emph{Points :} \zh{是} (verbe copule) ; \zh{个} classif. ; \zh{严重} (2 syll.) + \zh{的} + \zh{社会问题}.

\textbf{(2) « Hier, le fonctionnaire a secrètement pris de l'argent. »} \\
\zh{昨天，那位官员偷偷地拿了钱。} (Zuótiān, nà wèi guānyuán tōutōu de ná le qián.) \\
\emph{Points :} Temps \zh{昨天} en tête ; \zh{位} classif. respectueux ; \zh{偷偷地} (manière) + \zh{地} ; \zh{拿了} (verbe + \zh{了} accompli).

\textbf{(3) « Nous devons éduquer les jeunes pour qu'ils deviennent honnêtes. »} \\
\zh{我们应该教育年轻人，让他们变得诚实。} (Wǒmen yīnggāi jiàoyù niánqīngrén, ràng tāmen biàn de chéngshí.) \\
\emph{Points :} \zh{应该} (devoir) ; \zh{让} (causatif) + \zh{变得} + adjectif ; pas de \zh{了} (obligation générale).

\textbf{(4) « Si tout le monde respecte la loi, la corruption diminuera. »} \\
\zh{如果大家都遵守法律，腐败就会减少。} (Rúguǒ dàjiā dōu zūnshǒu fǎlǜ, fǔbài jiù huì jiǎnshǎo.) \\
\emph{Points :} \zh{如果...就...} (si... alors) ; \zh{都} (tous) après sujet ; \zh{会} (futur/probabilité).
\end{exoresolu}

\begin{exercice}[Thème à faire — Devoir / Examen blanc]
Traduis les phrases suivantes en chinois (caractères + pinyin).
\begin{enumerate}
\item La cupidité est une cause importante de la corruption.
\item L'année dernière, il a donné de l'argent à un policier.
\item Le gouvernement doit renforcer la surveillance (监督 jiāndū).
\item Parce qu'il est pauvre, il veut obtenir de l'argent rapidement.
\item Nous devrions être des citoyens (公民 gōngmín) honnêtes.
\end{enumerate}
\end{exercice}

\begin{corrige}[Exercice Thème]
\begin{enumerate}
\item \zh{贪婪是腐败的一个重要原因。} (Tānlán shì fǔbài de yí ge zhòngyào yuányīn.)
\item \zh{去年，他给了一个警察一些钱。} (Qùnián, tā gěi le yí ge jǐngchá yìxiē qián.)
\item \zh{政府应该加强监督。} (Zhèngfǔ yīnggāi jiāqiáng jiāndū.)
\item \zh{因为他很穷，所以他想快点儿得到钱。} (Yīnwèi tā hěn qióng, suǒyǐ tā xiǎng kuài diǎnr dédào qián.) \emph{Note : \zh{快点儿} / \zh{很快地} acceptés.}
\item \zh{我们应该做诚实的公民。} (Wǒmen yīnggāi zuò chéngshí de gōngmín.)
\end{enumerate}
\end{corrige}

\courssub{Version : chinois → français}

\begin{methode}[Procédure rigoureuse pour la version (Chinois $\rightarrow$ Français)]
\begin{enumerate}
\item \textbf{Lecture globale} : identifier la structure (Sujet - Verbe - Objet), les connecteurs, les mots-outils (\zh{了, 过, 的, 地, 得}).
\item \textbf{Segmenter} : couper en unités de sens (propositions).
\item \textbf{Traduire les mots-outils/connecteurs en premier} :
  \begin{itemize}
  \item \zh{因为...所以...} $\rightarrow$ « Parce que... » (ou « Comme... ») — \textbf{jamais « donc » en français}.
  \item \zh{如果...就...} $\rightarrow$ « Si... alors... ».
  \item \zh{了} $\rightarrow$ Passé composé / aspect accompli.
  \item \zh{过} $\rightarrow$ « Avoir déjà... ».
  \item \zh{越来越} + Adj $\rightarrow$ « De plus en plus... ».
  \end{itemize}
\item \textbf{Rédiger en français naturel} : accorder les participes, choisir les articles, utiliser le subjonctif si nécessaire (\zh{应该} $\rightarrow$ « devrait que... » / « doit... »). Ne jamais calquer l'ordre chinois.
\end{enumerate}
\end{methode}

\begin{exoresolu}[Version guidée — Corrigé commenté]
Énoncé : Traduis en français naturel.
\tcblower
\textbf{(1) \zh{由于监督不力，贪污现象越来越严重。}} \\
(Yóuyú jiāndū bùlì, tānwū xiànxiàng yuè lái yuè yánzhòng.) \\
\emph{Analyse :} \zh{由于} (préposition « en raison de ») + \zh{监督不力} (surveillance inefficace) ; \zh{贪污现象} (phénomène de détournement) ; \zh{越来越严重} (de plus en plus grave). \\
\emph{Traduction :} « En raison d'une surveillance inefficace, le phénomène de corruption devient de plus en plus grave. »

\textbf{(2) \zh{如果官员不贪污，老百姓就会信任政府。}} \\
(Rúguǒ guānyuán bù tānwū, lǎobǎixìng jiù huì xìnrèn zhèngfǔ.) \\
\emph{Analyse :} \zh{如果...就...} (si... alors) ; \zh{老百姓} (le peuple, les gens ordinaires) ; \zh{信任} (faire confiance à). \\
\emph{Traduction :} « Si les fonctionnaires ne sont pas corrompus, le peuple fera confiance au gouvernement. »

\textbf{(3) \zh{他已经决定不再接受贿赂了。}} \\
(Tā yǐjīng juédìng bù zài jiēshòu huìlù le.) \\
\emph{Analyse :} \zh{已经...了} (déjà... action accomplie/état nouveau) ; \zh{决定} (décider) + \zh{不再} (ne plus) + Verbe. \\
\emph{Traduction :} « Il a déjà décidé de ne plus accepter de pots-de-vin. »
\end{exoresolu}

\begin{exercice}[Version à faire — Devoir / Examen blanc]
Traduis en français naturel.
\begin{enumerate}
\item \zh{首先，我们要明白腐败的危害。}
\item \zh{因为法律不严，所以有些人敢贪污。}
\item \zh{政府已经采取了措施反腐败。}
\item \zh{做一个诚实的人，不贪钱，不受贿。}
\item \zh{只有加强教育，腐败才会减少。}
\end{enumerate}
\end{exercice}

\begin{corrige}[Exercice Version]
\begin{enumerate}
\item « D'abord, nous devons comprendre les dangers de la corruption. »
\item « Parce que la loi n'est pas stricte, certaines personnes osent se livrer à la corruption. » (Ou : « Comme la loi n'est pas stricte, certains osent... »)
\item « Le gouvernement a déjà pris des mesures pour lutter contre la corruption. »
\item « Être une personne honnête, ne pas convoiter l'argent, ne pas accepter de pots-de-vin. » (Style impératif / maxime).
\item « Ce n'est qu'en renforçant l'éducation que la corruption diminuera. » (Structure \zh{只有...才...} = « Ce n'est que si... que... »).
\end{enumerate}
\end{corrige}

\courssub{Production guidée et grille d'évaluation}

\begin{experience}[Production orale \& écrite : Débat structuré « Lutte contre la corruption »]
\textbf{Contexte} : Vous êtes délégués de classe au Lycée de Yaoundé / Douala / Buea. Vous préparez une affiche ou un court discours (1 min) pour la « Journée de l'intégrité ».
\textbf{Consignes} :
\begin{enumerate}
\item \textbf{Par groupes de 3} : Rédigez un paragraphe de 6 phrases (chinois + pinyin) utilisant \textbf{au moins 4 connecteurs} (\zh{首先, 其次, 最后, 因为, 所以, 因此, 但是}) et \textbf{5 mots du vocabulaire thème} (\zh{腐败, 贪婪, 法律, 诚实, 监督, 教育, 贫穷, 贿赂}).
\item \textbf{Écriture} : Chaque membre écrit 2 phrases au tableau (caractères soignés, ordre des traits).
\item \textbf{Oral} : Un porte-parole lit le texte ; les autres classes notent sur la grille ci-dessous.
\end{enumerate}
\textbf{Modèle d'amorce} : \zh{同学们，腐败害人害己。首先……}
\end{experience}

\begin{aretenir}[Grille d'évaluation — Expression écrite (Sur 20 pts)]
\begin{center}\begin{tabularx}{\linewidth}{|X|X|X|}\hline
\rowcolor{popblueL} Critère & Indicateurs de réussite & Points \\ \hline
\textbf{Vocabulaire (5)} & Emploi correct de 5+ mots thème ; pas de confusion \zh{钱/线, 诚/城, 法/去} ; registres adaptés. & /5 \\ \hline
\textbf{Grammaire / Structures (5)} & Ordre Sujet-Temps-Lieu-Verbe respecté ; \zh{因为...所以...} correct ; \zh{的/地/得} distingués ; \zh{了/过} justifiés ; classif. présents. & /5 \\ \hline
\textbf{Caractères / Écriture (4)} & Caractères lisibles, traits dans l'ordre, radicaux corrects, pas de caractères inventés. & /4 \\ \hline
\textbf{Cohérence / Connecteurs (3)} & Plan logique (Problème - Causes - Solution) ; 3+ connecteurs utilisés à bon escient. & /3 \\ \hline
\textbf{Ponctuation / Pinyin (3)} & Ponctuation chinoise 。 ， ； ？ ; Pinyin avec tons corrects sur les voyelles (pas de chiffres). & /3 \\ \hline
\rowcolor{popblueL} \textbf{TOTAL} & & \textbf{/20} \\ \hline
\end{tabularx}\end{center}
\end{aretenir}

\begin{savaistu}[Regards croisés : Cameroun — Chine \#30]
\textbf{Cameroun} : La \zh{国家反腐败委员会} (CONAC, Commission Nationale Anti-Corruption) créée en 2006 mène des campagnes « \emph{École sans corruption} » dans les lycées. Le code pénal (Art. 134) sanctionne la corruption passive et active.
\textbf{Chine} : La \zh{中央纪委国家监委} (Commission centrale de contrôle de discipline / Commission nationale de supervision) pilote la campagne « \zh{打虎拍蝇} » (dǎ hǔ pāi yíng — « frapper les tigres, chasser les mouches ») : cibler à la fois les hauts fonctionnaires (\zh{老虎}) et la petite corruption de base (\zh{苍蝇}). Depuis 2012, des millions de cadres ont été sanctionnés.
\textbf{Point commun} : Dans les deux pays, l'\textbf{éducation civique} (\zh{廉政教育} liánzhèng jiàoyù) et la \textbf{transparence administrative} (\zh{阳光政务} yángguāng zhèngwù) sont présentées comme les piliers durables de la prévention, au-delà de la seule répression.
\end{savaistu}

\begin{attention}[Les 5 erreurs fatales au Bac / Examen de classe — Récapitulatif]
\begin{enumerate}
\item \textbf{Calque « Parce que... donc »} : Chinois \zh{因为...所以...} = Français « \textbf{Parce que}..., ... » \textbf{OU} « ..., \textbf{donc}... ». \emph{Jamais les deux.}
\item \textbf{Ordre Temps/Lieu} : \zh{我在学校学习} (Moi à l'école étudier) $\neq$ *\zh{我学习在学校}. \textbf{Toujours} avant le verbe.
\item \textbf{Classificateur absent} : *\zh{一问题} $\rightarrow$ \zh{一个问题} ; *\zh{两官员} $\rightarrow$ \zh{两位官员} / \zh{两个官员}.
\item \textbf{Caractères « jumeaux »} : \zh{钱/线, 诚/城, 法/去, 败/贩, 职/植, 贿/货}. Une erreur = sens changé = 0 pt sur la phrase.
\item \textbf{Abus de \zh{了}} : Pas de \zh{了} avec \zh{是} (présent), \zh{应该/想/能} (modaux), \zh{喜欢/爱} (états). \zh{了} = accompli / changement d'état \textbf{réel}.
\end{enumerate}
\end{attention}

\newpage

\coursec{Exercices}

\courssub{Je vérifie mes connaissances}

\begin{exercice}[QCM : Vocabulaire et structures du thème]
Chaque item ne comporte qu'une seule réponse correcte. Choisis la bonne proposition.
\begin{enumerate}
    \item Quel mot signifie « corruption » (au sens général, abus de pouvoir) ?
    \begin{itemize}
        \item A. 贪污 (tānwū)
        \item B. 腐败 (fǔbài)
        \item C. 贩卖 (fànmài)
        \item D. 犯罪 (fànzuì)
    \end{itemize}
    \item Pour dire « La loi est stricte », on utilise quel adjectif pour « stricte / sévère » ?
    \begin{itemize}
        \item A. 严厉 (yánlì)
        \item B. 严格 (yángé)
        \item C. 严重 (yánzhòng)
        \item D. 严肃 (yánsù)
    \end{itemize}
    \item Structure correcte pour exprimer la cause : « \emph{Parce que} la corruption est grave, \emph{donc} il faut la résoudre. »
    \begin{itemize}
        \item A. 因为……但是……
        \item B. 虽然……所以……
        \item C. 因为……所以……
        \item D. 不但……而且……
    \end{itemize}
    \item Le radical du caractère \zh{贪} (tān, cupide) et de \zh{贿} (huì, pot-de-vin) est :
    \begin{itemize}
        \item A. 贝 (bèi, coquillage/argent)
        \item B. 负 (fù, porter)
        \item C. 刂 (dāo, couteau)
        \item D. 心 (xīn, cœur)
    \end{itemize}
\end{enumerate}
\end{exercice}

\begin{exercice}[Vrai ou Faux justifié : Lecture rapide]
Lis le court texte ci-dessous. Détermine si les affirmations sont vraies (V) ou fausses (F) et justifie ta réponse en français en te basant sur le texte.

\begin{textebox}[Texte : Les causes de la corruption]
腐败是社会的大问题。首先，原因很复杂。有些人不诚实，想贪污钱财。其次，法律不够严格，监督不够。最后，教育缺乏，公民不懂得诚实的重要性。我们要解决腐败，必须加强法律，加强教育，培养诚实的公民。
\end{textebox}

\begin{enumerate}
    \item Le texte affirme que la corruption est un petit problème. \hfill V / F
    \item L'honnêteté des citoyens est citée comme une solution. \hfill V / F
    \item Le connecteur « 其次 » (qícì) introduit la première cause. \hfill V / F
    \item « 监督不够 » (jiāndū bùgòu) signifie « la surveillance est insuffisante ». \hfill V / F
\end{enumerate}
\end{exercice}

\begin{exercice}[Vocabulaire : Association et classification]
Associe chaque mot chinois à sa traduction française, puis classe-les dans le tableau ci-dessous selon leur nature grammaticale (Nom, Verbe, Adjectif, Adverbe/Connecteur).

\begin{center}
\begin{tabularx}{\linewidth}{|X|X|}\hline
\rowcolor{popblueL} \textbf{Chinois (Caractères + Pinyin)} & \textbf{Français} \\ \hline
1. 腐败 (fǔbài) & a. Honnête \\ \hline
2. 贪污 (tānwū) & b. Éducation \\ \hline
3. 法律 (fǎlǜ) & c. Corruption (général) \\ \hline
4. 严格 (yángé) & d. Cause / Raison \\ \hline
5. 诚实 (chéngshí) & e. Corruption (détournement de fonds) \\ \hline
6. 教育 (jiàoyù) & f. Loi \\ \hline
7. 原因 (yuányīn) & g. Résoudre \\ \hline
8. 解决 (jiějué) & h. Strict / Sévère \\ \hline
9. 社会 (shèhuì) & i. Société \\ \hline
10. 公民 (gōngmín) & j. Citoyen \\ \hline
\end{tabularx}
\end{center}

\begin{center}
\begin{tabularx}{\linewidth}{|X|X|X|X|}\hline
\rowcolor{popblueL} \textbf{Nom} & \textbf{Verbe} & \textbf{Adjectif / Verbe d'état} & \textbf{Connecteur / Autre} \\ \hline
 &  &  &  \\ \hline
 &  &  &  \\ \hline
 &  &  &  \\ \hline
 &  &  &  \\ \hline
\end{tabularx}
\end{center}
\end{exercice}

\begin{exercice}[Pinyin $\rightarrow$ Caractères : Dictée ciblée]
Écris les caractères chinois correspondants au pinyin suivant. Attention aux tons et aux radicaux (clés) : \zh{贝} (argent), \zh{氵} (eau/loi), \zh{言} (parole/honnêteté), \zh{彳} (marche/loi).
\begin{enumerate}
    \item fǔbài \hfill \rule{3cm}{0.4pt}
    \item tānwū \hfill \rule{3cm}{0.4pt}
    \item fǎlǜ \hfill \rule{3cm}{0.4pt}
    \item yángé \hfill \rule{3cm}{0.4pt}
    \item chéngshí \hfill \rule{3cm}{0.4pt}
    \item jiàoyù \hfill \rule{3cm}{0.4pt}
    \item yuányīn \hfill \rule{3cm}{0.4pt}
    \item jiějué \hfill \rule{3cm}{0.4pt}
    \item shèhuì \hfill \rule{3cm}{0.4pt}
    \item gōngmín \hfill \rule{3cm}{0.4pt}
\end{enumerate}
\end{exercice}

\courssub{Je m'entraîne}

\begin{exercice}[Thème / Version : Phrases isolées]
Traduis les phrases suivantes. Soigne l'ordre des mots (Sujet + Temps + Lieu + Manière + Verbe + Complément) et l'emploi des connecteurs.

\begin{enumerate}
    \item La corruption est un fléau pour la société. \\
    \rule{\linewidth}{0.4pt}
    \item Parce que les lois ne sont pas assez strictes, certains osent détourner de l'argent. \\
    \rule{\linewidth}{0.4pt}
    \item L'éducation est la méthode fondamentale pour résoudre la corruption. \\
    \rule{\linewidth}{0.4pt}
    \item Non seulement il est honnête, mais en plus il respecte strictement la loi. \\
    \rule{\linewidth}{0.4pt}
    \item 首先，我们要加强法律监督。其次，要加强诚信教育。 \\
    \rule{\linewidth}{0.4pt}
    \item 只有公民都诚实守法，社会才会更美好。 \\
    \rule{\linewidth}{0.4pt}
\end{enumerate}
\end{exercice}

\begin{exercice}[Transformation de phrases : Connecteurs et structures]
Transforme les phrases selon le modèle indiqué. Écris la phrase complète en caractères + pinyin.

\begin{enumerate}
    \item \textbf{Modèle \emph{因为……所以……} $\rightarrow$ \emph{由于……因此……}} \\
    因为腐败严重，所以我们要反腐败。 \\
    $\rightarrow$ \rule{\linewidth}{0.4pt}
    \item \textbf{Modèle Phrase simple $\rightarrow$ \emph{不但……而且……}} \\
    法律严格。教育重要。 \\
    $\rightarrow$ \rule{\linewidth}{0.4pt}
    \item \textbf{Modèle Affirmatif $\rightarrow$ \emph{只有……才……}} (Seulement si... alors...) \\
    公民诚实，社会才干净。 \\
    $\rightarrow$ \rule{\linewidth}{0.4pt}
    \item \textbf{Modèle Cause $\rightarrow$ \emph{之所以……是因为……}} \\
    社会腐败，因为缺乏监督。 \\
    $\rightarrow$ \rule{\linewidth}{0.4pt}
\end{enumerate}
\end{exercice}

\begin{exercice}[Texte à trous : Mots-outils et caractères clés]
Complète le texte ci-dessous avec les éléments manquants (caractères, mots-outils \zh{的/得/地}, \zh{了/过}, connecteurs). Le pinyin est donné en indice.

\begin{textebox}[Texte à compléter : Lutte contre la corruption]
反腐败\_\_\_\_\_\_ (shì) 一项长期任务。\_\_\_\_\_\_ (Shǒuxiān)，法律必须\_\_\_\_\_\_ (yángé)；\_\_\_\_\_\_ (Qícì)，监督要\_\_\_\_\_\_ (qiánghuà)；\_\_\_\_\_\_ (Zuìhòu)，教育\_\_\_\_\_\_ (bìxū) 跟上。只有\_\_\_\_\_\_ (zhǐyǒu) 全社会\_\_\_\_\_\_ (gòngtóng) 努力，腐败现象\_\_\_\_\_\_ (cái) 会减少。每个公民\_\_\_\_\_\_ (dōu) 应该\_\_\_\_\_\_ (chéngshí) 守法，\_\_\_\_\_\_ (bù) 參與\_\_\_\_\_\_ (tānwū) 活动。
\end{textebox}

\textit{Mots-outils à placer : 的 / 得 / 地 / 了 / 过 / 并 / 才 / 都 / 不 / 是 / 要 / 共同 / 必须 / 因为 / 所以}
\end{exercice}

\begin{exercice}[Remise en ordre et recopie : Structure d'un paragraphe argumentatif]
Remets les phrases suivantes dans l'ordre logique pour former un paragraphe cohérent (Introduction $\rightarrow$ Causes $\rightarrow$ Solutions $\rightarrow$ Conclusion). Recopie ensuite le texte complet en caractères chinois sur ton cahier.

\begin{enumerate}
    \item 最后，加强教育，培养诚实的公民，是根本出路。
    \item 腐败危害社会，必须坚决反腐败。
    \item 首先，原因在于制度不完善，监督不够严格。
    \item 其次，个人道德缺失，贪图钱财，也是重要原因。
    \item 因此，我们要完善法律，加大惩罚力度。
\end{enumerate}
\end{exercice}

\courssub{Je me prépare au Bac}

\begin{exercice}[Compréhension de l'écrit : Texte original type Bac]
Lis le texte suivant attentivement. Réponds aux questions en français (sauf indication contraire).

\begin{textebox}[Texte : 反腐败，从我做起 — Discours d'un lycéen à Yaoundé]
尊敬的老师们，亲爱的同学们：
大家好！
今天，我想和大家谈谈“反腐败”。腐败像一把毒刀，插在社会的心脏上。在喀麦隆，也在中国，腐败都阻碍国家发展。为什么会有腐败？我认为主要有两个原因。一是制度不够完善，监督不够严格，给贪污留下了机会。二是部分人缺乏诚信教育，只想捞钱，不顾后果。
怎么办？首先，国家要健全法律，让贪官不敢贪、不能贪。其次，学校要加强德育，从小培养我们诚实守法的品德。最后，我们每个公民都要从自己做起，拒绝行贿，举报腐败。
同学们，青年强则国强。让我们携手共建廉洁社会！
谢谢大家。
\end{textebox}

\begin{enumerate}
    \item \textbf{Vocabulaire :} Donne le sens en français de : 阻碍 (zǔ'ài)、健全 (jiànquán)、廉洁 (liánjié)、携手 (xiéshǒu).
    \item \textbf{Grammaire :} Repère dans le texte deux structures \emph{不但……而且……} ou \emph{让/使 + Quelqu'un + Verbe} (causative) et recopie-les.
    \item \textbf{Compréhension globale :} Quel est le public visé par ce discours ? Quel en est le but principal ?
    \item \textbf{Analyse des arguments :} Selon l'orateur, quelles sont les deux causes principales de la corruption ? Quelles sont les trois solutions proposées (repère les connecteurs \emph{首先、其次、最后}) ?
    \item \textbf{Expression personnelle (30-50 caractères) :} En chinois, écris une phrase pour dire : « En tant qu'élève, je dois refuser la corruption. »
\end{enumerate}
\end{exercice}

\begin{exercice}[Thème et Version : Épreuve type Bac (Barème : 10 pts)]
\textbf{Consignes :} Traduction soignée. Respect de la structure chinoise (Thème) et du style français (Version). Caractères lisibles, pinyin non requis pour la réponse finale mais caractères obligatoires.

\textbf{A. Thème (Français $\rightarrow$ Chinois) — 5 phrases / 5 pts}
Traduis en chinois (caractères simplifiés).
\begin{enumerate}
    \item La corruption nuit gravement au développement du pays.
    \item Les causes sont complexes : lois imparfaites et manque d'éducation morale.
    \item Le gouvernement doit renforcer la surveillance et punir sévèrement les corrompus.
    \item L'école doit éduquer les élèves à être honnêtes dès le plus jeune âge.
    \item Chaque citoyen a la responsabilité de construire une société intègre.
\end{enumerate}

\textbf{B. Version (Chinois $\rightarrow$ Français) — 1 texte / 5 pts}
Traduis en français correct et naturel.
\begin{textebox}[Texte à traduire]
反腐败斗争不能只靠政府，更需要全民参与。因为腐败的根源在于权力无监督和人心贪婪。所以我们要建立严格的监督机制，同时加强诚信教育。只有人人诚实守法，社会才能公平正义。这是我们这一代青年的责任和使命。
\end{textebox}
\end{exercice}

\begin{exercice}[Expression écrite guidée : Sujet de rédaction type Bac]
\textbf{Sujet :} « \emph{Comment être un citoyen honnête et participer à la lutte contre la corruption ?} » (Comment être un citoyen honnête et participer à la lutte contre la corruption ?)

\textbf{Consignes :}
\begin{itemize}
    \item Rédige un paragraphe cohérent d'environ \textbf{80 à 100 caractères chinois} (ponctuation comprise).
    \item Utilise \textbf{au moins 4 connecteurs différents} parmi : \zh{首先、其次、最后、因此、因为……所以……、只有……才……、不但……而且……、虽然……但是……}。
    \item Intègre \textbf{au moins 5 mots du glossaire suivant} : \zh{腐败、贪污、法律、严格、诚实、教育、原因、解决、社会、公民、监督、责任}。
    \item Soigne l'écriture des caractères (ordre des traits, proportion) et la ponctuation chinoise (。 ， 、 ；)。
\end{itemize}

\begin{center}
\begin{tikzpicture}[node distance=1.2cm, every node/.style={align=center, font=\small}]
\node[draw=popblue, thick, rounded corners, fill=popblueL, minimum width=12cm] (title) {\textbf{GRILLE D'AUTO-ÉVALUATION (Bac Blanc)}};
\node[draw=poporange, thick, rounded corners, fill=poporangeL, minimum width=12cm, below=0.3cm of title, align=center] (criteria){
\begin{tabularx}{11.5cm}{|X|c|}\hline
\rowcolor{poporangeL} \textbf{Critères} & \textbf{Points} \\ \hline
Respect du nombre de caractères (80-100) & /2 \\ \hline
Utilisation correcte de 4 connecteurs minimum & /3 \\ \hline
Intégration correcte de 5 mots du glossaire & /3 \\ \hline
Grammaire / Ordre des mots / Structures (把、被、连词) & /4 \\ \hline
Vocabulaire précis / Caractères corrects (radicaux, traits) & /3 \\ \hline
Cohérence du discours / Argumentation & /3 \\ \hline
Ponctuation chinoise / Présentation & /2 \\ \hline
\rowcolor{poporangeL} \textbf{TOTAL} & \textbf{/20} \\ \hline
\end{tabularx}
};
\end{tikzpicture}
\end{center}
\end{exercice}

\newpage

\coursec{Corrigés des exercices}

\begin{corrige}[Exercice 1 — QCM : Vocabulaire et structures du thème]
\begin{enumerate}
    \item \textbf{B. 腐败 (fǔbài)}.
    \zh{腐败} est le terme général pour « corruption » (fléau social, pourriture du système). \zh{贪污} (tānwū) désigne spécifiquement le « détournement de fonds / malversation » (acte précis d'un fonctionnaire). \zh{贩卖} = trafic, \zh{犯罪} = crime.
    \item \textbf{B. 严格 (yángé)}.
    \zh{严格} = strict, rigoureux (règles, lois, exigences ; on dit \zh{严格的法律}, \zh{严格要求}). \zh{严厉} (yánlì) = sévère, dur (personne, ton, punition : \zh{严厉的批评}). \zh{严重} (yánzhòng) = grave (problème, conséquence : \zh{问题严重}). \zh{严肃} (yánsù) = solennel, sérieux (attitude, occasion).
    \item \textbf{C. 因为……所以…… (yīnwèi… suǒyǐ…)}.
    Structure cause-effet standard : \zh{因为} + cause, \zh{所以} + conséquence. \zh{虽然……所以……} est fautif (on ne peut pas lier concession et conséquence ainsi ; la concession exige \zh{虽然……但是……}), \zh{因为……但是……} mélange cause et opposition, \zh{不但……而且……} = progression (non seulement… mais encore).
    \item \textbf{A. 贝 (bèi, coquillage/argent)}.
    \zh{贪} (cupide) et \zh{贿} (pot-de-vin) ont tous deux le radical \zh{贝} (coquillage, monnaie antique), lié à l'argent et à la richesse — on le retrouve aussi dans \zh{财} (richesse), \zh{贵} (cher), \zh{贫} (pauvre). \zh{刂} (couteau) et \zh{心} (cœur) n'apparaissent pas dans ces deux caractères.
\end{enumerate}
\end{corrige}

\begin{corrige}[Exercice 2 — Vrai ou Faux justifié : Lecture rapide]
\begin{enumerate}
    \item \textbf{F (Faux)}. Le texte dit : \zh{腐败是社会的大问题} (Fǔbài shì shèhuì de dà wèntí, « La corruption est un \textbf{grand} problème de la société »), pas un petit problème.
    \item \textbf{V (Vrai)}. La dernière phrase propose : \zh{培养诚实的公民} (péiyǎng chéngshí de gōngmín, « cultiver des citoyens honnêtes ») comme solution.
    \item \textbf{F (Faux)}. \zh{首先} (shǒuxiān, « premièrement ») introduit la première cause. \zh{其次} (qícì, « deuxièmement ») introduit la \textbf{deuxième} cause.
    \item \textbf{V (Vrai)}. \zh{监督} (jiāndū) = surveillance, contrôle ; \zh{不够} (bùgòu) = insuffisant (se place \emph{après} l'adjectif : \zh{严格} $\rightarrow$ \zh{不够严格}). Donc « la surveillance est insuffisante ».
\end{enumerate}
\end{corrige}

\begin{corrige}[Exercice 3 — Vocabulaire : Association et classification]
\textbf{1. Association (Chinois $\rightarrow$ Français) :}
\begin{center}
\begin{tabularx}{\linewidth}{|X|X|}\hline
\rowcolor{popblueL} \textbf{Chinois} & \textbf{Français} \\ \hline
1. 腐败 (fǔbài) & c. Corruption (général) \\ \hline
2. 贪污 (tānwū) & e. Détournement de fonds / malversation \\ \hline
3. 法律 (fǎlǜ) & f. Loi \\ \hline
4. 严格 (yángé) & h. Strict / rigoureux \\ \hline
5. 诚实 (chéngshí) & a. Honnête \\ \hline
6. 教育 (jiàoyù) & b. Éducation \\ \hline
7. 原因 (yuányīn) & d. Cause / raison \\ \hline
8. 解决 (jiějué) & g. Résoudre \\ \hline
9. 社会 (shèhuì) & i. Société \\ \hline
10. 公民 (gōngmín) & j. Citoyen \\ \hline
\end{tabularx}
\end{center}

\textbf{2. Classification par nature grammaticale :}
\begin{center}
\begin{tabularx}{\linewidth}{|X|X|X|X|}\hline
\rowcolor{popblueL} \textbf{Nom} & \textbf{Verbe} & \textbf{Adjectif / Verbe d'état} & \textbf{Connecteur / Autre} \\ \hline
腐败 (fǔbài) & 贪污 (tānwū) & 严格 (yángé) & 因为……所以…… \\ \hline
法律 (fǎlǜ) & 解决 (jiějué) & 诚实 (chéngshí) & 首先 / 其次 \\ \hline
原因 (yuányīn) & 教育 (jiàoyù)* &  &  \\ \hline
社会 (shèhuì) &  &  &  \\ \hline
公民 (gōngmín) &  &  &  \\ \hline
\end{tabularx}
\end{center}
\emph{Note :} \zh{教育} peut fonctionner comme nom (\zh{教育很重要}, « l'éducation est importante ») ou comme verbe (\zh{教育孩子}, « éduquer les enfants »). \zh{腐败} est nom (\zh{反对腐败}, « lutter contre la corruption ») ou verbe d'état (\zh{这个官员很腐败}, « ce fonctionnaire est corrompu »). On les classe ici selon l'usage le plus fréquent du thème.
\end{corrige}

\begin{corrige}[Exercice 4 — Pinyin $\rightarrow$ Caractères : Dictée ciblée]
\begin{enumerate}
    \item \zh{腐败} (fǔbài) — Radicaux : \zh{月} (chair/viande, en bas de \zh{腐}, sous la phonétique \zh{府} fǔ) ; \zh{贝} (argent, à gauche de \zh{败}). Ordre des traits : pour \zh{腐}, on écrit d'abord l'extérieur \zh{广} puis l'intérieur ; pour \zh{败}, gauche (\zh{贝}) avant droite (\zh{攵}).
    \item \zh{贪污} (tānwū) — Radical \zh{贝} (argent) pour les deux : \zh{贪} = \zh{今} (jīn, phonétique) au-dessus de \zh{贝} ; \zh{污} = \zh{氵} (eau) + \zh{亏}. \emph{Attention :} le radical de \zh{污} est \zh{氵} (eau souillée), celui de \zh{贪} est \zh{贝}.
    \item \zh{法律} (fǎlǜ) — \zh{氵} (eau) pour \zh{法} (origine : l'eau qui coule droit = la loi) ; \zh{彳} (pas, marche) pour \zh{律}.
    \item \zh{严格} (yángé) — \zh{严} : radical \zh{一} (caractère à écrire de haut en bas, la barre horizontale d'abord) ; \zh{木} (bois) pour \zh{格} (\zh{木} + \zh{各} gè, phonétique).
    \item \zh{诚实} (chéngshí) — \zh{讠} (parole, forme simplifiée de \zh{言}) pour \zh{诚} ; \zh{宀} (toit) pour \zh{实} (forme simplifiée de 實).
    \item \zh{教育} (jiàoyù) — \zh{攵} (action, taper) pour \zh{教} ; \zh{月} (chair) pour \zh{育} (en bas, sous \zh{云}).
    \item \zh{原因} (yuányīn) — \zh{厂} (falaise) pour \zh{原} ; \zh{囗} (enceinte) pour \zh{因} : on écrit le cadre extérieur, puis \zh{大} à l'intérieur, puis on ferme le cadre (règle : « fermer en dernier »).
    \item \zh{解决} (jiějué) — \zh{角} (corne) pour \zh{解} ; \zh{冫} (glace, deux points) pour \zh{决}. \emph{Piège :} \zh{决} s'écrit avec \zh{冫} (deux points), pas avec \zh{氵} (trois points) !
    \item \zh{社会} (shèhuì) — \zh{礻} (autel, forme de \zh{示}) pour \zh{社} : un point, un crochet horizontal, un vertical, un point — ne pas confondre avec \zh{衤} (vêtement, deux points) ; \zh{人} (homme) pour \zh{会} (forme simplifiée de 會).
    \item \zh{公民} (gōngmín) — \zh{八} pour \zh{公} ; \zh{氏} (clan) pour \zh{民}. \emph{Piège :} \zh{民} n'a pas le radical \zh{文} ; ne pas ajouter de trait en haut à droite (sinon on écrit un caractère fautif proche de \zh{氓}).
\end{enumerate}
\end{corrige}

\begin{corrige}[Exercice 5 — Thème / Version : Phrases isolées]
\begin{enumerate}
    \item \zh{腐败危害社会。} (Fǔbài wéihài shèhuì.) — \emph{La corruption nuit à la société.} (\zh{危害} = nuire à, être un fléau pour ; verbe transitif direct, sans préposition).
    \item \zh{因为法律不够严格，所以有些人敢贪污钱财。} (Yīnwèi fǎlǜ bùgòu yángé, suǒyǐ yǒuxiē rén gǎn tānwū qiáncái.) — Structure \zh{因为……所以……} ; \zh{不够} + adjectif ; \zh{敢} (oser) + verbe ; \zh{有些} + nom = « certains ».
    \item \zh{教育是解决腐败的根本方法。} (Jiàoyù shì jiějué fǔbài de gēnběn fāngfǎ.) — \zh{根本} (fondamental) ; \zh{方法} (méthode). Le \zh{的} lie le long attribut verbal \zh{解决腐败} au nom \zh{方法} : attribut + \zh{的} + nom.
    \item \zh{他不但诚实，而且严格遵守法律。} (Tā bùdàn chéngshí, érqiě yángé zūnshǒu fǎlǜ.) — Structure \zh{不但……而且……} ; \zh{遵守} (respecter, observer) + \zh{法律} ; \zh{严格} sert ici d'adverbe de manière devant le verbe.
    \item \emph{Premièrement, nous devons renforcer la surveillance légale. Deuxièmement, il faut renforcer l'éducation à l'honnêteté.} — \zh{首先，我们要加强法律监督。其次，要加强诚信教育。} (Shǒuxiān, wǒmen yào jiāqiáng fǎlǜ jiāndū. Qícì, yào jiāqiáng chéngxìn jiàoyù.) — \zh{加强} (renforcer) + nom ; \zh{诚信} (intégrité, honnêteté).
    \item \emph{Seulement si les citoyens sont tous honnêtes et respectent la loi, la société sera plus belle.} — \zh{只有公民都诚实守法，社会才会更美好。} (Zhǐyǒu gōngmín dōu chéngshí shǒufǎ, shèhuì cái huì gèng měihǎo.) — Structure \zh{只有……才……} (condition nécessaire) ; \zh{诚实守法} (honnête et respectueux de la loi, expression coordonnée) ; \zh{更} + adjectif = « plus ».
\end{enumerate}
\end{corrige}

\begin{corrige}[Exercice 6 — Transformation de phrases : Connecteurs et structures]
\begin{enumerate}
    \item \zh{由于腐败严重，因此我们要坚决反对腐败。} (Yóuyú fǔbài yánzhòng, yīncǐ wǒmen yào jiānjué fǎnduì fǔbài.) — \zh{由于} (préposition de cause, soutenu) + proposition ; \zh{因此} (adverbe de conséquence) en tête de la seconde proposition. \zh{反对} (s'opposer à, combattre) est le verbe correct ; on évite de dire \zh{反腐败} seul sans verbe.
    \item \zh{我们不但要完善法律，而且要重视教育。} (Wǒmen bùdàn yào wánshàn fǎlǜ, érqiě yào zhòngshì jiàoyù.) — \zh{不但……而且……} lie deux propositions de \emph{même sujet} : le sujet \zh{我们} se place \emph{avant} \zh{不但} et n'est pas répété. \zh{重视} (accorder de l'importance à).
    \item \zh{只有公民诚实，社会才干净。} (Zhǐyǒu gōngmín chéngshí, shèhuì cái gānjìng.) — \zh{只有……才……} : condition nécessaire. \zh{干净} (propre, sain) au sens figuré = « intègre ».
    \item \zh{社会之所以腐败，是因为缺乏监督。} (Shèhuì zhīsuǒyǐ fǔbài, shì yīnwèi quēfá jiāndū.) — Structure \zh{之所以……是因为……} (« la raison pour laquelle… c'est parce que… ») : le résultat d'abord, la cause ensuite ; mise en relief de la cause. \zh{缺乏} (manquer de) + objet. \emph{Attention au pinyin :} \zh{之所以} se lit zhīsuǒyǐ, à ne pas confondre avec \zh{只要} (zhǐyào, « pourvu que »).
\end{enumerate}
\end{corrige}

\begin{corrige}[Exercice 7 — Texte à trous : Mots-outils et caractères clés]
\textbf{Texte complété :}
\begin{textebox}[Texte complété : Lutte contre la corruption]
反腐败是一项长期任务。首先，法律必须严格；其次，监督要加强；最后，教育必须跟上。只有共同努力，腐败现象才会减少。每个公民都应该诚实守法，不参与贪污活动。\\
Fǎn fǔbài shì yí xiàng chángqī rènwu. Shǒuxiān, fǎlǜ bìxū yángé ; qícì, jiāndū yào jiāqiáng ; zuìhòu, jiàoyù bìxū gēnshàng. Zhǐyǒu gòngtóng nǔlì, fǔbài xiànxiàng cái huì jiǎnshǎo. Měige gōngmín dōu yīnggāi chéngshí shǒufǎ, bù cānyù tānwū huódòng.\\
« Lutter contre la corruption est une tâche de longue haleine. Premièrement, les lois doivent être strictes ; deuxièmement, la surveillance doit être renforcée ; enfin, l'éducation doit suivre. Ce n'est qu'en œuvrant ensemble que les phénomènes de corruption diminueront. Chaque citoyen doit être honnête et respecter la loi, et ne pas participer à des activités de détournement. »
\end{textebox}

\textbf{Justifications :}
\begin{itemize}
    \item \zh{是} (verbe « être ») : « La lutte anti-corruption \textbf{est} une tâche… » ; classificateur \zh{项} (xiàng) pour \zh{任务} (tâche, projet).
    \item \zh{首先 / 其次 / 最后} : connecteurs d'énumération (premièrement / deuxièmement / enfin).
    \item \zh{必须严格 / 必须跟上} : \zh{必须} (devoir, obligation forte) + verbe ou adjectif ; \zh{跟上} (suivre, être à la hauteur).
    \item \zh{加强} (jiāqiáng, renforcer) : verbe qui forme la collocation standard avec \zh{监督} : \zh{加强监督} (renforcer la surveillance).
    \item \zh{只有……才……} : structure de condition nécessaire ; \zh{共同} (ensemble, en commun) modifie \zh{努力}.
    \item \zh{都} : avec \zh{每个}, « chaque… tous » (quantification universelle : \zh{每个公民都…}).
    \item \zh{诚实守法} : expression coordonnée figée (« honnête et respectueux de la loi »), sans \zh{地} entre les deux verbes d'état.
    \item \zh{不} : négation devant \zh{参与} (cānyù, participer). \emph{Attention à l'écriture :} \zh{参与} en caractères simplifiés (la forme 參與 est traditionnelle, à ne pas utiliser).
    \item Mots-outils \zh{的/得/地/了/过} non requis ici : pas de complément du nom à lier, pas de complément de degré, pas d'aspect accompli ou expérientiel sur le verbe principal.
\end{itemize}
\end{corrige}

\begin{corrige}[Exercice 8 — Remise en ordre et recopie : Structure d'un paragraphe argumentatif]
\textbf{Ordre logique :} 2 $\rightarrow$ 3 $\rightarrow$ 4 $\rightarrow$ 5 $\rightarrow$ 1

\textbf{Texte complet recopié :}
\begin{textebox}[Paragraphe argumentatif ordonné]
腐败危害社会，必须坚决反对腐败。首先，原因在于制度不完善，监督不够严格。其次，个人道德缺失，贪图钱财，也是重要原因。因此，我们要完善法律，加大惩罚力度。最后，加强教育，培养诚实的公民，是根本出路。\\
Fǔbài wéihài shèhuì, bìxū jiānjué fǎnduì fǔbài. Shǒuxiān, yuányīn zàiyú zhìdù bù wánshàn, jiāndū bùgòu yángé. Qícì, gèrén dàodé quēshī, tāntú qiáncái, yě shì zhòngyào yuányīn. Yīncǐ, wǒmen yào wánshàn fǎlǜ, jiādà chéngfá lìdù. Zuìhòu, jiāqiáng jiàoyù, péiyǎng chéngshí de gōngmín, shì gēnběn chūlù.\\
« La corruption nuit à la société, il faut la combattre résolument. Premièrement, la cause réside dans l'imperfection du système et l'insuffisance de la surveillance. Deuxièmement, le manque de morale individuelle et la soif d'argent sont aussi des causes importantes. C'est pourquoi nous devons perfectionner les lois et renforcer les sanctions. Enfin, renforcer l'éducation et former des citoyens honnêtes est la voie fondamentale. »
\end{textebox}

\textbf{Analyse structurelle :}
\begin{itemize}
    \item \textbf{Introduction (thèse) :} phrase 2 — constat + prise de position.
    \item \textbf{Argument 1 (cause structurelle) :} phrase 3 — \zh{首先} + \zh{制度不完善、监督不够严格}. Structure utile : \zh{原因在于…} (« la cause réside dans… »).
    \item \textbf{Argument 2 (cause individuelle) :} phrase 4 — \zh{其次} + \zh{道德缺失、贪图钱财} ; \zh{也是} marque l'ajout (« est aussi »).
    \item \textbf{Solution immédiate (conséquence) :} phrase 5 — \zh{因此} + \zh{完善法律、加大惩罚力度} (« perfectionner les lois, accroître l'intensité des sanctions »).
    \item \textbf{Conclusion (solution fondamentale) :} phrase 1 — \zh{最后} + \zh{加强教育、培养诚实的公民、根本出路} (« voie de sortie fondamentale »).
\end{itemize}
\end{corrige}

\begin{corrige}[Exercice 9 — Compréhension de l'écrit : Texte original type Bac]
\begin{enumerate}
    \item \textbf{Vocabulaire :}
    \begin{itemize}
        \item \zh{阻碍} (zǔ'ài) : entraver, gêner, faire obstacle à.
        \item \zh{健全} (jiànquán) : perfectionner, rendre complet et sain (système, loi, organisation) ; ici verbe : \zh{健全法律}.
        \item \zh{廉洁} (liánjié) : intègre, honnête, incorruptible (qualité morale, administration) ; \zh{廉洁社会} = « société intègre ».
        \item \zh{携手} (xiéshǒu) : main dans la main, unir ses efforts, collaborer étroitement.
    \end{itemize}
    \item \textbf{Grammaire — Structures repérées :}
    \begin{enumerate}
        \item \textbf{Structure causative \zh{让/使 + quelqu'un + verbe} :} \zh{让贪官不敢贪、不能贪} (ràng tānguān bù gǎn tān, bù néng tān) — « faire en sorte que les fonctionnaires corrompus n'osent pas et ne puissent pas se corrompre ». \zh{让} = faire, amener à ; le sujet de la seconde proposition (\zh{贪官}) est à la fois objet de \zh{让} et sujet du verbe suivant (construction pivot).
        \item \textbf{Progression :} le texte n'utilise pas explicitement \zh{不但……而且……}, mais la phrase \zh{在喀麦隆，也在中国，腐败都阻碍国家发展} (Zài Kāmàilóng, yě zài Zhōngguó, fǔbài dōu zǔ'ài guójiā fāzhǎn) exprime une progression de portée (« non seulement au Cameroun, mais aussi en Chine ») grâce à \zh{也……都……}. Pour exprimer la même idée avec la structure étudiée : \zh{腐败不但阻碍喀麦隆的发展，而且阻碍中国的发展。}
    \end{enumerate}
    \item \textbf{Compréhension globale :}
    \begin{itemize}
        \item \textbf{Public visé :} \zh{老师们，同学们} (professeurs et camarades élèves) $\rightarrow$ un auditoire scolaire (lycéens), lors d'un discours ou d'une allocution.
        \item \textbf{But principal :} sensibiliser la jeunesse à la lutte contre la corruption (\zh{反腐败，从我做起}, « lutter contre la corruption, en commençant par soi-même ») et appeler à l'action collective (\zh{携手共建廉洁社会}, « construire main dans la main une société intègre »).
    \end{itemize}
    \item \textbf{Analyse des arguments :}
    \begin{itemize}
        \item \textbf{Deux causes principales :}
              1. \zh{制度不够完善，监督不够严格} (système imparfait, surveillance insuffisante) — cause structurelle, institutionnelle.
              2. \zh{部分人缺乏诚信教育，只想捞钱，不顾后果} (une partie des gens manque d'éducation à l'intégrité, ne pense qu'à s'enrichir, sans se soucier des conséquences) — cause morale, individuelle.
        \item \textbf{Trois solutions (avec connecteurs) :}
              1. \zh{首先} : \zh{国家要健全法律，让贪官不敢贪、不能贪} (perfectionner la loi, dissuader).
              2. \zh{其次} : \zh{学校要加强德育，从小培养我们诚实守法的品德} (renforcer l'éducation morale à l'école, dès le plus jeune âge).
              3. \zh{最后} : \zh{我们每个公民都要从自己做起，拒绝行贿，举报腐败} (responsabilité citoyenne : refuser de verser des pots-de-vin, dénoncer la corruption).
    \end{itemize}
    \item \textbf{Expression personnelle (30-50 caractères) :}
    \zh{作为学生，我必须拒绝腐败，诚实守法，还要劝身边的人不参与贪污。} (Zuòwéi xuéshēng, wǒ bìxū jùjué fǔbài, chéngshí shǒufǎ, hái yào quàn shēnbiān de rén bù cānyù tānwū.) — 33 caractères (ponctuation comprise). Conforme.
    \emph{Variantes possibles :} \zh{做诚实学生，拒绝贪污，从我做起。} / \zh{青年当自强，拒绝腐败，共建廉洁社会。}
\end{enumerate}
\end{corrige}

\begin{corrige}[Exercice 10 — Thème et Version : Épreuve type Bac (Barème : 10 pts)]
\textbf{A. Thème (Français $\rightarrow$ Chinois) — 5 pts (1 pt/phrase)}
\begin{enumerate}
    \item \zh{腐败严重危害国家发展。} (Fǔbài yánzhòng wéihài guójiā fāzhǎn.) — \zh{严重} (gravement) se place \emph{devant} le verbe ; \zh{危害} (nuire à) est transitif direct.
    \item \zh{原因复杂：法律不完善，缺乏德育。} (Yuányīn fùzá : fǎlǜ bù wánshàn, quēfá déyù.) — \zh{不完善} (imparfait) ; \zh{德育} (éducation morale), terme standard ; \zh{缺乏} + nom (manquer de).
    \item \zh{政府必须加强监督，严惩贪腐分子。} (Zhèngfǔ bìxū jiāqiáng jiāndū, yánchéng tānfǔ fènzǐ.) — \zh{加强监督} (renforcer la surveillance) ; \zh{严惩} (punir sévèrement) ; \zh{贪腐分子} (éléments corrompus). \emph{Pinyin :} \zh{贪腐} se lit tānfǔ.
    \item \zh{学校应从小教育学生诚实守信。} (Xuéxiào yīng cóngxiǎo jiàoyù xuéshēng chéngshí shǒuxìn.) — \zh{从小} (dès le plus jeune âge), complément de temps placé devant le verbe ; \zh{诚实守信} (honnête et digne de confiance).
    \item \zh{每个公民都有责任建设廉洁社会。} (Měige gōngmín dōu yǒu zérèn jiànshè liánjié shèhuì.) — \zh{每个……都……} (chaque…) ; \zh{有责任} + verbe (avoir le devoir de) ; \zh{建设} (construire) ; \zh{廉洁社会} (société intègre).
\end{enumerate}

\textbf{Barème indicatif Thème :} 1 pt par phrase. 0,5 si erreur de vocabulaire mineure ; 0 si erreur de structure majeure (ordre des mots, connecteur manquant) ou caractère faux.

\textbf{B. Version (Chinois $\rightarrow$ Français) — 5 pts}
\textbf{Traduction proposée :}
\begin{quote}
« La lutte contre la corruption ne peut pas reposer uniquement sur le gouvernement ; elle exige d'autant plus la participation de tous. Car la racine de la corruption réside dans l'absence de contrôle du pouvoir et dans la cupidité des cœurs. C'est pourquoi nous devons établir un mécanisme de surveillance strict, tout en renforçant l'éducation à l'intégrité. Ce n'est que si chacun est honnête et respecte la loi que la société pourra être juste et équitable. C'est là la responsabilité et la mission de notre génération de jeunes. »
\end{quote}

\textbf{Points de vigilance (Barème) :}
\begin{itemize}
    \item \zh{不能只靠……更需要……} : « ne pas pouvoir compter seulement sur… mais avoir encore plus besoin de… » (progression \zh{不只……更……}). — 1 pt
    \item \zh{根源在于……和……} : « la cause profonde réside dans… et… » (\zh{在于} = résider dans, verbe soutenu). — 1 pt
    \item \zh{权力无监督 / 人心贪婪} : « l'absence de surveillance du pouvoir / la cupidité des cœurs » ; \zh{无} = sans, absence de (registre soutenu pour \zh{没有}). — 1 pt
    \item \zh{建立……机制，同时加强……} : « établir un mécanisme…, tout en renforçant… » (\zh{同时} = en même temps, par ailleurs). — 1 pt
    \item \zh{只有……才能……} : « ce n'est que si… que… pourra… » (condition nécessaire) ; \zh{公平正义} (juste et équitable). — 1 pt
\end{itemize}
\end{corrige}

\begin{corrige}[Exercice 11 — Expression écrite guidée : Sujet de rédaction type Bac]
\textbf{Proposition de rédaction modèle (environ 95 caractères) :}
\begin{textebox}[Rédaction modèle]
作为公民，首先，我们要诚实守法，不参与贪污；其次，要监督权力，因为监督不力，所以容易滋生腐败；最后，要支持法律严惩腐败分子。只有人人尽责任，社会才会廉洁。\\
Zuòwéi gōngmín, shǒuxiān, wǒmen yào chéngshí shǒufǎ, bù cānyù tānwū ; qícì, yào jiāndū quánlì, yīnwèi jiāndū bù lì, suǒyǐ róngyì zīshēng fǔbài ; zuìhòu, yào zhīchí fǎlǜ yánchéng fǔbài fènzǐ. Zhǐyǒu rénrén jìn zérèn, shèhuì cái huì liánjié.\\
« En tant que citoyens, premièrement, nous devons être honnêtes et respecter la loi, sans participer aux détournements ; deuxièmement, nous devons surveiller le pouvoir, car si la surveillance fait défaut, la corruption se développe facilement ; enfin, nous devons soutenir la sévérité de la loi envers les corrompus. Ce n'est que si chacun remplit son devoir que la société sera intègre. »
\end{textebox}
\textbf{Comptage :} environ 95 caractères (ponctuation comprise). Conforme à la consigne (80-100).

\textbf{Vérification des consignes :}
\begin{itemize}
    \item \textbf{Connecteurs (5 utilisés) :} \zh{首先}、\zh{其次}、\zh{最后}、\zh{因为……所以……}、\zh{只有……才……}. (\textbf{Validé, $\ge$ 4}.)
    \item \textbf{Mots du glossaire (8 utilisés) :} \zh{公民}、\zh{诚实}、\zh{贪污}、\zh{腐败}、\zh{监督}、\zh{法律}、\zh{责任}、\zh{社会}. (\textbf{Validé, $\ge$ 5}.)
    \item \textbf{Grammaire et structures :} \zh{作为……} (préposition, « en tant que ») ; \zh{不} + verbe (négation) ; \zh{因为……所以……} ; \zh{只有……才……} ; \zh{要} + verbe (devoir/volonté). Ordre Sujet + Temps/Lieu + Manière + Verbe respecté.
    \item \textbf{Ponctuation :} \zh{；} pour séparer les propositions coordonnées longues, \zh{，} à l'intérieur des propositions, \zh{。} final.
\end{itemize}

\textbf{Grille d'auto-évaluation — Commentaire pour l'élève :}
\begin{itemize}
    \item \textbf{Respect du nombre de caractères (/2) :} environ 95 caractères $\rightarrow$ 2/2.
    \item \textbf{Connecteurs (/3) :} 5 connecteurs distincts utilisés correctement $\rightarrow$ 3/3.
    \item \textbf{Glossaire (/3) :} 8 mots intégrés naturellement $\rightarrow$ 3/3.
    \item \textbf{Grammaire (/4) :} structures complexes (\zh{因为……所以……}, \zh{只有……才……}, \zh{作为}) sans erreur d'ordre $\rightarrow$ 4/4.
    \item \textbf{Vocabulaire et caractères (/3) :} caractères complexes (\zh{监、督、贪、腐、责、任、廉}) écrits correctement, radicaux respectés $\rightarrow$ 3/3.
    \item \textbf{Cohérence (/3) :} argumentation logique (individuel $\rightarrow$ institutionnel $\rightarrow$ légal et collectif) $\rightarrow$ 3/3.
    \item \textbf{Ponctuation et présentation (/2) :} ponctuation chinoise standard, écriture soignée $\rightarrow$ 2/2.
    \item \textbf{TOTAL : 20/20.}
\end{itemize}

\textbf{Points de vigilance pour le Bac :}
\begin{enumerate}
    \item \textbf{Ne pas écrire en style « télégraphique »} (sujet + verbe seulement) : utilisez des compléments (lieu, manière, but) et des connecteurs.
    \item \textbf{Attention aux caractères « pièges » :} \zh{监} (radical \zh{皿}, récipient — pas \zh{血}) ; \zh{督} (radical \zh{目}, l'œil, en bas) ; \zh{贪} (\zh{今} au-dessus de \zh{贝} — ne pas écrire \zh{令} à la place de \zh{今}) ; \zh{腐} (\zh{府} sous \zh{广}, avec \zh{月} en bas) ; \zh{责} (radical \zh{贝}) ; \zh{任} (\zh{亻} + \zh{壬}) ; \zh{廉} (sous \zh{广}, ne pas oublier les traits intérieurs).
    \item \textbf{Caractères simplifiés uniquement :} écrire \zh{参与} (et non la forme traditionnelle 參與), \zh{会} (et non 會), \zh{实} (et non 實).
    \item \textbf{Connecteurs :} ne les juxtaposez pas sans logique. \zh{首先/其次/最后} pour l'énumération ; \zh{因为/所以} pour la cause interne ; \zh{因此} pour la conséquence ; \zh{只有/才} pour la condition nécessaire globale.
    \item \textbf{\zh{的/得/地} :} \zh{诚实守法} (pas de \zh{地} entre verbes d'état coordonnés) ; \zh{严格地监督} possible avec \zh{地} si l'adjectif devient adverbe de manière, mais la collocation \zh{加强监督} reste préférable ; \zh{诚实的公民} avec \zh{的} (attribut + nom).
\end{enumerate}
\end{corrige}

\newpage

\coursec{Fiche bilan}

\begin{popfigure}
\begin{tikzpicture}[mindmap, grow cyclic, every node/.style={concept, align=center, font=\small}, root concept/.append style={concept color=popblue, font=\bfseries}, level 1/.append style={level distance=3.2cm, sibling angle=51}, level 2/.append style={level distance=2.2cm}]
\node[root concept, align=center]{Leçon 30\\ La corruption\\ 腐败 fǔbài}
  child[concept color=poporange]{ node[align=center]{Lexique clé\\ 腐败 贪污\\ 受贿 法律} }
  child[concept color=popgreen]{ node[align=center]{Causes\\ 因为…所以…\\ 由于 yóuyú} }
  child[concept color=poppurple]{ node[align=center]{Conséquences\\ 导致 dǎozhì\\ 影响 yǐngxiǎng} }
  child[concept color=poppink]{ node[align=center]{Solutions\\ 应该 必须\\ 加强 打击} }
  child[concept color=popgold]{ node[align=center]{Connecteurs\\ 首先 其次\\ 最后 总之} }
  child[concept color=popblue!70]{ node[align=center]{Caractères\\ 腐 败 贪 污\\ 贿 律} }
  child[concept color=popgreen!70]{ node[align=center]{Traduction\\ thème et\\ version} };
\end{tikzpicture}
\legende{Carte mentale de la leçon : écrire un texte simple sur les facteurs de la corruption.}
\end{popfigure}

\begin{bilan}
\textbf{1. Lexique clé du thème}\par
\begin{center}
\begin{tabularx}{\linewidth}{|X|X|X|X|}
\hline
\rowcolor{popblueL} Caractères & Pinyin & Nature & Traduction \\
\hline
腐败 & fǔbài & nom / adj. & corruption ; corrompu \\
\hline
贪污 & tānwū & verbe & détourner (des fonds) \\
\hline
受贿 & shòuhuì & verbe & accepter des pots-de-vin \\
\hline
法律 & fǎlǜ & nom & loi, législation \\
\hline
原因 & yuányīn & nom & cause, raison \\
\hline
导致 & dǎozhì & verbe & entraîner, conduire à \\
\hline
影响 & yǐngxiǎng & nom / verbe & influence ; influencer \\
\hline
发展 & fāzhǎn & nom / verbe & développement ; se développer \\
\hline
加强 & jiāqiáng & verbe & renforcer \\
\hline
打击 & dǎjī & verbe & combattre, réprimer \\
\hline
诚实 & chéngshí & adj. & honnête \\
\hline
公民 & gōngmín & nom & citoyen \\
\hline
\end{tabularx}
\end{center}

\textbf{2. Grammaire à connaître par cœur}\par
\begin{center}
\begin{tabularx}{\linewidth}{|X|X|X|}
\hline
\rowcolor{popblueL} Structure & Exemple (caractères + pinyin) & Traduction \\
\hline
因为 A，所以 B & 因为法律不严，所以腐败很多。Yīnwèi fǎlǜ bù yán, suǒyǐ fǔbài hěn duō. & Parce que la loi n'est pas stricte, la corruption est répandue. \\
\hline
由于 + cause，+ conséquence & 由于贫穷，一些人贪污。Yóuyú pínqióng, yìxiē rén tānwū. & À cause de la pauvreté, certains détournent des fonds. \\
\hline
A 导致 B & 腐败导致国家不发展。Fǔbài dǎozhì guójiā bù fāzhǎn. & La corruption empêche le développement du pays. \\
\hline
应该 / 必须 + verbe & 我们应该加强教育。Wǒmen yīnggāi jiāqiáng jiàoyù. & Nous devons renforcer l'éducation. \\
\hline
为了 + but，+ action & 为了打击腐败，政府制定了法律。Wèile dǎjī fǔbài, zhèngfǔ zhìdìngle fǎlǜ. & Pour combattre la corruption, le gouvernement a établi des lois. \\
\hline
\end{tabularx}
\end{center}

\textbf{3. Connecteurs pour structurer le texte}\par
\begin{itemize}
  \item \cle{首先} shǒuxiān (d'abord) → \cle{其次} qícì (ensuite) → \cle{最后} zuìhòu (enfin) ;
  \item \cle{但是} dànshì (mais), \cle{所以} suǒyǐ (donc), \cle{而且} érqiě (de plus) ;
  \item conclusion : \cle{总之} zǒngzhī (en somme), \cle{我认为} wǒ rènwéi (je pense que).
\end{itemize}

\textbf{4. Écriture des caractères complexes}\par
\begin{itemize}
  \item \zh{腐} (14 traits) : clé 广 (abri) en haut-gauche, puis 府 + 肉 (viande) à l'intérieur ; écrire l'extérieur avant l'intérieur.
  \item \zh{败} (8 traits) : 贝 (coquillage, argent) à gauche + 攵 (frapper) à droite ; gauche avant droite.
  \item \zh{贪} (8 traits) : 今 en haut + 贝 en bas ; du haut vers le bas. Ne pas confondre avec \zh{贫} pín (pauvre) : 分 + 贝.
  \item \zh{贿} (10 traits) : 贝 à gauche + 有 à droite. Tous les caractères liés à l'argent portent la clé \cle{贝}.
  \item \zh{律} (9 traits) : 彳 (pas, marcher) à gauche + 聿 à droite ; horizontal avant vertical.
\end{itemize}

\textbf{5. Caractères proches : ne pas confondre}\par
\begin{itemize}
  \item \zh{贪} tān (convoiter) ≠ \zh{贫} pín (pauvre) ;
  \item \zh{败} bài (perdre, corrompre) ≠ \zh{贩} fàn (vendre) ;
  \item \zh{律} lǜ (loi) ≠ \zh{津} jīn (Tianjin, gué) ;
  \item \zh{原} yuán (origine) ≠ \zh{愿} yuàn (vouloir).
\end{itemize}

\textbf{6. Méthode express de production écrite}\par
\begin{enumerate}
  \item Introduction : définir le problème (\zh{腐败是一个严重的问题。}).
  \item Causes : 2 ou 3 facteurs avec 因为 / 由于.
  \item Conséquences : 导致 / 影响.
  \item Solutions : 应该 / 必须 / 为了.
  \item Conclusion : 总之 + avis personnel.
\end{enumerate}
\end{bilan}

\begin{aretenir}[Checklist avant le Bac]
\begin{itemize}
  \item $\square$ Je connais par cœur le lexique du thème : 腐败、贪污、受贿、法律、导致、影响、加强、打击.
  \item $\square$ Je sais écrire les caractères complexes 腐、败、贪、贿、律 dans le bon ordre des traits.
  \item $\square$ Je distingue les caractères proches : 贪/贫、败/贩、律/津、原/愿.
  \item $\square$ Je maîtrise la structure 因为…所以… sans oublier la virgule chinoise ，.
  \item $\square$ Je sais utiliser 由于 + cause et A 导致 B pour exprimer les conséquences.
  \item $\square$ J'emploie 应该 / 必须 / 为了 pour proposer des solutions.
  \item $\square$ Je structure mon texte avec 首先、其次、最后、总之.
  \item $\square$ Je place les compléments de temps et de lieu avant le verbe.
  \item $\square$ Je vérifie les tons en pinyin : fǔbài, tānwū, fǎlǜ, dǎozhì.
  \item $\square$ En version, je traduis d'abord le sens global, puis je rédige en français correct.
  \item $\square$ En thème, je repère les mots-outils nécessaires (的、了、得) avant de traduire.
  \item $\square$ Je relis mon texte : ponctuation chinoise pleine chasse ，。？ et caractères simplifiés corrects.
\end{itemize}
\end{aretenir}

\findecours

\end{document}
